% Prototype C at volume scale — originals verso, English recto.
% Auto-generated by tools/transcript2tex.py --layout loeb.
\documentclass[10pt,twoside]{book}
% Shared preamble for the M2 layout prototypes (xelatex + Cardo)
% ⚠ KDP flagged insufficient gutter at 625pp (needs >=0.75in clear from the trim
% edge). The true cause: \plnumber's \llap ruler (below) pokes INTO the inner
% margin by design, so the nominal 0.80in inner left only ~0.66in actually clear
% -- verified by measuring page 21's rendered ink, not by guessing. Fixed by
% moving 0.15in from outer to inner, which leaves \textwidth (inner+outer
% unchanged) untouched -- no reflow, no page-count shift, no cover rebuild.
\usepackage[paperwidth=6in,paperheight=9in,
  inner=0.95in,outer=0.50in,top=0.80in,bottom=0.85in]{geometry}
\usepackage{fontspec}
\setmainfont{Cardo}[
  BoldFont={* Bold}, ItalicFont={* Italic},
  Ligatures=TeX]
\usepackage{graphicx}
\usepackage{xcolor}
\definecolor{rulegray}{gray}{0.55}

% a sense-line: #1 = indent in em, #2 = text; long lines wrap with hanging indent
\newlength{\plhang}\setlength{\plhang}{5em}

% ---- sense-line numbering, which is what ANCHORS both foot-bands ----------
% The 1853's own apparatus keys every reading as PAGE. LINE ("34. 27"), and our
% transcripts are line-faithful, so OUR SENSE-LINE COUNT IS ITS LINE COUNT.
% Verified at printed 34: Τὰ ἔργα τῶν χειρῶν σου μὴ παρίδῃς is line 27, exactly
% as class A says. The apparatus was already speaking this coordinate; the page
% simply was not showing it.
%
% ⚠ ONE NUMBERING SERVES ALL THREE COLUMNS. Part I's Latin recto is line-for-line
% with its Greek verso, and our English is line-keyed to the Greek — verified at
% printed 34, where Greek 27 and English 27 are the same line and both pages run
% to 30. So the number set beside the Greek reads across the whole opening, and
% the Latin is deliberately NOT numbered: a second set of figures on the same
% leaf would say the same thing twice.
%
% ⚠ Numbering counts \pl calls, i.e. SENSE-lines, not typeset lines. That is the
% correct unit and it is also the only stable one — 839 of Part I's lines turn in
% the two columns, and a turned line must not take a number of its own.
\newcounter{plline}\newcounter{plmod}
\newif\ifplnum \plnumfalse
\newcommand{\plreset}{\setcounter{plline}{0}\setcounter{plmod}{0}}

% ⚠ EVERY FIVE LINES IS NOT AN ANCHOR — it is a RULER. A note keyed to line 27 is
% unfindable if the margin only prints 25 and 30: the reader must count, and
% nothing on line 27 says a note is there to count towards. Wilson, on the second
% sample: "there is still no visible anchor."
%
% So there are TWO devices, and they are deliberately different alphabets:
%   · the MARGIN carries ARABIC line figures — the ruler, and the 1853's own
%     PAGE. LINE coordinate, which the apparatus has always cited;
%   · the LINE-END carries a ROMAN superscript — the anchor, which says "there is
%     a note on me" at the place the eye already is.
% Roman against arabic so the two can never be read as the same series. The band
% below repeats the roman, so a marker is matched by looking, not by counting.
%
% A line that carries a note prints its arabic figure too, whatever the every-five
% rule says — the ruler should not go silent exactly where it is wanted.
% ⚠ THE LATIN COLUMN IS NUMBERED ON ITS OWN COUNT, and this was got wrong first
% time. The Latin was left unnumbered on the ground that it is line-for-line with
% the Greek so a second set of figures would repeat them. tools/check_lineparity.py
% shows that is true of 111 of Part I's 131 openings and FALSE of 20 — the plate
% turns its own long lines and the Latin is the wordier column, so it can run five
% or six lines longer than its Greek (printed 48/49, 62/63, 142/143). One set of
% figures would therefore have named the wrong Latin line on one opening in seven,
% and nothing on the page would have shown it.
%
% The MARKER SERIES still runs continuously across the leaf: the Greek's noted
% lines take i…k and the Latin's take k+1…m, so a roman is unique on the page and
% the band needs no column label. \setnotedcol carries the offset.
\newcommand{\notedlines}{}
\newcounter{plofs}
\newcommand{\setnoted}[1]{\renewcommand{\notedlines}{#1}\setcounter{plofs}{0}}
\newcommand{\setnotedcol}[2]{\renewcommand{\notedlines}{#1}\setcounter{plofs}{#2}}
\newif\ifplnoted
\newcounter{plpos}\newcounter{plfound}
\makeatletter
% Walks the noted-line list and records BOTH that this line is noted and its
% position in the list, which is the roman numeral the band will repeat.
\newcommand{\plcheck}[1]{\plnotedfalse
  \setcounter{plpos}{0}\setcounter{plfound}{0}%
  \def\pl@empty{}\ifx\notedlines\pl@empty\else
    \@for\pl@x:=\notedlines\do{\stepcounter{plpos}%
      \ifnum\pl@x=#1\relax\plnotedtrue\setcounter{plfound}{\value{plpos}}\fi}%
  \fi}
\makeatother

% ⚠ NOT \numexpr modulo: \numexpr divides with ROUNDING, so n/5*5=n is true for
% n=3 as well as n=5. Count down in a second counter instead.
\newcommand{\plnumber}{%
  \ifplnum
    \stepcounter{plline}\stepcounter{plmod}%
    \plcheck{\value{plline}}%
    % ⚠ EVERY FIVE ONLY. An earlier pass also printed a figure on every noted
    % line, which was right while the arabic WAS the anchor and became clutter
    % the moment the roman marker took that job — printed 34 came out with
    % figures on eighteen of thirty lines. The two devices divide cleanly: roman
    % anchors, arabic measures.
    \ifnum\value{plmod}=5
      \setcounter{plmod}{0}%
      \llap{{\tiny\color{rulegray}\theplline}\hspace{0.6em}}%
    \fi
  \fi}

% The roman marker, set at the END of the line — \plcheck has already run for this
% line inside \plnumber, so \ifplnoted and \plfound still hold its values.
\newcommand{\plmark}{%
  \ifplnum\ifplnoted
    \textsuperscript{\normalfont
      \romannumeral\numexpr\value{plfound}+\value{plofs}\relax}%
  \fi\fi}
% The number laps into the margin BEFORE the line's own indent, so it sits at the
% block's left edge whatever the indent level — otherwise a deeply indented line
% would carry its figure into the middle of the column.
\newcommand{\pl}[2]{\par\noindent\hangindent=\plhang\hangafter=1
  \plnumber\hspace*{#1em}#2\plmark}
% vertical breathing space between prayer paragraphs
\newcommand{\parasep}{\par\vspace{0.55\baselineskip}}
% small-caps language label used by the stacked layout
\newcommand{\langhead}[1]{\par\vspace{0.4\baselineskip}%
  {\centering\scshape\color{rulegray}#1\par}\vspace{0.35\baselineskip}}
% unit title
\newcommand{\unittitle}[1]{\par\vspace{1.2\baselineskip}%
  {\centering\large\scshape #1\par}\vspace{0.8\baselineskip}}
% thin separator rule
\newcommand{\thinrule}{\par\vspace{0.5\baselineskip}%
  {\centering\color{rulegray}\rule{0.35\textwidth}{0.4pt}\par}\vspace{0.5\baselineskip}}
\raggedbottom
\setlength{\parindent}{0pt}
% bidi must be loaded LAST (provides \RL for the Hebrew headings)
\usepackage{bidi}

\usepackage{fancyhdr}
\usepackage{ifthen}
\setlength{\headheight}{14pt}

\newcommand{\headL}{ΕΥΧΑΙ ΚΑΘΗΜΕΡΙΝΑΙ.\ $\cdot$\ PRECES QUOTIDIANÆ.}
\newcommand{\headR}{THE DAILY PRAYERS.}
\newcommand{\sethead}[1]{\renewcommand{\headL}{#1}\renewcommand{\headR}{#1}}
\pagestyle{fancy}\fancyhf{}
\fancyhead[CE]{\small\headL}
\fancyhead[CO]{\small\headR}
\fancyhead[LE,RO]{\small\thepage}
\fancyhead[RE,LO]{\footnotesize\color{rulegray}[\,\rightmark\,]}
\renewcommand{\headrulewidth}{0pt}

\newcommand{\parthead}[1]{\clearpage\thispagestyle{empty}\vspace*{2.2in}%
  \addcontentsline{toc}{part}{#1}%
  {\centering\Huge\scshape #1\par}\clearpage}

\newwrite\fitfile
\immediate\openout\fitfile=\jobname.fit
\newcommand{\versoalign}{\clearpage
  \ifodd\value{page}
    \null\thispagestyle{empty}\write\fitfile{B \thepage}\clearpage
  \fi}

% A minipage cannot break across leaves: it overruns in silence. So every unit is
% boxed, its height written out, and only then set — python compares the height
% against \textheight. This is the measurement the mirror got from page marks.
\newcommand{\originals}[3]{%
  \setbox0=\vbox{\noindent
    \begin{minipage}[t]{0.485\textwidth}
      \footnotesize\setlength{\plhang}{1.5em}\plreset\plnumtrue #2
    \end{minipage}\hfill
    \begin{minipage}[t]{0.485\textwidth}
      \footnotesize\setlength{\plhang}{1.5em}\plreset\plnumtrue
      \setnoted{}#3
    \end{minipage}\par}%
  % ⚠ \ht0 ALONE IS WRONG HERE and silently reported ~5.7pt for every Part I
  % verso until 2026-08-09. This vbox's content is ONE LINE of two tall boxes, so
  % its height is just that line's height above the baseline and the whole column
  % sits in the DEPTH. Measure ht+dp. (\originalsone and \enleaf are multi-line
  % vboxes and were never affected.)
  \immediate\write\fitfile{H O #1 \the\dimexpr\ht0+\dp0\relax}\box0}

\newcommand{\originalsone}[3]{%
  \setbox0=\vbox{\hsize=\textwidth\footnotesize\setlength{\plhang}{1.5em}\plreset\plnumtrue #2}%
  \immediate\write\fitfile{H O #1 \the\ht0}%
  \ifthenelse{\equal{#3}{epigraph}}{\dropto{0.16}\box0}{\box0}}

\newcommand{\enleaf}[3]{%
  \setbox0=\vbox{\hsize=\textwidth\small\plreset\plnumtrue #2}%
  \immediate\write\fitfile{H E #1 \the\ht0}%
  \ifthenelse{\equal{#3}{epigraph}}{\dropto{0.16}\box0}{\box0}}

% ---- the two foot-bands -------------------------------------------------
% The Loeb leaves roughly 44% of every leaf empty, and a Loeb's foot is where an
% apparatus belongs. The verso band carries the apparatus criticus against the
% originals; the recto band carries explanatory notes against the English.
%
% ⚠ The band is set at the FOOT, which means the leaf must be composed as a box
% of exactly \textheight with \vfill between body and band. \vfil alone will not
% do it: the preamble sets \raggedbottom, whose own bottom glue competes with it
% and lands the band a third of the way up. Same trap as the display leaves.
%
% ⚠ The band is measured INTO the fit record (\dimexpr body + band), not beside
% it. A band that overruns is exactly as silent as a minipage that overruns, and
% the point of the instrument is that nothing is unmeasured — which is how Part
% I's verso went 131 leaves without ever being tested.
% ⚠ The \par after the rule is load-bearing: without it the rule is just an inline
% box and the band's first word sets on the same line as the rule.
\newcommand{\bandrule}{\par\vspace{5pt}%
  {\noindent\color{rulegray}\rule{0.28\textwidth}{0.4pt}\par}\vspace{3.5pt}}

\newcommand{\apparatusband}[1]{%
  \bandrule{\scriptsize\setlength{\parindent}{0pt}\setlength{\parskip}{1pt}%
    \raggedright #1\par}}

\newcommand{\notesband}[1]{%
  \bandrule{\footnotesize\setlength{\parindent}{0pt}\setlength{\parskip}{2pt}%
    #1\par}}

% Compose a leaf whose body sits at the top and whose band sits at the foot.
\newcommand{\bandleaf}[2]{\vbox to \textheight{#1\vfill#2}}

% #5 = the Latin column's own noted lines, #6 = how many romans the Greek column
% already spent, so the series runs on rather than restarting at i.
\newcommand{\originalsn}[6]{%
  \setbox0=\vbox{\noindent
    \begin{minipage}[t]{0.485\textwidth}
      \footnotesize\setlength{\plhang}{1.5em}\plreset\plnumtrue #2
    \end{minipage}\hfill
    \begin{minipage}[t]{0.485\textwidth}
      % ⚠ numbered on ITS OWN count — see the preamble. It is NOT reliably
      % line-for-line with the Greek: 20 of Part I's 131 openings differ, by as
      % much as six lines, because the plate turns its own long lines.
      \footnotesize\setlength{\plhang}{1.5em}%
      \setnotedcol{#5}{#6}\plreset\plnumtrue #3
    \end{minipage}\par}%
  \setbox2=\vbox{\hsize=\textwidth\apparatusband{#4}}%
  \immediate\write\fitfile{H O #1 \the\dimexpr\ht0+\dp0+\ht2+\dp2\relax}%
  \bandleaf{\box0}{\box2}}

\newcommand{\originalsonen}[4]{%
  \setbox0=\vbox{\hsize=\textwidth\footnotesize\setlength{\plhang}{1.5em}\plreset\plnumtrue #2}%
  \setbox2=\vbox{\hsize=\textwidth\apparatusband{#4}}%
  \immediate\write\fitfile{H O #1 \the\dimexpr\ht0+\ht2+\dp2\relax}%
  \bandleaf{\box0}{\box2}}

\newcommand{\enleafn}[4]{%
  \setbox0=\vbox{\hsize=\textwidth\small\plreset\plnumtrue #2}%
  \setbox2=\vbox{\hsize=\textwidth\notesband{#4}}%
  \immediate\write\fitfile{H E #1 \the\dimexpr\ht0+\ht2+\dp2\relax}%
  \bandleaf{\box0}{\box2}}

% Place a display block a fixed way down the leaf. NOT \vfil: the preamble sets
% \raggedbottom, whose own bottom glue competes with any \vfil pair and lands the
% block about a third of the way down by accident. A title page sits where it is
% put, so put it.
\newcommand{\dropto}[1]{\vspace*{#1\textheight}}

% A leaf composed as one box of exactly the text height, so \vfill inside it
% divides the leaf instead of competing with \raggedbottom's bottom glue.
\newcommand{\fullleaf}[1]{\thispagestyle{empty}%
  \vbox to \textheight{#1}\clearpage}

% A display leaf: the 1853's own title and half-title pages, which are not running
% text and must not be set as though they were. \pl is redefined to drop its
% indent and centre, so the fragment renders as the page it is.
\newcommand{\titleleaf}[2]{\thispagestyle{empty}%
  \dropto{0.28}%
  \begingroup\centering
    \renewcommand{\parasep}{\par\vspace{2.4\baselineskip}}%
    \renewcommand{\pl}[2]{\par{\Large\scshape ##2}}%
    #2\par
  \endgroup}

\begin{document}
\frontmatter
\pagestyle{plain}
\immediate\write\fitfile{T \the\textheight}

% half-title
\fullleaf{\vspace*{0.30\textheight}
  {\centering{\LARGE\scshape Preces Privatae}\par}\vfill}
\fullleaf{}

% title page
\fullleaf{\vspace*{0.14\textheight}
  {\centering
  {\Huge\scshape Preces Privatae}\\[0.35em]
  {\large\itshape Private Prayers}\\[2.4em]
  {\Large Lancelot Andrewes}\\[0.6em]
  {\small\itshape Bishop of Winchester, 1555--1626}\\[3em]
  {\small Greek, Latin and Hebrew as printed in the}\\
  {\small edition of J.\,H. Parker, Oxford, 1853,}\\
  {\small reprinting the Sheldonian text of 1675,}\\[0.5em]
  {\small with a new English translation}\\[3em]{\small\scshape Part II \textemdash{} Preces Quotidianae}\par}
  \vfill
  {\centering\small\scshape Wroot Press\par}}

% colophon
\fullleaf{\vfill
  {\raggedright\footnotesize
  The Greek, Latin and Hebrew are transcribed from the 1853 Parker
  edition, which is in the public domain.\par\medskip
  The English translation, the apparatus and the editorial matter are
  \copyright{} Wroot Press, and are issued under a Creative Commons
  Attribution\,--\,NonCommercial 4.0 licence.\par\medskip
  Set in Cardo.\par}}

\tableofcontents
\clearpage

% The editor's preface. DRAFT — Wilson's to edit; the voice should be his.
% Every claim here is grounded in the repo: CONVENTIONS.md for the rules,
% apparatus/THE-LATIN-QUESTION.md for printed 34/35 and 174/175,
% apparatus/CLASS-A-ledger.md for the dropped Hebrew, CLASS-B-unit-boundaries.md
% for the blank leaves after printed 184, EDITION-SHAPE.md for the layout.
% If you change a claim here, change it there too.
\begingroup
% The body of the book is unindented sense-lines, so the preface separates its
% paragraphs by space rather than by indent. Scoped: it must not leak into the text.
\setlength{\parskip}{0.55\baselineskip}
\setlength{\parindent}{0pt}
\addcontentsline{toc}{section}{Preface}
\markboth{Preface}{Preface}
{\centering\Large\scshape Preface\par}
\vspace{1.5\baselineskip}

Lancelot Andrewes kept a book of prayers in his own hand, in Greek, in Latin, and
in Hebrew, and did not publish it. What survives is a series of witnesses to it:
manuscripts that differ from one another, and printed editions that differ from
the manuscripts. This edition sets out one of those witnesses whole --- the text
printed at Oxford by J.\,H. Parker in 1853, itself a reprint of the Sheldonian
edition of 1675 --- and puts a new English translation beside it.

The English translations of Andrewes now in print all descend from one
nineteenth-century construal. No edition currently in print, and none in the
public domain, presents the Greek and the Latin at all. That is the gap this book
is for. A reader who wants to know what Andrewes actually wrote, rather than what
a translator made of it, has until now had to find a scan.

\vspace{0.8\baselineskip}
{\scshape The text.}\quad
What is printed here on the left-hand pages is a record of the 1853 plate, not an
improvement on it. Accents, breathings, punctuation, capitals and ligatures stand
as they were set. Where the 1853 gives a scripture reference that is wrong, the
wrong reference is printed and the note says so; where a sort is broken or a word
has been lost to a defect in the plate, the loss is marked rather than mended.
Nothing has been silently corrected. An edition that quietly repairs its source
is no longer a witness to anything, and this one is meant to be usable as
evidence.

The consequence is that a few things in this book are known to be wrong and are
left wrong on purpose. They are marked, and the reasons are given.

\vspace{0.8\baselineskip}
{\scshape Which language is Andrewes'.}\quad
In the first part the Greek is Andrewes' own text and the Latin faces it. The
translation here follows the Greek, and treats the Latin as the first witness to
how the Greek was construed by someone much nearer to it than we are. Where the
two genuinely diverge, the Greek is followed and the Latin's reading is recorded.
In the second and third parts, which are largely Latin, the Latin is the text and
is translated as such. Where Hebrew stands in the text it is kept, and translated;
the English layer never silently drops a script.

The 1853 editor was modest about his Latin. His apparatus says of the Greek-only
manuscript he collated that it lacks the Latin, \emph{quam Editor ipse confecisse
videtur} --- which the editor seems to have made himself. That would make the
Latin of the first half of Part I an editor's Latin, and worth less as a witness
than it looks.

It is worth more than he claimed. At printed 34--35 the manuscript reads
πλάσμα where the printed Greek reads τὰ ἔργα, and the facing Latin reads
\emph{Figmenta} --- which departs from the printed Greek, and departs also from
the Vulgate of the psalm it is quoting, which reads \emph{opera manuum tuarum}.
There is no route to \emph{Figmenta} except from a Greek text reading πλάσμα. At
printed 174--175 the Latin \emph{Ablue} renders ἀπόλουσαι, wash thou away, where
the printed Greek has ἀπόλυσαι, release thou --- one letter apart, and the Latin
saw the other one.
The Latin column is therefore not simply downstream of the Greek beside it: in
places it sees a manuscript the printed Greek does not. That is a small finding
with a large consequence for how this book should be read, and it is why the
Latin is set as a column and not as a footnote.

\vspace{0.8\baselineskip}
{\scshape The English.}\quad
The translation is new, and it is line-keyed: one English line to one line of the
original, at the same indentation, with the same brace-groups and columns. A
reader should be able to run a finger down the page and stay in step. That
constraint costs something in fluency and is worth its cost, because the shape of
these prayers on the page --- the lists, the widening and narrowing, the single
words given a line to themselves --- is not decoration. It is how Andrewes prayed.

The register is the English of the Authorised Version and the Prayer Book:
\emph{thou} and \emph{thee} to God, and the cadences an English reader already
carries. The syntax, though, stays plain. Archaism is confined to the pronouns and
verb-endings of address; there are no Jacobean inversions the Greek does not
force. For the Psalms the crib is Coverdale's Psalter, the Prayer Book's, ahead of
the Authorised Version --- not from preference but because Coverdale worked
through the Latin and the Greek, as Andrewes did, and so agrees with him in
exactly the places where the Authorised Version, made from the Hebrew, does not.
Andrewes never quoted an English Bible. He was quoting the Septuagint and the
Vulgate, and Coverdale stands nearer to that than any later English does.

\vspace{0.8\baselineskip}
{\scshape The page.}\quad
The originals are on the left-hand page and the English on the right. In Part I,
where the 1853 sets Greek and Latin on facing pages, both share the left-hand page
in two columns, so that the whole of an opening of the 1853 stands under the eye
at once with its translation beside it.

An earlier plan reproduced the 1675 opening exactly --- Greek verso, Latin recto,
page for page, with the English as a register along the foot. It was the more
beautiful arrangement and it was given up, for a reason worth recording: it made
every page of the book depend on every page fitting. It also had no room for the
English in those parts of the book that are in one language only, where there is
no facing page to spend. The present arrangement never depends on a page fitting,
and gives the English a page of its own throughout.

The page numbers running along the outer edge are this edition's. The numbers in
brackets on the inner edge are the 1853's, and every page of it is here and
findable by them.

\vspace{0.8\baselineskip}
{\scshape What is not here.}\quad
Two things, and it is better to say so.

The 1853 does not print all of Andrewes' Hebrew. Its own apparatus records Hebrew
standing in the manuscript at more than a dozen places where the printed page has
none --- among them a phrase of rabbinic anthropology, the \emph{tohu wa-bohu} of
the second verse of Genesis, and a short Greek-and-Hebrew lexicon of the words for
sin. This edition prints the 1853, so that Hebrew is recorded in the apparatus
and is not restored to the text. Restoring it would be a different book: a
critical text of Andrewes rather than a faithful setting of a printed edition, and
it would need the manuscripts themselves.

And the manuscript has silences the print does not show. Three blank leaves stand
in it after what is here printed page 184, in the middle of the seventh day. What
that silence means --- whether a division has been lost, or something was meant to
be added --- cannot be settled from a printed edition. It is recorded and left
open.

\vspace{0.8\baselineskip}
{\scshape Acknowledgement.}\quad
F.\,E. Brightman's translation of 1903 remains the best English understanding of
this book, and has been consulted throughout --- after drafting, never before, and
never copied. His arrangement is his own: he judged the second part of the 1675
unusable as printed and rearranged the whole for practical use. That judgement is
defensible, and its consequence is that the \emph{Preces Privatae} as a book, in
its own order, has been out of reach of English readers for a long time. Restoring
that order is much of what this edition is for. Where his notes have corrected a
reference or settled an attribution here, the note says so.

\vspace{1.2\baselineskip}
{\raggedleft\itshape Wroot Press\par}
\endgroup
\clearpage


\mainmatter
\pagestyle{fancy}

\parthead{Part II --- Preces Quotidianae}
\sethead{§1 \emph{Duo in me cognosco}}
% ---------- printed 267 ----------
\versoalign
\markright{267}
\setnoted{}
\addcontentsline{toc}{section}{§1 \emph{Duo in me cognosco}}
\originalsone{267}{\input{fragments/la267}}{}
\clearpage
\markright{267}
\setnoted{6}
\enleafn{267}{\input{fragments/en267}}{}{{\scriptsize \textsuperscript{i}\,Ps. lxxviii. 39 — \emph{a wind that passeth away, and cometh not again}\par}}
\clearpage
% ---------- printed 268 ----------
\versoalign
\markright{268}
\setnoted{}
\originalsone{268}{\input{fragments/la268}}{}
\clearpage
\markright{268}
\setnoted{}
\enleaf{268}{\input{fragments/en268}}{}
\clearpage
% ---------- printed 269 ----------
\versoalign
\markright{269}
\setnoted{}
\originalsone{269}{\input{fragments/la269}}{}
\clearpage
\markright{269}
\setnoted{}
\enleaf{269}{\input{fragments/en269}}{}
\clearpage
\sethead{§2 \emph{Preces Matutinæ}}
% ---------- printed 270 ----------
\versoalign
\markright{270}
\setnoted{}
\addcontentsline{toc}{section}{§2 \emph{Preces Matutinæ}}
\originalsone{270}{\input{fragments/la270}}{}
\clearpage
\markright{270}
\setnoted{2,5,7,9,11,17,21,23}
\enleafn{270}{\input{fragments/en270}}{}{{\scriptsize \textsuperscript{i}\,Ps. lxv. 2 — \emph{O thou that hearest prayer, unto thee shall all flesh come} \textperiodcentered\ \textsuperscript{ii}\,Ps. lv. 17 — \emph{evening, and morning, and at noon, will I pray} \textperiodcentered\ \textsuperscript{iii}\,Ps. v. 3 — \emph{my voice shalt thou hear in the morning} \textperiodcentered\ \textsuperscript{iv}\,Ps. cxli. 2 — \emph{let my prayer be set forth before thee as incense} \textperiodcentered\ \textsuperscript{v}\,Ps. lxiii. 7 — \emph{because thou hast been my help} \textperiodcentered\ \textsuperscript{vi}\,Lam. iii. 41 — \emph{let us lift up our heart with our hands unto God}, which the prayer turns to the singular \textperiodcentered\ \textsuperscript{vii}\,Ps. cxxiii. 2 — \emph{until that he have mercy upon us} \textperiodcentered\ \textsuperscript{viii}\,Ps. cxix. 132 — \emph{as thou usest to do unto those that love thy name}\par}}
\clearpage
% ---------- printed 271 ----------
\versoalign
\markright{271}
\setnoted{}
\originalsone{271}{\input{fragments/la271}}{}
\clearpage
\markright{271}
\setnoted{2,4,6,8,15,16,17,20,24,26,28}
\enleafn{271}{\input{fragments/en271}}{}{{\scriptsize \textsuperscript{i}\,Ps. xci. 11 — \emph{he shall give his angels charge over thee} \textperiodcentered\ \textsuperscript{ii}\,Ps. xxv. 4 — \emph{teach me thy paths} \textperiodcentered\ \textsuperscript{iii}\,Ps. cxix. 133 — \emph{let not any iniquity have dominion over me} \textperiodcentered\ \textsuperscript{iv}\,Ps. xvii. 5 — \emph{that my footsteps slip not} \textperiodcentered\ \textsuperscript{v}\,Tit. iii. 4 — the Greek word Andrewes keeps, φιλάνθρωπε, \emph{the love of God toward man} \textperiodcentered\ \textsuperscript{vi}\,Jas. v. 11 — πολυεύσπλαγχνος, \emph{the Lord is very pitiful} \textperiodcentered\ \textsuperscript{vii}\,2 Cor. i. 3 — \emph{the Father of mercies} \textperiodcentered\ \textsuperscript{viii}\,Luke xv. 19 — the prodigal: \emph{no more worthy to be called thy son} \textperiodcentered\ \textsuperscript{ix}\,Mark ix. 24 — \emph{help thou mine unbelief} — here made \emph{help thou mine impenitence} \textperiodcentered\ \textsuperscript{x}\,Luke i. 78 — \emph{through the tender mercy of our God} \textperiodcentered\ \textsuperscript{xi}\,Eph. i. 7 — \emph{according to the riches of his grace}\par}\par\vspace{3pt}\textsuperscript{ix}\,\textbf{He alters the verse as he prays it, and the alteration is the petition.} The father in Mark cries \emph{help thou mine unbelief}; Andrewes writes \emph{adjuva Tu impœnitentiam meam} — help thou mine \textbf{impenitence}. The substitution is deliberate and is not a misquotation: what he wants help against is not doubt but hardness. The English keeps his word and the note keeps the Gospel's, because the reader needs to hear both to hear either.}
\clearpage
% ---------- printed 272 ----------
\versoalign
\markright{272}
\setnoted{}
\originalsone{272}{\input{fragments/la272}}{}
\clearpage
\markright{272}
\setnoted{3,4,5,6,8,9,10,14,17,18}
\enleafn{272}{\input{fragments/en272}}{}{{\scriptsize \textsuperscript{i}\,Eph. ii. 4 — \emph{for his great love wherewith he loved us} \textperiodcentered\ \textsuperscript{ii}\,Luke xviii. 13 — \emph{God be merciful to me a sinner} \textperiodcentered\ \textsuperscript{iii}\,1 Tim. i. 15 — \emph{of whom I am chief} \textperiodcentered\ \textsuperscript{iv}\,Ps. xlii. 7 — \emph{deep calleth unto deep} \textperiodcentered\ \textsuperscript{v}\,Rom. v. 20 — \emph{where sin abounded, grace did much more abound} \textperiodcentered\ \textsuperscript{vi}\,Rom. xii. 21 — \emph{overcome evil with good} \textperiodcentered\ \textsuperscript{vii}\,Jas. ii. 13 — \emph{mercy rejoiceth against judgment} \textperiodcentered\ \textsuperscript{viii}\,Matt. xvi. 16 — \emph{thou art the Christ, the Son of the living God} \textperiodcentered\ \textsuperscript{ix}\,John i. 29 — \emph{behold the Lamb of God, which taketh away the sin of the world} \textperiodcentered\ \textsuperscript{x}\,Luke xix. 10 — \emph{to seek and to save that which was lost}\par}\par\vspace{3pt}\textsuperscript{v}\,\textbf{\textup{Rom. v. 20} again, and cited right.} This is the third page to give the verse correctly — with printed 103 and 111's own facing evidence — against the \textup{Rom. v. 2} at printed 111. ⚠ And note what he does with it: the Vulgate's indicative \emph{superabundavit gratia}, grace \textbf{did} abound, becomes the subjunctive \emph{superabundet}, let it abound. He turns a statement of what happened into a request that it happen again.}
\clearpage
\sethead{§3 \emph{Intercessio}}
% ---------- printed 273 ----------
\versoalign
\markright{273}
\setnoted{}
\addcontentsline{toc}{section}{§3 \emph{Intercessio}}
\originalsone{273}{\input{fragments/la273}}{}
\clearpage
\markright{273}
\setnoted{6,18,20}
\enleafn{273}{\input{fragments/en273}}{}{{\scriptsize \textsuperscript{i}\,Tit. i. 5 — \emph{set in order the things that are wanting} \textperiodcentered\ \textsuperscript{ii}\,2 Tim. ii. 15 — ὀρθοτομεῖν, \emph{rightly dividing}; Gal. ii. 14 — \emph{walked not uprightly} \textperiodcentered\ \textsuperscript{iii}\,Rom. xii. 3 — \emph{not to think of himself more highly than he ought}\par}}
\clearpage
% ---------- printed 274 ----------
\versoalign
\markright{274}
\setnoted{}
\originalsone{274}{\input{fragments/la274}}{}
\clearpage
\markright{274}
\setnoted{6,23,25,26}
\enleafn{274}{\input{fragments/en274}}{}{{\scriptsize \textsuperscript{i}\,Ps. cxvii. 25 — the plate's figure, which is the \textbf{Vulgate's} numbering of the psalm our Bibles call cxviii: \emph{Save now, I beseech thee, O LORD: O LORD, I beseech thee, send now prosperity.} ⚠ \textbf{One verse supplies both lines}, and the reference stands under the second because it covers the pair. \textperiodcentered\ \textsuperscript{ii}\,Ps. lxxix. 13 — \emph{render unto our neighbours seven-fold into their bosom}; Prayer Book numbering \textperiodcentered\ \textsuperscript{iii}\,Ps. xxv. 12 — \emph{his seed shall inherit the land}; Prayer Book numbering \textperiodcentered\ \textsuperscript{iv}\,Ps. xli. 1 — \emph{blessed is he that considereth the poor}\par}\par\vspace{3pt}\textsuperscript{i}\,⚠⚠ \textbf{Two lines of Hebrew, and they are the two halves of one verse.} \textup{\RL{הושיעה נא}}, \emph{hoshi'ah-na}, \textbf{save now} — the cry the crowds took up on Palm Sunday, which English kept as \textbf{Hosanna}; and \textup{\RL{הצליחה נא}}, \emph{hatzlichah-na}, \textbf{prosper now}, which English never kept at all. The plate gives each half its own Latin line — \emph{salvum fac} and \emph{prosperare} — and prints the reference once, under the second. ⚠ \textbf{The half we made an acclamation of is a petition}, and its twin, which asks for success, dropped out of the liturgy entirely. One of the few places the 1853 let Andrewes' third script through. ✅ Both words re-read at native resolution 2026-08-10 and confirmed letter for letter.}
\clearpage
% ---------- printed 275 ----------
\versoalign
\markright{275}
\setnoted{}
\originalsone{275}{\input{fragments/la275}}{}
\clearpage
\markright{275}
\setnoted{21,23,25,29}
\enleafn{275}{\input{fragments/en275}}{}{{\scriptsize \textsuperscript{ii}\,Ps. cxlv. 10 — \emph{all thy works praise thee, and thy saints give thanks unto thee} \textperiodcentered\ \textsuperscript{iii}\,Ps. xcii. 1 — the psalm \emph{for the Sabbath day} \textperiodcentered\ \textsuperscript{iv}\,Ps. cxlv. 1 — \emph{I will magnify thee, O God my King}\par}\par\vspace{3pt}\textsuperscript{i}\,\textbf{The Thanksgiving opens on the Sabbath psalm, and it opens where the intercession stopped.} The petitions above end \emph{we beseech Thee, hear us, O Lord} — eight numbered asks, the last of them for the pleader and his kindred together — and then, on the same leaf, the section changes and the same voice begins to give thanks. \textbf{The 1853 sets both on one page and this edition keeps them there}; the join is the point, and a reader who meets \emph{Thanksgiving} at the head of a fresh leaf loses it. Ps. xcii is titled \emph{for the Sabbath day}, so the praise begins on the day appointed for it — morning loving-kindness and night-season truth in the next breath.}
\clearpage
\sethead{§4 \emph{Gratiarum Actio}}
% ---------- printed 276 ----------
\versoalign
\markright{276}
\setnoted{}
\addcontentsline{toc}{section}{§4 \emph{Gratiarum Actio}}
\originalsone{276}{\input{fragments/la276}}{}
\clearpage
\markright{276}
\setnoted{5,7,8,11,13,16,17,19,20,22,24,26}
\enleafn{276}{\input{fragments/en276}}{}{{\scriptsize \textsuperscript{i}\,Rom. iv. 17 — \emph{calleth those things which be not as though they were} \textperiodcentered\ \textsuperscript{ii}\,Col. i. 16 — \emph{visible and invisible} \textperiodcentered\ \textsuperscript{iii}\,Heb. i. 3 — \emph{upholding all things by the word of his power} \textperiodcentered\ \textsuperscript{iv}\,Acts xiv. 17 — \emph{filling our hearts with food and gladness} \textperiodcentered\ \textsuperscript{v}\,Ps. cxix. 91 — \emph{for all are thy servants} \textperiodcentered\ \textsuperscript{vi}\,Gen. ii. 7 — \emph{the breath of life} \textperiodcentered\ \textsuperscript{vii}\,Gen. i. 26 — \emph{in our image, after our likeness} \textperiodcentered\ \textsuperscript{viii}\,Ps. viii. 6 — \emph{thou madest him to have dominion} \textperiodcentered\ \textsuperscript{ix}\,Gen. ii. 15 — \emph{and put him into the garden of Eden} \textperiodcentered\ \textsuperscript{x}\,Acts xi. 18 — \emph{granted repentance unto life} \textperiodcentered\ \textsuperscript{xi}\,Gen. iii. 15 — the promised seed; 2 Pet. i. 4 — \emph{partakers of the divine nature} \textperiodcentered\ \textsuperscript{xii}\,Rom. i. 19 — \emph{that which may be known of God}\par}}
\clearpage
% ---------- printed 277 ----------
\versoalign
\markright{277}
\setnoted{}
\originalsone{277}{\input{fragments/la277}}{}
\clearpage
\markright{277}
\setnoted{3,5,7,10,13,15,16,17,21,23,28}
\enleafn{277}{\input{fragments/en277}}{}{{\scriptsize \textsuperscript{i}\,Gal. iv. 4 — \emph{when the fulness of the time was come} \textperiodcentered\ \textsuperscript{ii}\,Heb. ii. 16 — \emph{he took on him the seed of Abraham} \textperiodcentered\ \textsuperscript{iii}\,Phil. ii. 7 — \emph{took upon him the form of a servant} \textperiodcentered\ \textsuperscript{iv}\,Eph. v. 2 — \emph{an offering and a sacrifice to God} \textperiodcentered\ \textsuperscript{v}\,Gal. iii. 13 — \emph{made a curse for us} \textperiodcentered\ \textsuperscript{vi}\,Isa. v. 4 — the vineyard: \emph{what could have been done more?} \textperiodcentered\ \textsuperscript{vii}\,2 Pet. i. 4 — \emph{partakers of the divine nature} \textperiodcentered\ \textsuperscript{viii}\,2 Cor. ii. 14 — \emph{the savour of his knowledge} \textperiodcentered\ \textsuperscript{ix}\,1 Pet. iii. 2 — \emph{your chaste conversation coupled with fear} \textperiodcentered\ \textsuperscript{x}\,Heb. xi. 37 — of those who suffered \emph{even unto blood} \textperiodcentered\ \textsuperscript{xi}\,Heb. ix. 15 — \emph{they which are called might receive the promise}\par}\par\vspace{3pt}\textsuperscript{vi}\,⚠ \textbf{Isaiah's question is asked of Christ, not of Israel.} \emph{What could have been done more to my vineyard, that I have not done in it?} is God's complaint against a people that yielded wild grapes; Andrewes turns it into praise — \emph{nihil non faciens quo facto opus erat}, leaving nothing undone that needed doing. \textbf{The reproach becomes a commendation} by changing who is speaking of whom, and that inversion is the sort of thing his method does constantly.}
\clearpage
% ---------- printed 278 ----------
\versoalign
\markright{278}
\setnoted{}
\originalsone{278}{\input{fragments/la278}}{}
\clearpage
\markright{278}
\setnoted{1,3,4,6,7,9,17,19,21,26,28,30}
\enleafn{278}{\input{fragments/en278}}{}{{\scriptsize \textsuperscript{i}\,1 Tim. iii. 15 — \emph{the pillar and ground of the truth} \textperiodcentered\ \textsuperscript{ii}\,Matt. xvi. 18 — \emph{the gates of hell shall not prevail against it} \textperiodcentered\ \textsuperscript{iii}\,1 Tim. vi. 20 — \emph{keep that which is committed to thy trust} \textperiodcentered\ \textsuperscript{iv}\,Col. ii. 5 — \emph{your order, and the stedfastness of your faith} \textperiodcentered\ \textsuperscript{v}\,2 Sam. vii. 13 — \emph{I will stablish the throne of his kingdom for ever} \textperiodcentered\ \textsuperscript{vi}\,Ps. cxlvii. 14 — \emph{filleth thee with the finest of the wheat} \textperiodcentered\ \textsuperscript{vii}\,Ps. cv. 22 — \emph{and teach his senators wisdom} \textperiodcentered\ \textsuperscript{viii}\,Jer. iii. 15 — \emph{pastors which shall feed you with knowledge and understanding} \textperiodcentered\ \textsuperscript{ix}\,Isa. ii. 4 — \emph{their spears into pruninghooks} \textperiodcentered\ \textsuperscript{x}\,Tit. iii. 5 — \emph{the renewing of the Holy Ghost} \textperiodcentered\ \textsuperscript{xi}\,Wisd. xi. 23 — \emph{thou winkest at the sins of men, because they should amend} \textperiodcentered\ \textsuperscript{xii}\,Isa. xxx. 18 — \emph{therefore will the LORD wait, that he may be gracious unto you}\par}}
\clearpage
% ---------- printed 279 ----------
\versoalign
\markright{279}
\setnoted{}
\originalsone{279}{\input{fragments/la279}}{}
\clearpage
\markright{279}
\setnoted{1,2,3,4,5,8,11,12}
\enleafn{279}{\input{fragments/en279}}{}{{\scriptsize \textsuperscript{i}\,Rom. ii. 5 — \emph{after thy hardness and impenitent heart} \textperiodcentered\ \textsuperscript{ii}\,Acts ii. 37 — at Pentecost: \emph{they were pricked in their heart} \textperiodcentered\ \textsuperscript{iii}\,Deut. xxxii. 29 — \emph{that they would consider their latter end} \textperiodcentered\ \textsuperscript{iv}\,Heb. x. 2 — \emph{the worshippers once purged should have had no more conscience of sins} \textperiodcentered\ \textsuperscript{v}\,Hos. ii. 15 — \emph{the valley of Achor for a door of hope} \textperiodcentered\ \textsuperscript{vi}\,1 John i. 9 — \emph{he is faithful and just to forgive us our sins}; Matt. xviii. 32 — \emph{I forgave thee all that debt}; John xx. 23 — \emph{whose soever sins ye remit}; Matt. xvi. 19 — \emph{the keys of the kingdom} \textperiodcentered\ \textsuperscript{vii}\,Isa. xxxviii. 12 — Hezekiah: \emph{from day even to night wilt thou make an end of me} \textperiodcentered\ \textsuperscript{viii}\,Ps. cii. 24 — \emph{take me not away in the midst of my days}\par}}
\clearpage
\sethead{§5 \emph{Deprecatio}}
% ---------- printed 280 ----------
\versoalign
\markright{280}
\setnoted{}
\addcontentsline{toc}{section}{§5 \emph{Deprecatio} (\emph{Ne perdas})}
\originalsone{280}{\input{fragments/la280}}{}
\clearpage
\markright{280}
\setnoted{4,5,10,22,23,26,27}
\enleafn{280}{\input{fragments/en280}}{}{{\scriptsize \textsuperscript{i}\,Isa. lxiii. 15 — \emph{look down from heaven, and behold from the habitation of thy holiness} \textperiodcentered\ \textsuperscript{ii}\,Ps. cxiii. 5 — \emph{who dwelleth on high, and humbleth himself to behold} \textperiodcentered\ \textsuperscript{iii}\,Ps. xliv. 1 — \emph{our fathers have told us, what work thou didst in their days} \textperiodcentered\ \textsuperscript{iv}\,Rev. vi. 16 — \emph{hide us from the face of him that sitteth on the throne} \textperiodcentered\ \textsuperscript{v}\,Matt. xxv. 33 — the goats set on the left hand \textperiodcentered\ \textsuperscript{vi}\,Matt. viii. 12 — \emph{cast out into outer darkness} \textperiodcentered\ \textsuperscript{vii}\,Jude 6 — the angels that kept not their first estate, \emph{reserved in everlasting chains}\par}}
\clearpage
% ---------- printed 281 ----------
\versoalign
\markright{281}
\setnoted{}
\originalsone{281}{\input{fragments/la281}}{}
\clearpage
\markright{281}
\setnoted{2,8,10,11,13,15,16}
\enleafn{281}{\input{fragments/en281}}{}{{\scriptsize \textsuperscript{i}\,Rev. xiv. 11 — \emph{the smoke of their torment ascendeth up for ever} \textperiodcentered\ \textsuperscript{ii}\,Eph. iv. 18 — \emph{the blindness of their heart} \textperiodcentered\ \textsuperscript{iii}\,Prov. vii. 13 — the strange woman, \emph{with an impudent face}; Isa. xlviii. 4 — \emph{thy brow brass} \textperiodcentered\ \textsuperscript{iv}\,1 Tim. iv. 2 — \emph{conscience seared with a hot iron} \textperiodcentered\ \textsuperscript{v}\,Tit. i. 16 — \emph{unto every good work reprobate} \textperiodcentered\ \textsuperscript{vi}\,1 John v. 16 — \emph{a sin unto death} \textperiodcentered\ \textsuperscript{vii}\,Matt. xii. 32 — the word against the Holy Ghost\par}\par\vspace{3pt}\textsuperscript{iii}\,\textbf{A soft forehead and a hard one, and both are prayed against.} Proverbs gives the \emph{impudent face} of the strange woman; Isaiah gives the brow of brass. Andrewes braces them as the two ways a face can fail — shamelessness that feels nothing and obstinacy that yields nothing — and asks to be delivered from either. The English keeps the brace rather than choosing.}
\clearpage
% ---------- printed 282 ----------
\versoalign
\markright{282}
\setnoted{}
\originalsone{282}{\input{fragments/la282}}{}
\clearpage
\markright{282}
\setnoted{15,16,17,18,19,20,22}
\enleafn{282}{\input{fragments/en282}}{}{{\scriptsize \textsuperscript{i}\,1 Tim. vi. 5 — \emph{men of corrupt minds, and destitute of the truth} \textperiodcentered\ \textsuperscript{ii}\,Hos. xi. 5 — \emph{the Assyrian shall be his king, because they refused to return} \textperiodcentered\ \textsuperscript{iii}\,1 Kings xii. 27 — Jeroboam, who feared the people would return to Jerusalem \textperiodcentered\ \textsuperscript{iv}\,1 Kings xii. 13 — Rehoboam, who answered the people roughly \textperiodcentered\ \textsuperscript{v}\,Josh. vii. 24 — the valley of Achor, where Achan and all that he had were brought up \textperiodcentered\ \textsuperscript{vi}\,Judg. ix. 23 — \emph{God sent an evil spirit between Abimelech and the men of Shechem} \textperiodcentered\ \textsuperscript{vii}\,Ps. cxx. 2 — \emph{deliver my soul from lying lips}; Ps. xci. 3 — \emph{the snare of the fowler}\par}\par\vspace{3pt}\textsuperscript{ii}\,\textbf{Four names for four ways a people is lost, and none of them is explained.} Asshur is the foreign power a nation is given over to; Jeroboam the ruler who divides worship to keep a throne; Rehoboam the ruler who answers roughly and loses ten tribes; the evil spirit of Shechem the discord God sends between a man and his own city. As at printed 204, the deprecation depends entirely on the reader knowing the histories, and the English can only name them.}
\clearpage
% ---------- printed 283 ----------
\versoalign
\markright{283}
\setnoted{}
\originalsone{283}{\input{fragments/la283}}{}
\clearpage
\markright{283}
\setnoted{4,9,22,23,25}
\enleafn{283}{\input{fragments/en283}}{}{{\scriptsize \textsuperscript{i}\,2 Cor. vii. 1 — \emph{all filthiness of the flesh and spirit} \textperiodcentered\ \textsuperscript{ii}\,2 Sam. xxiv. 16 — the angel at the threshingfloor: \emph{it is enough; stay now thine hand} \textperiodcentered\ \textsuperscript{iii}\,Matt. xxv. 21 — \emph{enter thou into the joy of thy lord} \textperiodcentered\ \textsuperscript{iv}\,John xvi. 24 — \emph{that your joy may be full} \textperiodcentered\ \textsuperscript{v}\,Matt. xxv. 34 — \emph{then shall the King say unto them on his right hand}\par}}
\clearpage
% ---------- printed 284 ----------
\versoalign
\markright{284}
\setnoted{}
\originalsone{284}{\input{fragments/la284}}{}
\clearpage
\markright{284}
\setnoted{1,13,16,18,19}
\enleafn{284}{\input{fragments/en284}}{}{{\scriptsize \textsuperscript{i}\,Luke xxiii. 43 — \emph{to day shalt thou be with me in paradise}; Luke xvi. 22 — \emph{carried by the angels into Abraham's bosom} \textperiodcentered\ \textsuperscript{ii}\,Eph. vi. 14 — the whole armour, which prayer crowns at verse 18 \textperiodcentered\ \textsuperscript{iii}\,John xv. 16 — \emph{that your fruit should remain} \textperiodcentered\ \textsuperscript{iv}\,Ps. xvii. 15 — \emph{I shall be satisfied, when I awake, with thy likeness} \textperiodcentered\ \textsuperscript{v}\,2 Kings xx. 6 — \emph{I will add unto thy days fifteen years}\par}}
\clearpage
% ---------- printed 285 ----------
\versoalign
\markright{285}
\setnoted{}
\originalsone{285}{\input{fragments/la285}}{}
\clearpage
\markright{285}
\setnoted{6,8,11}
\enleafn{285}{\input{fragments/en285}}{}{{\scriptsize \textsuperscript{i}\,Ps. cxxxii. 15 — \emph{I will abundantly bless her provision} \textperiodcentered\ \textsuperscript{ii}\,Ps. cxlvii. 14 — \emph{filleth thee with the finest of the wheat} \textperiodcentered\ \textsuperscript{iii}\,Ps. cxlvii. 13 — \emph{he hath blessed thy children within thee}\par}}
\clearpage
\sethead{§6 \emph{Sacrificium Vespertinum}}
% ---------- printed 286 ----------
\versoalign
\markright{286}
\setnoted{}
\addcontentsline{toc}{section}{§6 \emph{Sacrificium Vespertinum} + \emph{Horologium}}
\originalsone{286}{\input{fragments/la286}}{}
\clearpage
\markright{286}
\setnoted{3,6,8,9,11,12,15,16,17,20,21}
\enleafn{286}{\input{fragments/en286}}{}{{\scriptsize \textsuperscript{i}\,Ps. cxix. 55 — \emph{I have thought upon thy Name, O LORD, in the night} \textperiodcentered\ \textsuperscript{ii}\,Job xxxv. 10 — Elihu: none asketh \emph{Where is God my maker, who giveth songs in the night?} \textperiodcentered\ \textsuperscript{iii}\,Ps. lxv. 8 — the harvest psalm: the outgoings of the morning and evening made to rejoice \textperiodcentered\ \textsuperscript{iv}\,Ps. cxxvii. 3 — \emph{he giveth his beloved sleep}, verse 3 on the Prayer Book's numbering \textperiodcentered\ \textsuperscript{v}\,Acts i. 7 — the Ascension: \emph{it is not for you to know the times or the seasons} \textperiodcentered\ \textsuperscript{vi}\,Mark iv. 38 — the disciples wake Christ in the storm \textperiodcentered\ \textsuperscript{vii}\,Luke ii. 8 — the shepherds in the field by night, which fixes the hour of the Nativity \textperiodcentered\ \textsuperscript{viii}\,Titus iii. 5 — \emph{the washing of regeneration, and renewing of the Holy Ghost} \textperiodcentered\ \textsuperscript{ix}\,Gal. iv. 19 — Paul in travail again \emph{until Christ be formed in you}; Eph. iv. 13 — \emph{unto a perfect man} \textperiodcentered\ \textsuperscript{x}\,Mark xvi. 2 — the women at the sepulchre, \emph{at the rising of the sun} \textperiodcentered\ \textsuperscript{xi}\,Rom. vi. 4 — the baptismal argument: \emph{even so we also should walk in newness of life}\par}\par\vspace{3pt}\textsuperscript{iv}\,\textbf{The plate's \emph{cxxvii. 3} is not a slip.} On the Prayer Book's numbering \emph{he giveth his beloved sleep} is verse 3; the Authorised Version makes it verse 2. With the two Prayer Book references at printed 433 this is a further sign of which Psalter lay open. And Andrewes does not ask simply for sleep: \emph{somnum sanitatis}, the sleep of soundness — the health that sleep is the evidence of, rather than the sleep itself. Neither English Bible has the word, and our English keeps it.\\[1pt]\textsuperscript{v}\,\textbf{The hours are Roman, counted from sunrise.} The third is about nine in the morning, the sixth noon, the ninth about three in the afternoon, and the \emph{Horologium} walks them in order, giving each the Gospel event that falls there. A reader counting from midnight will put Pentecost before dawn and the darkness over the land in the small hours. Andrewes explains none of it, because his reader did not need it explained.}
\clearpage
% ---------- printed 287 ----------
\versoalign
\markright{287}
\setnoted{}
\originalsone{287}{\input{fragments/la287}}{}
\clearpage
\markright{287}
\setnoted{2,3,6,7,9,12,13,14,17,22,23,26}
\enleafn{287}{\input{fragments/en287}}{}{{\scriptsize \textsuperscript{i}\,Acts ii. 15 — Peter at Pentecost: \emph{it is but the third hour of the day} \textperiodcentered\ \textsuperscript{ii}\,Ps. li. 11 — David's Miserere: \emph{take not thy Holy Spirit from me} \textperiodcentered\ \textsuperscript{iv}\,Matt. xxvii. 45 — the darkness from the sixth hour unto the ninth \textperiodcentered\ \textsuperscript{v}\,Col. ii. 14 — the handwriting that was against us, blotted out and nailed to the cross \textperiodcentered\ \textsuperscript{vi}\,Acts x. 9, 11 — Peter on the housetop at the sixth hour; the sheet let down from heaven \textperiodcentered\ \textsuperscript{vii}\,Gal. ii. 15 — Paul to Peter at Antioch: \emph{sinners of the Gentiles} \textperiodcentered\ \textsuperscript{viii}\,Acts x. 16 — the sheet received up again into heaven \textperiodcentered\ \textsuperscript{ix}\,John iv. 52 — the nobleman's son: \emph{at the seventh hour the fever left him} \textperiodcentered\ \textsuperscript{x}\,Mark xv. 34 — at the ninth hour, \emph{Eloi, Eloi, lama sabachthani}; Heb. ii. 9 — \emph{that he should taste death for every man} \textperiodcentered\ \textsuperscript{xi}\,Col. iii. 5 — \emph{mortify therefore your members which are upon the earth} \textperiodcentered\ \textsuperscript{xii}\,Acts iii. 1 — Peter and John go up to the temple \emph{at the hour of prayer, being the ninth hour}\par}\par\vspace{3pt}\textsuperscript{iii}\,\textbf{The sixth hour and the sixth day are one figure.} \emph{Hora sexta dieque sexto} — man was made on the sixth day and remade at the sixth hour; Andrewes sets the two sixes side by side and leaves the join to the reader, as he leaves every join in this office. The hour is doing double duty on this page besides: the sixth is the darkness over the land and, five lines later, Peter's housetop vision. That is the \emph{Horologium}'s method throughout — an hour is not one event but everything the Gospels put there.}
\clearpage
% ---------- printed 288 ----------
\versoalign
\markright{288}
\setnoted{}
\originalsone{288}{\input{fragments/la288}}{}
\clearpage
\markright{288}
\setnoted{2,8,14,16,17,20,26,28}
\enleafn{288}{\input{fragments/en288}}{}{{\scriptsize \textsuperscript{i}\,John i. 41 — Andrew, having found Him, to his brother Simon: \emph{we have found the Messias} \textperiodcentered\ \textsuperscript{ii}\,Matt. xx. 6 — the householder at the eleventh hour: \emph{why stand ye here all the day idle?} \textperiodcentered\ \textsuperscript{iii}\,John xiii. 2 — \emph{and supper being ended}, at the washing of the feet \textperiodcentered\ \textsuperscript{v}\,1 Cor. xi. 29 — \emph{eateth and drinketh damnation to himself, not discerning the Lord's body} \textperiodcentered\ \textsuperscript{vi}\,Mark xv. 42 — \emph{when the even was come}, Joseph of Arimathea asks for the body \textperiodcentered\ \textsuperscript{vii}\,John xx. 22, 23 — the risen Christ breathes on them: \emph{whose soever sins ye remit}\par}\par\vspace{3pt}\textsuperscript{iv}\,\textbf{The bequests of the New Testament — \emph{legata Novi Testamenti}, and the Latin says what English cannot.} \emph{Testamentum} is at once the covenant and the last will, so the Supper is the reading of a will and what it conveys are \emph{legata}, legacies. English divided the word long ago: a testament is a book, a bequest is a legacy, and \emph{the legacies of the new covenant} would sound like a figure of speech where the Latin is simply the law of wills. Our English keeps \emph{bequests} rather than spending the double sense to make the covenant plain.\\[1pt]\textsuperscript{viii}\,\textbf{He asks for half of a power that was given whole.} John xx. 23 confers both the remitting of sins and the retaining, and the prayer takes the one and expressly declines the other — \emph{sed quoad remissionem, ne ad retentionem, Domine}. Nothing on the page argues the asymmetry; it is enacted rather than defended, as the eleventh-hour stanza above enacts its own and asks the wage of a full day after standing idle through all of it.}
\clearpage
% ---------- printed 289 ----------
\versoalign
\markright{289}
\setnoted{}
\originalsone{289}{\input{fragments/la289}}{}
\clearpage
\markright{289}
\setnoted{1,2,3,4,9,14,20,24,25,26,27,28}
\enleafn{289}{\input{fragments/en289}}{}{{\scriptsize \textsuperscript{i}\,Ps. cxix. 62 — \emph{at midnight I will rise to give thanks unto thee} \textperiodcentered\ \textsuperscript{ii}\,Acts xvi. 25 — Paul and Silas at midnight in the prison at Philippi \textperiodcentered\ \textsuperscript{iii}\,Job xxxv. 10 — Elihu: \emph{who giveth songs in the night} \textperiodcentered\ \textsuperscript{iv}\,Ps. lxiii. 6 — \emph{when I remember thee upon my bed} \textperiodcentered\ \textsuperscript{v}\,Matt. xxv. 6 — the ten virgins: \emph{at midnight there was a cry made} \textperiodcentered\ \textsuperscript{vi}\,Matt. xxvi. 74 — Peter's third denial, and immediately the cock crew \textperiodcentered\ \textsuperscript{vii}\,Matt. xxiv. 50 — \emph{in a day when he looketh not for him, and in an hour that he is not aware of}; Luke xii. 46 — the same, of the unfaithful steward \textperiodcentered\ \textsuperscript{viii}\,Ps. xliii. 3 — \emph{O send out thy light}; Ps. lxxiv. 16 — \emph{thou hast prepared the light and the sun} \textperiodcentered\ \textsuperscript{ix}\,Matt. v. 45 — \emph{he maketh his sun to rise on the evil and on the good} \textperiodcentered\ \textsuperscript{x}\,Eph. iv. 18 — the Gentiles, \emph{the blindness of their heart} \textperiodcentered\ \textsuperscript{xi}\,Ps. iv. 6 — \emph{Lord, lift thou up the light of thy countenance upon us} \textperiodcentered\ \textsuperscript{xii}\,Ps. xxxvi. 9 — \emph{in thy light shall we see light}\par}\par\vspace{3pt}\textsuperscript{viii}\,\textbf{The office ends on a ladder of light, and it is built from two psalms the editor had to finish.} The plate cites \emph{Ps.} xliii. 3 for \emph{Qui emittis lucem, creas auroram}, and only the first half is there — \emph{O send out thy light} is Ps. xliii, but \emph{createst the dawn} is Ps. lxxiv. 16, which the 1853 editor supplied in his own square brackets. From there the rungs mount without a break: light sent out, the sun on good and evil, the mind's blindness enlightened, the countenance lifted, light seen in light — and the last rung stands on the next page, \emph{and at the last, in the light of grace, the light of glory}. The page-turn is the 1853's, not the ladder's.}
\clearpage
% ---------- printed 290 ----------
\versoalign
\markright{290}
\setnoted{}
\originalsone{290}{\input{fragments/la290}}{}
\clearpage
\markright{290}
\setnoted{2,4,5,6,10,12,14,16,18,24,25,26,27}
\enleafn{290}{\input{fragments/en290}}{}{{\scriptsize \textsuperscript{i}\,Ps. cxxxvi. 25 — the psalm of \emph{his mercy endureth for ever}, at the last mercy it names: food to all flesh \textperiodcentered\ \textsuperscript{ii}\,Ps. lxxi. 5 — an old man's prayer: \emph{thou art my trust from my youth} \textperiodcentered\ \textsuperscript{iii}\,Acts xiv. 17 — Paul at Lystra: God \emph{filling our hearts with food and gladness} \textperiodcentered\ \textsuperscript{iv}\,Heb. xiii. 9 — \emph{that the heart be stablished with grace, not with meats} \textperiodcentered\ \textsuperscript{v}\,Ps. cii. 11 — \emph{my days are gone like a shadow} \textperiodcentered\ \textsuperscript{vi}\,Eccles. xi. 8 — the Preacher to the young man: \emph{the days of darkness, for they shall be many} \textperiodcentered\ \textsuperscript{vii}\,John ix. 4 — at the healing of the man born blind: \emph{the night cometh, when no man can work} \textperiodcentered\ \textsuperscript{viii}\,Matt. xxv. 30 — the unprofitable servant, \emph{cast into outer darkness} \textperiodcentered\ \textsuperscript{ix}\,Luke xxiv. 29 — Emmaus: \emph{abide with us, for it is toward evening} \textperiodcentered\ \textsuperscript{x}\,John xiv. 16 — \emph{another Comforter}, the ἄλλος the plate keeps in Greek \textperiodcentered\ \textsuperscript{xi}\,1 John ii. 27 — \emph{the anointing which ye have received of him} \textperiodcentered\ \textsuperscript{xii}\,Eph. iv. 30 — \emph{the holy Spirit of God, whereby ye are sealed} \textperiodcentered\ \textsuperscript{xiii}\,2 Cor. v. 5 — \emph{the earnest of the Spirit}; and i. 22, the sealing and the earnest together\par}\par\vspace{3pt}\textsuperscript{x}\,\textbf{ἄλλος is left standing in Greek because the word is the argument.} John xiv. 16 promises \emph{another} Comforter, and Greek has a word for another-of-the-same-kind as against another-of-a-different-sort; this is the first, and the whole doctrine of the Spirit as a second Advocate beside the one ascended rests on it. Andrewes prints the bare adjective in Greek beside the Latin title and lets it carry the weight, which is why the English apposes it rather than absorbing it into \emph{another}. The four titles it heads — Paraclete, Unction, Seal, Earnest — are the New Testament's names for the Spirit set out as a schema, the same device as the \emph{Opera} brace above, and the evening office ends not in petition but in a list.}
\clearpage
\sethead{§7 \emph{Allegatio}}
% ---------- printed 291 ----------
\versoalign
\markright{291}
\setnoted{}
\addcontentsline{toc}{section}{§7 \emph{Allegatio} (the Pleading)}
\originalsone{291}{\input{fragments/la291}}{}
\clearpage
\markright{291}
\setnoted{15,17,19,20,21,22,23}
\enleafn{291}{\input{fragments/en291}}{}{{\scriptsize \textsuperscript{i}\,Ps. ciii. 8–12 — the great mercy-psalm, taken five verses together \textperiodcentered\ \textsuperscript{ii}\,Ps. lxxxvi. 5 — \emph{plenteous in mercy unto all them that call upon thee} \textperiodcentered\ \textsuperscript{iii}\,Ps. cxlv. 9 — \emph{his mercies are over all his works} \textperiodcentered\ \textsuperscript{iv}\,Mic. vii. 18 — the prophet's close: \emph{he delighteth in mercy} \textperiodcentered\ \textsuperscript{v}\,2 Cor. i. 3 — \emph{the Father of mercies, and the God of all comfort} \textperiodcentered\ \textsuperscript{vi}\,Ps. lix. 17 — \emph{God is my defence, and the God of my mercy} \textperiodcentered\ \textsuperscript{vii}\,Isa. xxviii. 21 — \emph{that he may do his work, his strange work}\par}\par\vspace{3pt}\textsuperscript{vii}\,\textbf{The plea's first ground is that punishment is not what God is for.} Isaiah says God rises up \emph{ut faciat opus suum, alienum opus suum} — to do his work, his strange work — and Andrewes braces the two halves so the asymmetry is visible on the page: shewing mercy is his \emph{own} work, punishing a work \emph{strange and foreign} to him. The whole \emph{Allegatio} is a plea at law, and this is the opening argument: not that the prisoner deserves mercy, but that mercy is the judge's proper business and severity his borrowed one.}
\clearpage
% ---------- printed 292 ----------
\versoalign
\markright{292}
\setnoted{}
\originalsone{292}{\input{fragments/la292}}{}
\clearpage
\markright{292}
\setnoted{2,3,5,7,9,10,11,13,14,18,23,24,25}
\enleafn{292}{\input{fragments/en292}}{}{{\scriptsize \textsuperscript{i}\,Num. xiv. 17, 18 — Moses pleading for Israel after the spies, when God would have disinherited them \textperiodcentered\ \textsuperscript{iii}\,John xx. 17 — the risen Christ to Mary Magdalene \textperiodcentered\ \textsuperscript{iv}\,Josh. vii. 9 — Joshua on his face after Ai: \emph{what wilt thou do unto thy great name?} \textperiodcentered\ \textsuperscript{v}\,John i. 29 — the Baptist, seeing Him come \textperiodcentered\ \textsuperscript{vi}\,Job xix. 25 — Job on the ash-heap, answering his comforters \textperiodcentered\ \textsuperscript{vii}\,John iv. 42 — the Samaritans of Sychar, after two days \textperiodcentered\ \textsuperscript{viii}\,1 John ii. 1 — \emph{if any man sin}, which is the condition the Advocate answers \textperiodcentered\ \textsuperscript{ix}\,Heb. vii. 25 — \emph{he ever liveth to make intercession for them} \textperiodcentered\ \textsuperscript{x}\,Matt. iii. 16 — the Baptism in Jordan \textperiodcentered\ \textsuperscript{xi}\,Ps. cxix. 49 — \emph{remember thy word unto thy servant} \textperiodcentered\ \textsuperscript{xii}\,Titus i. 2 — \emph{God, that cannot lie, promised before the world began} \textperiodcentered\ \textsuperscript{xiii}\,Heb. vi. 17 — the oath sworn to Abraham, God swearing by Himself for want of a greater\par}\par\vspace{3pt}\textsuperscript{ii}\,\textbf{These are not devotions but GROUNDS, and the difference is the whole section.} Heads [3]–[5] run through the Names of the Father, of Christ and of the Spirit — Lamb, Redeemer, Saviour, Mediator, Advocate, Intercessor, High Priest; Dove, Unction, Comforter — and each is entered as a reason why the plea must succeed. Two of them are forensic terms outright: an \emph{Advocate} is counsel and an \emph{Intercessor} files the motion. A reader who takes the list for a litany of titles will miss that Andrewes is arguing, and arguing from who his judge is rather than from anything he has to offer.}
\clearpage
% ---------- printed 293 ----------
\versoalign
\markright{293}
\setnoted{}
\originalsone{293}{\input{fragments/la293}}{}
\clearpage
\markright{293}
\setnoted{2,4,7,9,10,12,18,20,21,24}
\enleafn{293}{\input{fragments/en293}}{}{{\scriptsize \textsuperscript{i}\,Rom. iii. 3 — \emph{shall their unbelief make the faith of God without effect?} \textperiodcentered\ \textsuperscript{ii}\,2 Tim. ii. 13 — \emph{if we believe not, yet he abideth faithful} \textperiodcentered\ \textsuperscript{iii}\,Ps. xxii. 4 — the psalm of the Passion, where it turns from the sufferer to the fathers \textperiodcentered\ \textsuperscript{iv}\,Ps. xxv. 5 — \emph{thy tender mercies, which have been ever of old}; Prayer Book numbering \textperiodcentered\ \textsuperscript{v}\,Ps. lxxxix. 48 — \emph{Lord, where are thy old loving-kindnesses?}; Prayer Book numbering \textperiodcentered\ \textsuperscript{vi}\,Ecclus. ii. 11 — Ben Sira to the young man: \emph{look at the generations of old, and see} \textperiodcentered\ \textsuperscript{vii}\,Ps. cxxxviii. 8 — \emph{forsake not the works of thine own hands} \textperiodcentered\ \textsuperscript{viii}\,Isa. lxiv. 8 — \emph{we are the clay, and thou our potter} \textperiodcentered\ \textsuperscript{ix}\,Wisd. xi. 25 — \emph{thou hatest nothing that thou hast made} \textperiodcentered\ \textsuperscript{x}\,Gen. i. 26 — the sixth day: \emph{let us make man in our image, after our likeness}\par}\par\vspace{3pt}\textsuperscript{vi}\,\textbf{\textup{\emph{Syr.}} is Ecclesiasticus, and the gloss is the editor's.} The plate cites the book as \emph{Syr.} — for Sirach — and \textup{[i. e. \emph{Ecclus.}]} in square brackets is the 1853 editor supplying the familiar name, not Andrewes. The distinction is kept everywhere in this edition: printed in the plate, supplied by the editor, owed to a later scholar are three different weights of evidence, and only the first says anything about Andrewes.\\[1pt]\textsuperscript{vii}\,⚠ \textbf{The verse the Latin question turned on returns here, and reads differently.} At printed 34 the Latin renders this same \emph{Ps.} cxxxviii. 8 as \textbf{\emph{Figmenta} manuum tuarum} — a word that answers neither the printed Greek (τὰ ἔργα) nor the Vulgate (\emph{opera manuum tuarum}), and can only come from a Greek text reading πλάσμα, which is what the Wright apograph reads. It is the strongest single evidence that the 1853's Latin is not merely construed from its own Greek. Here in Part II the same verse is plain \textbf{\emph{Opus} manuum Tuarum}. The two are not in conflict — Part I sets a Latin against a Greek and Part II has no Greek to answer — but the contrast is worth a reader's eye, and it belongs to the introduction's case rather than to a correction of either page.}
\clearpage
% ---------- printed 294 ----------
\versoalign
\markright{294}
\setnoted{}
\originalsone{294}{\input{fragments/la294}}{}
\clearpage
\markright{294}
\setnoted{2,4,6,9,10,11,14,18,20,22,23,24}
\enleafn{294}{\input{fragments/en294}}{}{{\scriptsize \textsuperscript{i}\,Col. iii. 10 — the new man \emph{renewed in knowledge after the image of him that created him} \textperiodcentered\ \textsuperscript{ii}\,1 Cor. vi. 20 — \emph{ye are not your own, for ye are bought with a price} \textperiodcentered\ \textsuperscript{iii}\,1 Pet. i. 19 — \emph{a lamb without blemish and without spot} \textperiodcentered\ \textsuperscript{iv}\,Jer. xiv. 9 — the great drought, the prophet pleading: \emph{we are called by thy name; leave us not} \textperiodcentered\ \textsuperscript{v}\,Dan. ix. 19 — Daniel at the evening sacrifice, on the seventy years \textperiodcentered\ \textsuperscript{vi}\,Acts ix. 15 — the Lord to Ananias of Saul: \emph{he is a chosen vessel unto me} \textperiodcentered\ \textsuperscript{vii}\,1 Cor. vi. 15 — the argument against fornication, that the body is a member of Christ \textperiodcentered\ \textsuperscript{viii}\,Ps. cxix. 94 — \emph{I am thine, oh save me} \textperiodcentered\ \textsuperscript{ix}\,Ps. cxvi. 16 — \emph{I am thy servant, and the son of thine handmaid} \textperiodcentered\ \textsuperscript{x}\,Mark iv. 38 — the disciples wake Him in the storm: \emph{carest thou not that we perish?} \textperiodcentered\ \textsuperscript{xi}\,Luke xvii. 10 — \emph{we are unprofitable servants} \textperiodcentered\ \textsuperscript{xii}\,Luke xv. 24 — the father of the prodigal: \emph{this my son was dead, and is alive again}\par}\par\vspace{3pt}\textsuperscript{x}\,\textbf{A reproach turned into a plea, and answered in the same line.} The disciples' \emph{carest thou not that we perish?} is spoken in fear and something like accusation; Andrewes takes the question over as a ground of pleading and then answers it himself — \emph{Imo pertinet}, nay, Thou carest. He does the same with the storm at printed 286, where the verse serves the hour of prayer rather than the plea. The \emph{Preces} uses a text more than once and never for the same thing twice.\\[1pt]\textsuperscript{xi}\,\textbf{The whole head turns on a concessive, twice.} \emph{Servus inutilis, servus tamen. Filius perditus, tamen filius.} An unprofitable servant, yet a servant; a lost son, yet a son. Nothing is claimed about the pleader's worth — the argument is that the relation survives the failure, which is why the two texts chosen are the ones where a servant is called useless and a son is called dead. The English keeps \emph{yet} in both, because the plea is carried entirely by that word.}
\clearpage
% ---------- printed 295 ----------
\versoalign
\markright{295}
\setnoted{}
\originalsone{295}{\input{fragments/la295}}{}
\clearpage
\markright{295}
\setnoted{4,5,7,9,15,17,21,23}
\enleafn{295}{\input{fragments/en295}}{}{{\scriptsize \textsuperscript{i}\,Ps. vi. 2 — the first of the penitential psalms: \emph{have mercy upon me, for I am weak} \textperiodcentered\ \textsuperscript{ii}\,Ps. lxxxix. 47 — \emph{remember how short my time is} \textperiodcentered\ \textsuperscript{iii}\,Ps. lxxviii. 39 — of the generation in the wilderness, that God spared them \textperiodcentered\ \textsuperscript{iv}\,Ps. ciii. 14 — \emph{he knoweth whereof we are made} \textperiodcentered\ \textsuperscript{v}\,Ps. lxxix. 8 — after the temple was defiled: \emph{we are brought very low} \textperiodcentered\ \textsuperscript{vi}\,Ps. cvi. 44 — the psalm of Israel's repeated rebellions, at the point where He hears anyway \textperiodcentered\ \textsuperscript{vii}\,Ps. li. 17 — David's Miserere: \emph{a broken and a contrite heart} \textperiodcentered\ \textsuperscript{viii}\,Ps. xxxviii. 18 — \emph{I will declare mine iniquity; I will be sorry for my sin}\par}\par\vspace{3pt}\textsuperscript{ii}\,\textbf{The same Vulgate line as printed 34, doing different work.} There \emph{Memorare quæ mea substantia} was stacked against \emph{Substantia mea apud te est} four lines above it, and the whole join depended on the shared \emph{substantia} — a pun the Authorised Version dissolves, rendering one \emph{my hope is in thee} and the other \emph{remember how short my time is}. Here the verse stands alone, pleading the weakness of nature, and there is no join to protect; so the English can be plain where at 34 it had to keep the word.\\[1pt]\textsuperscript{iii}\,⚠ \textbf{A counter-instance worth keeping, on the question of which Psalter was to hand.} Printed 433's two references are right only on the Prayer Book's numbering, and this pass has since found nine more places where the plate counts with the BCP against the AV. This is not one of them: \emph{a wind that passeth away} is verse \textbf{39} in the Authorised Version and \textbf{40} in the Prayer Book Psalter, and the plate prints 39. The evidence is real but it is not uniform, and an introduction that claims the 1853 counts by the Prayer Book throughout will be answerable for this page.}
\clearpage
% ---------- printed 296 ----------
\versoalign
\markright{296}
\setnoted{}
\originalsone{296}{\input{fragments/la296}}{}
\clearpage
\markright{296}
\setnoted{2,3,4,5,6,10,15,18,20}
\enleafn{296}{\input{fragments/en296}}{}{{\scriptsize \textsuperscript{i}\,Ps. lxxxvi. 3 — \emph{I cry unto thee daily} \textperiodcentered\ \textsuperscript{ii}\,Ps. lxxx. 4 — \emph{how long wilt thou be angry at the prayer of thy servant?} \textperiodcentered\ \textsuperscript{iii}\,Matt. xviii. 32 — the lord to the unmerciful servant: \emph{I forgave thee all that debt} \textperiodcentered\ \textsuperscript{v}\,Luke vi. 37 — \emph{forgive, and ye shall be forgiven} \textperiodcentered\ \textsuperscript{vi}\,Mark xi. 25 — \emph{when ye stand praying, forgive} \textperiodcentered\ \textsuperscript{vii}\,Ps. cxix. 20 — \emph{my soul breaketh for the longing that it hath unto thy judgments} \textperiodcentered\ \textsuperscript{viii}\,Neh. i. 11 — Nehemiah praying before he speaks to the king \textperiodcentered\ \textsuperscript{ix}\,Luke xii. 47 — the servant that knew his lord's will and did it not\par}\par\vspace{3pt}\textsuperscript{ii}\,\textbf{The plate reads \emph{servi Tui}, and that is the Vulgate against both English Psalters.} Ps. lxxx. 4 has \emph{the prayer of thy people} in the Authorised Version and in the Prayer Book alike; the Vulgate has \emph{orationem servi tui}, thy servant's prayer, and Andrewes takes it. The difference is not decorative — the whole \emph{Allegatio} is one man pleading in the first person, and a plea about \emph{thy people} would loosen it at the point where it needs to be tightest. The English keeps \emph{thy servant}, which is CONVENTIONS §9 yielding where he reaches past the received Psalter.\\[1pt]\textsuperscript{iv}\,\textbf{The only ground in the plea that rests on something the pleader has done — and he argues the other side of it himself.} Every other head pleads God's nature, God's names, God's promise, or the pleader's weakness. This one pleads a condition met: \emph{because we forgive.} And the next lines print the reverse without softening it — \emph{but if ye forgive not, neither will your Father forgive you.} A plea at law that enters the strongest objection to its own case is doing something a devotion would not; the head is not a claim of merit but a reminder of terms.}
\clearpage
% ---------- printed 297 ----------
\versoalign
\markright{297}
\setnoted{}
\originalsone{297}{\input{fragments/la297}}{}
\clearpage
\markright{297}
\setnoted{1,3,5,13,16,19,21,23,27}
\enleafn{297}{\input{fragments/en297}}{}{{\scriptsize \textsuperscript{i}\,Ps. xxx. 9 — \emph{what profit is there in my blood?} \textperiodcentered\ \textsuperscript{ii}\,Ps. vi. 5 — \emph{in the grave who shall give thee thanks?} \textperiodcentered\ \textsuperscript{iii}\,Ps. lxxxviii. 10 — the darkest of the psalms, which ends without comfort \textperiodcentered\ \textsuperscript{iv}\,Isa. xxxviii. 18 — Hezekiah's writing, after his sickness and recovery \textperiodcentered\ \textsuperscript{v}\,Ps. lxxxix. 46 — Ethan: how short a time God has made men for \textperiodcentered\ \textsuperscript{vi}\,Ps. cxliii. 2 — \emph{enter not into judgment with thy servant} \textperiodcentered\ \textsuperscript{vii}\,Ps. cxxx. 3 — \emph{if thou, LORD, shouldest mark iniquities} \textperiodcentered\ \textsuperscript{viii}\,Job ix. 3 — Job answering Bildad: \emph{he cannot answer him one of a thousand} \textperiodcentered\ \textsuperscript{ix}\,Joel ii. 17 — the priests weeping between the porch and the altar\par}\par\vspace{3pt}\textsuperscript{i}\,\textbf{The head pleads from the silence of Sheol, and its title says how to read it.} Five texts in a row argue that the dead do not praise God, so that God Himself loses by the pleader's death — the grave cannot praise, the dust cannot declare the truth, they that go down into the pit cannot hope for it. This is the Old Testament's view of the dead, and Andrewes neither corrects it nor builds doctrine on it: the head is headed \emph{No profit [to God]}, and it is an argument advanced in a plea, not a statement of what he believes about the resurrection. Isaiah's answer stands two lines later in his own catena — \emph{the living, the living, he shall praise thee.}\\[1pt]\textsuperscript{vii}\,⚠ \textbf{The volume proves another of its own misprints, and this is the fourth time.} Here \textup{\emph{Psal.} cxxx. 3} is cited correctly for \emph{if thou shouldest mark iniquities}; at printed 40 the same verse is cited \textup{\emph{Psal.} cxxx. 30}, in a psalm that has eight verses. Both stand as printed. The pattern is now well established — printed 25 confirming the reading at 409, §17 confirming \textup{Job. xiv. 4} against Part III's \textup{xiv. 14}, §37 confirming \textup{Is. xxxviii. 15} — and it is worth more than any conjecture, because the corroborating witness is the book itself.}
\clearpage
% ---------- printed 298 ----------
\versoalign
\markright{298}
\setnoted{}
\originalsone{298}{\input{fragments/la298}}{}
\clearpage
\markright{298}
\setnoted{3,5,9,13,19,24}
\enleafn{298}{\input{fragments/en298}}{}{{\scriptsize \textsuperscript{i}\,Ps. lxxiv. 18 — the psalm written after the sanctuary was burned \textperiodcentered\ \textsuperscript{ii}\,Exod. xxxii. 12 — Moses after the golden calf: \emph{wherefore should the Egyptians speak?} \textperiodcentered\ \textsuperscript{iii}\,Num. xiv. 16 — Moses after the spies, foretelling what the nations will say \textperiodcentered\ \textsuperscript{iv}\,Ps. lxxix. 9 — \emph{help us, O God of our salvation, for the glory of thy name} \textperiodcentered\ \textsuperscript{v}\,Ps. li. 13 — David's Miserere, at its turn outward: \emph{then will I teach transgressors thy ways} \textperiodcentered\ \textsuperscript{vi}\,1 Tim. i. 16 — Paul on his own case: \emph{a pattern to them which should hereafter believe}\par}\par\vspace{3pt}\textsuperscript{ii}\,\textbf{Two pleas of Moses, and they are the Bible's own precedent for this whole argument.} After the golden calf and again after the spies, Moses turns God's threat aside not by defending Israel but by asking what Egypt and Canaan will say — that God brought them out for mischief, or could not bring them in. Andrewes sets the two side by side and lets them do the same work here: the evil that would follow is not the pleader's ruin but the reproach of the Name. It is the hinge of the \emph{Allegatio}, and the next head turns the same argument over — the \textbf{good} that would follow, thanksgiving, the conversion of others, and a pattern for those who come after.}
\clearpage
% ---------- printed 299 ----------
\versoalign
\markright{299}
\setnoted{}
\originalsone{299}{\input{fragments/la299}}{}
\clearpage
\markright{299}
\setnoted{2,4,5,6,7,9,11,13,15,16,18,23,24}
\enleafn{299}{\input{fragments/en299}}{}{{\scriptsize \textsuperscript{i}\,Isa. xliii. 25 — \emph{I, even I, am he that blotteth out thy transgressions for mine own sake} \textperiodcentered\ \textsuperscript{ii}\,Dan. ix. 19 — Daniel: \emph{defer not, for thine own sake, O my God} \textperiodcentered\ \textsuperscript{iii}\,Rom. iii. 25 — \emph{whom God hath set forth to be a propitiation through faith in his blood} \textperiodcentered\ \textsuperscript{iv}\,Ps. lxxxiv. 9 — \emph{behold, O God our shield, and look upon the face of thine anointed} \textperiodcentered\ \textsuperscript{v}\,Ps. cxxxii. 10 — \emph{turn not away the face of thine anointed} \textperiodcentered\ \textsuperscript{vi}\,Matt. xv. 22 — the woman of Canaan, crying after Him \textperiodcentered\ \textsuperscript{vii}\,2 Sam. xix. 23 — David sparing Shimei on the day he came back over Jordan \textperiodcentered\ \textsuperscript{viii}\,Isa. lxi. 1 — the text Christ read at Nazareth \textperiodcentered\ \textsuperscript{ix}\,Luke iv. 18 — and read it of Himself \textperiodcentered\ \textsuperscript{x}\,Matt. ix. 13 — at Matthew's table: \emph{I came not to call the righteous, but sinners} \textperiodcentered\ \textsuperscript{xi}\,John iii. 17 — \emph{God sent not his Son into the world to condemn the world} \textperiodcentered\ \textsuperscript{xii}\,John xvii. 5 — the high-priestly prayer, before the Passion \textperiodcentered\ \textsuperscript{xiii}\,Gen. xiv. 18 — Melchizedek, priest of the most high God\par}\par\vspace{3pt}\textsuperscript{vii}\,\textbf{From kindred: the plea addresses the Son of David, and then cites what David did.} The Canaanite woman's \emph{have mercy on me, O thou Son of David} names the kinship; the next text is David himself sparing Shimei — the man who had cursed him — on the day he came back over Jordan. The argument is that the son will do as the father did, and it is made entirely by putting the two verses next to each other. The head that follows, \emph{from office}, drops kinship for job description: the Spirit is upon Him \textbf{because} He was anointed to heal the broken-hearted, so mercy is not a favour asked but the work He was sent to do.}
\clearpage
\sethead{§8 \emph{Confessio Laudis}}
% ---------- printed 300 ----------
\versoalign
\markright{300}
\setnoted{}
\addcontentsline{toc}{section}{§8 \emph{Confessio Laudis}}
\originalsone{300}{\input{fragments/la300}}{}
\clearpage
\markright{300}
\setnoted{2,4,7,9,12,13,15,16,20,22}
\enleafn{300}{\input{fragments/en300}}{}{{\scriptsize \textsuperscript{i}\,Eccl. v. 8 — the Preacher, on oppression: \emph{higher than the highest regardeth} \textperiodcentered\ \textsuperscript{ii}\,Gen. xxi. 33 — Abraham at Beersheba, after the well was sworn for \textperiodcentered\ \textsuperscript{iii}\,Jer. xxiii. 24 — against the prophets who say they have dreamed \textperiodcentered\ \textsuperscript{iv}\,Ps. cxxxix. 7 — \emph{whither shall I go then from thy spirit?} \textperiodcentered\ \textsuperscript{v}\,John xxi. 17 — Peter, asked the third time: \emph{Lord, thou knowest all things} \textperiodcentered\ \textsuperscript{vi}\,1 Kings viii. 39 — Solomon at the dedication of the temple \textperiodcentered\ \textsuperscript{vii}\,Luke i. 37 — Gabriel to Mary, of her cousin Elisabeth \textperiodcentered\ \textsuperscript{viii}\,Gen. xvii. 1 — to Abram at ninety-nine years old \textperiodcentered\ \textsuperscript{ix}\,Rom. xi. 33 — where Paul breaks off the argument on Israel into praise \textperiodcentered\ \textsuperscript{x}\,Ps. cxvii. 2 — the shortest psalm, entire in two verses\par}\par\vspace{3pt}\textsuperscript{ii}\,\textbf{A schoolman's list of the attributes, proved entirely out of stories.} Eternity, omnipresence, omniscience, omnipotence, the depth of wisdom, unshaken truth — the headings could stand in any scholastic handbook, and not one of them is proved by a definition. Eternity is Abraham planting a grove at Beersheba; omniscience is Peter on the shore, asked a third time; omnipotence is Gabriel in a house in Nazareth; the God who knows all hearts is quoted from Solomon's prayer. The \emph{Confessio} argues like a schoolman and evidences like a preacher.\\[1pt]\textsuperscript{x}\,⚠ \textbf{The same numeral means two different psalms in this volume, and this page proves it.} Here \textup{\emph{Psal.} cxvii. 2} is right on the Authorised Version's numbering — Ps. cxvii has two verses and this is the second. At printed 274 the plate cites \textup{\emph{Ps.} cxvii. 25} for the Hosanna, which is impossible on that numbering and can only be the \textbf{Vulgate's} cxvii, our cxviii. Neither is an error; the 1853 simply does not hold to one Psalter's numbers, which is why every note in this edition that touches a psalm number says which system it means, and why the plate's own figure is never altered.}
\clearpage
% ---------- printed 301 ----------
\versoalign
\markright{301}
\setnoted{}
\originalsone{301}{\input{fragments/la301}}{}
\clearpage
\markright{301}
\setnoted{2,4,6,9,10,11,12,14,17,20,22}
\enleafn{301}{\input{fragments/en301}}{}{{\scriptsize \textsuperscript{i}\,Matt. xxiv. 35 — on the Mount of Olives, of the end \textperiodcentered\ \textsuperscript{ii}\,Ps. cxi. 3 — \emph{his righteousness endureth for ever} \textperiodcentered\ \textsuperscript{iii}\,Ps. xlii. 7 — \emph{deep calleth unto deep at the noise of thy waterspouts} \textperiodcentered\ \textsuperscript{iv}\,2 Cor. x. 1 — Paul, beginning his defence of himself to Corinth \textperiodcentered\ \textsuperscript{v}\,Gen. xviii. 32 — Abraham's last plea for Sodom, come down to ten \textperiodcentered\ \textsuperscript{vi}\,Mic. vii. 18 — the closing doxology of Micah's prophecy \textperiodcentered\ \textsuperscript{vii}\,Acts xvii. 30 — Paul on Mars' hill, of the times of ignorance \textperiodcentered\ \textsuperscript{viii}\,Rom. ii. 4 — to the man who judges another and does the same \textperiodcentered\ \textsuperscript{ix}\,Ps. lxxviii. 38 — of the generation in the wilderness, forgiven again \textperiodcentered\ \textsuperscript{x}\,Hos. vi. 4 — \emph{what shall I do unto thee?}, asked of both kingdoms \textperiodcentered\ \textsuperscript{xi}\,Neh. ix. 28, 30 — the Levites' confession, while the Law is read\par}\par\vspace{3pt}\textsuperscript{iii}\,\textbf{A verse about drowning, read of mercy.} \emph{Abyssus abyssum invocat} — deep calleth unto deep — is spoken in Ps. xlii by a man going under: the waterspouts, the waves and billows passing over him. Andrewes sets it beneath the heading \emph{The Fountain, the Ocean, the Deep of mercy}, so the two deeps become the depth of misery and the depth of pity, each calling to the other. The reading is not in the psalm and he does not argue for it; he places the verse and lets the heading do the work, which is how most of this catalogue proceeds.}
\clearpage
% ---------- printed 302 ----------
\versoalign
\markright{302}
\setnoted{}
\originalsone{302}{\input{fragments/la302}}{}
\clearpage
\markright{302}
\setnoted{2,4,6,8,11,13,16,17,20,21,23}
\enleafn{302}{\input{fragments/en302}}{}{{\scriptsize \textsuperscript{i}\,Ps. ciii. 10 — \emph{he hath not dealt with us after our sins} \textperiodcentered\ \textsuperscript{ii}\,Isa. xl. 2 — the opening of the book of comfort: \emph{speak ye comfortably to Jerusalem} \textperiodcentered\ \textsuperscript{iii}\,Ps. ciii. 13 — \emph{like as a father pitieth his own children} \textperiodcentered\ \textsuperscript{iv}\,Joel ii. 13 — \emph{rend your heart, and not your garments} \textperiodcentered\ \textsuperscript{v}\,Ps. ciii. 9 — \emph{he will not alway be chiding} \textperiodcentered\ \textsuperscript{vi}\,Matt. xviii. 32 — the lord to the unmerciful servant \textperiodcentered\ \textsuperscript{vii}\,2 Cor. v. 19 — \emph{not imputing their trespasses unto them} \textperiodcentered\ \textsuperscript{ix}\,Luke xv. 22, 23 — the father to his servants, before the son can finish his speech \textperiodcentered\ \textsuperscript{xi}\,Luke vi. 35 — \emph{he is kind unto the unthankful and to the evil}\par}\par\vspace{3pt}\textsuperscript{viii}\,\textbf{Propitiation is proved by a parable, and by the part of it where nobody speaks.} The heading is the most technical word in the catalogue, and under it Andrewes sets not a sacrifice but the best robe, the ring, the fatted calf — the father's instructions to his servants, given while the son is still reciting the confession he had prepared. What the heading calls propitiation the text shows as an interruption.\\[1pt]\textsuperscript{x}\,\textbf{Χρηστός is left in Greek because the word is one letter from the Name.} \emph{Chrestos} means kind, useful, serviceable, and it stands so close to \emph{Christos} that ancient writers inside and outside the Church confused the two on purpose and by accident. Andrewes prints the Greek as the heading and lets the Latin gloss follow; the English apposes it rather than translating it away, because \emph{Kind} alone loses the reason this attribute is set among the names of Christ at all.}
\clearpage
% ---------- printed 303 ----------
\versoalign
\markright{303}
\setnoted{}
\originalsone{303}{\input{fragments/la303}}{}
\clearpage
\markright{303}
\setnoted{2,3,4,14,15,17,19,21}
\enleafn{303}{\input{fragments/en303}}{}{{\scriptsize \textsuperscript{i}\,Matt. xx. 9 — the labourers hired at the eleventh hour, paid a whole day's wage \textperiodcentered\ \textsuperscript{ii}\,Luke xxiii. 43 — to the thief, on the day itself \textperiodcentered\ \textsuperscript{iv}\,Exod. xv. 11 — the Song of the Sea, sung on the far bank \textperiodcentered\ \textsuperscript{vi}\,Heb. i. 14 — \emph{are they not all ministering spirits, sent forth to minister?} \textperiodcentered\ \textsuperscript{vii}\,1 Thess. iv. 16 — \emph{with the voice of the archangel, and with the trump of God} \textperiodcentered\ \textsuperscript{viii}\,1 Pet. iii. 22 — \emph{angels and authorities and powers being made subject unto him}\par}\par\vspace{3pt}\textsuperscript{iii}\,\textbf{Fourteen works of mercy in seven pairs, set in two columns so that each is answered across the page.} The blind enlightened and the prisoners loosed; the naked clothed and the fallen raised; the scattered gathered and the living fed; the hungry sustained and the dead quickened. Every one is a scriptural participle with its own reference, and the whole table is a single sentence still governed by the \emph{Propter} four pages back — these are not petitions but reasons for praise. The English keeps the two columns and the references the plate prints inside them, so the tags below add nothing where the page already says everything.\\[1pt]\textsuperscript{v}\,\textbf{The confession turns here from what God does to whom He does it by.} \emph{Confessio in Angelis} — God glorified \textbf{in} His angels — and the nine orders follow in rank, each with its office: angels bear the charge of men, archangels announce the greater things, virtues work wonders, powers hold off the demons, principalities govern, dominions give, thrones judge. The governing \emph{in} is the whole point and is kept through every heading: the orders are not praised, God is praised in them.}
\clearpage
\sethead{§9 \emph{Confessio} [in Angelis et Sanctis]}
% ---------- printed 304 ----------
\versoalign
\markright{304}
\setnoted{}
\addcontentsline{toc}{section}{§9 \emph{Confessio} [in Angelis et Sanctis]}
\originalsone{304}{\input{fragments/la304}}{}
\clearpage
\markright{304}
\setnoted{2,3,4,6,14}
\enleafn{304}{\input{fragments/en304}}{}{{\scriptsize \textsuperscript{i}\,Col. i. 16 — \emph{thrones, or dominions, or principalities, or powers} \textperiodcentered\ \textsuperscript{iii}\,Gen. iii. 24 — the cherubim set at the east of the garden, with the flaming sword \textperiodcentered\ \textsuperscript{iv}\,Isa. vi. 2 — the year that king Uzziah died, in the temple\par}\par\vspace{3pt}\textsuperscript{ii}\,\textbf{Cherubim for knowledge, Seraphim for love, and the etymologies are the argument.} \emph{Cherub} was read as \emph{fulness of knowledge} and \emph{seraph} as \emph{the burning one}, so the two highest orders stand for the two things a soul wants — to know and to love — and the prayer that follows asks not to admire them but to be \textbf{joined unto their choirs}. The three proof-texts are chosen for the same reason: Paul names the ranks, Genesis shows a cherub set to guard, Isaiah shows a seraph burning above the throne.\\[1pt]\textsuperscript{v}\,\textbf{A double column: the order, and what that order gives.} Patriarchs and Faith, Prophets and Hope, Apostles and Labours, Martyrs and Blood, Confessors and Zeal, Doctors and Studies — and where the pairing fails the plate leaves the right-hand column empty rather than inventing (the Evangelists and the Ascetics have no gift set against them, and the Therapeutæ share the Tears with the line above by a brace). The empty cells are the 1853's and are kept. ⚠ \emph{Therapeutæ} are not physicians: the name is Philo's, for the contemplative ascetics of Egypt, and Andrewes sets them beside the Ascetics as a distinct order.}
\clearpage
% ---------- printed 305 ----------
\versoalign
\markright{305}
\setnoted{}
\originalsone{305}{\input{fragments/la305}}{}
\clearpage
\markright{305}
\setnoted{3,4,5}
\enleafn{305}{\input{fragments/en305}}{}{\textsuperscript{i}\,\textbf{From here the section is one sentence, and the notes below it thin out for that reason.} \emph{Confessio Laudum} is a single catalogue governed by the standing \emph{Pro} — For the Light, For the forming of Man, For the promise of the Seed — running from the first day of creation to the nurture of this particular man without a break in the grammar. The plate prints its own reference inside each cell of each brace and our English keeps them there, so a scripture band beneath would only repeat what the reader has already got, line for line. \textbf{An empty foot on these leaves is the right outcome, not a gap.}\\[1pt]\textsuperscript{ii}\,\textbf{The forming of man is given a brace inside the brace, and the inner one is the point.} Everything else in the list is a single item; man's making alone divides into \emph{habita deliberatione} and \emph{manibus suis} — of set counsel, and with His own hands. The counsel is \emph{let us make man} and the hands are the potter's, and Andrewes separates them because the one answers why and the other how. The cell that closes the brace has no reference at all and needs none: \emph{Peccante, non tamen derelicto} — having sinned, yet not forsaken.\\[1pt]\textsuperscript{iii}\,\textbf{Natural knowledge and revealed are set in one brace as two gifts of one giver.} The promise of the Seed is Gen. iii. 15, the first promise of the gospel; beside it stand \emph{that which may be known of God} and \emph{the work of the law written in their hearts}, both from the opening chapters of Romans, where they belong to the argument that leaves the heathen without excuse. Andrewes takes them out of that argument and enters them as mercies, then adds the oracles of the Prophets, the melody of the Psalms, the prudence of the Sentences and the experience of the Histories — so the ways God has made Himself known run from conscience to poetry without a rank among them.}
\clearpage
\sethead{§10 \emph{Confessio Laudum}}
% ---------- printed 306 ----------
\versoalign
\markright{306}
\setnoted{}
\addcontentsline{toc}{section}{§10 \emph{Confessio Laudum}}
\originalsone{306}{\input{fragments/la306}}{}
\clearpage
\markright{306}
\setnoted{2,3,4}
\enleafn{306}{\input{fragments/en306}}{}{\textsuperscript{i}\,\textbf{The technical words of the Incarnation, and two of them will stop a reader.} \emph{Exinanitio} is the Latin for St Paul's κένωσις — Christ \emph{emptied himself}, Phil. ii. 7 — and it is not self-abasement in general but a specific act, the laying aside of what was His by right. \emph{Apprehensio seminis Abraham} is Heb. ii. 16: He took hold not of angels but of the seed of Abraham, and \emph{apprehensio} keeps the grasping in it, the deliberate seizing of one nature and not another. Between them stands \emph{mysterium magnum pietatis}, the great mystery of godliness, which is 1 Tim. iii. 16's own name for the whole. The English keeps the technical weight rather than smoothing it, because the brace is a list of doctrines, not of feelings.\\[1pt]\textsuperscript{ii}\,\textbf{Two braces read across each other, and the second dates the first.} \emph{For all the good that He did and the evil that He suffered} — and then, in its own brace beside it, \emph{from the manger even to the cross.} The construction is not decoration: it says that the doing and the suffering are coextensive with the whole life, so there is no stretch of it that was not both. This is the \emph{Preces}' commonest device and its most easily lost — a brace answering a brace across the measure, which reads as two lists on a page and as one sentence in the mouth.\\[1pt]\textsuperscript{iii}\,\textbf{The economy in order, and the last three items are not events but dispositions.} The catalogue walks the infancy in sequence — Incarnation, poor Nativity, the manger, the Circumcision that subjected Him to the Law, the first-fruits of the blood, the lovable Name of JESUS, the showing to sinners of the Gentiles, the Presentation, the flight into Egypt — and then, numbered 1, 2, 3, it ends on \emph{eagerness to hear, avidity to seek, humility in obeying his parents}, all resting on the single verse that says He was subject unto them. \textbf{A list of what God did becomes a list of what a boy was like}, and the transition is unmarked.}
\clearpage
% ---------- printed 307 ----------
\versoalign
\markright{307}
\setnoted{}
\originalsone{307}{\input{fragments/la307}}{}
\clearpage
\markright{307}
\setnoted{3,5,7,8,10,11}
\enleafn{307}{\input{fragments/en307}}{}{{\scriptsize \textsuperscript{i}\,Heb. xii. 3 — \emph{consider him that endured such contradiction of sinners against himself} \textperiodcentered\ \textsuperscript{ii}\,Luke iv. 29 — at Nazareth, after He read Isaiah: they led Him to the brow of the hill \textperiodcentered\ \textsuperscript{iii}\,John viii. 59 — they took up stones; and x. 31, 33 — again, \emph{for blasphemy} \textperiodcentered\ \textsuperscript{v}\,John xviii. 40 — \emph{not this man, but Barabbas}\par}\par\vspace{3pt}\textsuperscript{iv}\,\textbf{The thanksgiving does not stop at the insults; it takes them in.} The standing \emph{Pro} is still governing here, so what the page renders is \emph{thanks for} being cast down headlong for a good sermon, stoned for a good work, called a Samaritan, a glutton, a demoniac, an impostor — and put after Barabbas. Andrewes does not soften the four names or gather them into a phrase about suffering: he lists them as they were said, in a brace, and files them among the mercies. \textbf{The English keeps each word and the brace, because a summary would turn a catalogue of what was shouted into a sentiment about rejection.}\\[1pt]\textsuperscript{vi}\,\textbf{Ten kinds of speech and act, in a column of pairs.} Sermons and homilies, colloquies and diatribes, intercessions and prayers — then, without a partner, examples, signs, mysteries, keys. The pairing is by register: public preaching against household teaching, conversation against sustained argument, pleading for others against praying. \emph{Diatriba} here is the schools' word for a discourse worked out at length, not invective. And the four unpaired items rise from what He said to what He did, showed, instituted and gave away — \emph{Claves}, the keys, being the last thing on the list because it is the thing He handed over.}
\clearpage
% ---------- printed 308 ----------
\versoalign
\markright{308}
\setnoted{}
\originalsone{308}{\input{fragments/la308}}{}
\clearpage
\markright{308}
\setnoted{4,5,7,10,11,12,13,15,17,19}
\enleafn{308}{\input{fragments/en308}}{}{{\scriptsize \textsuperscript{iii}\,John iii. 17 — to Nicodemus, who came by night \textperiodcentered\ \textsuperscript{iv}\,Matt. ix. 13 — at Matthew's table, to the Pharisees who asked why He ate with publicans \textperiodcentered\ \textsuperscript{v}\,Luke ix. 56 — after the Samaritan village would not receive Him, and the sons of thunder offered fire \textperiodcentered\ \textsuperscript{vi}\,Luke xix. 10 — to Zacchæus, in his house \textperiodcentered\ \textsuperscript{vii}\,Matt. xx. 28 — after the mother of Zebedee's children asked for the two seats \textperiodcentered\ \textsuperscript{viii}\,Matt. xi. 28 — after the woes on Chorazin and Bethsaida \textperiodcentered\ \textsuperscript{ix}\,Luke xxiii. 34 — from the cross, of the men driving the nails \textperiodcentered\ \textsuperscript{x}\,Matt. xx. 14 — the householder to the labourer who murmured at the eleventh hour\par}\par\vspace{3pt}\textsuperscript{i}\,\textbf{Six parables chosen on one principle.} The two debtors, the man half-dead, the Publican, the servant owing ten thousand talents, the lost sheep, the eleventh-hour labourer — every one turns on somebody who cannot pay, cannot walk, cannot plead, or came too late. The brace is not a survey of the parables but an argument assembled out of them, and it is the same argument the \emph{Allegatio} made at printed 288 from the eleventh hour.\\[1pt]\textsuperscript{ii}\,⚠ \textbf{The Sayings are all one saying, and they are stacked so the negative turns into the positive.} \emph{Not} to condemn the world — \emph{not} to judge it — \emph{not} to call the righteous — \emph{not} to destroy men's lives: four refusals in a row, each followed by what He came for instead, until \emph{the Son of man is come to seek and to save that which was lost} arrives without a negative at all and the series can end on the ransom. Andrewes prints no comment and needs none; the order does the work. \textbf{A later pass that tidies the repetitions of \emph{veni non… sed} destroys the figure}, exactly as tidying the divided Psalm xci at printed 431/433 would.}
\clearpage
% ---------- printed 309 ----------
\versoalign
\markright{309}
\setnoted{}
\originalsone{309}{\input{fragments/la309}}{}
\clearpage
\markright{309}
\setnoted{2,5,10,13,14,15,16,24,25,26}
\enleafn{309}{\input{fragments/en309}}{}{{\scriptsize \textsuperscript{i}\,Matt. xv. 28 — the woman of Canaan, after He had answered her not a word: \emph{O woman, great is thy faith} \textperiodcentered\ \textsuperscript{ii}\,Luke viii. 48 — the woman who touched the hem in the press, and was found out \textperiodcentered\ \textsuperscript{iii}\,Luke xxii. 62 — Peter, gone out and wept bitterly \textperiodcentered\ \textsuperscript{iv}\,John iii. 1–9 — Nicodemus, who came by night and asked how these things could be \textperiodcentered\ \textsuperscript{vi}\,Luke xv. 2 — the murmuring of the Pharisees and scribes, from which the words are taken \textperiodcentered\ \textsuperscript{viii}\,Phil. ii. 8 — \emph{obedient unto death, even the death of the cross} \textperiodcentered\ \textsuperscript{ix}\,Luke xxii. 15 — at the table: \emph{with desire I have desired to eat this passover with you} \textperiodcentered\ \textsuperscript{x}\,Matt. xxvi. 36 — Gethsemane, the garden over the brook; John xix. 13 — Gabbatha, the Pavement where Pilate sat to give judgment; and Golgotha, the place of a skull\par}\par\vspace{3pt}\textsuperscript{v}\,⚠ \textbf{A sneer is made the title of the catalogue, and it is printed in Greek to keep it a quotation.} \emph{Οὗτος δέχεται ἁμαρτωλούς} — \emph{this man receiveth sinners} — was said against Him by the Pharisees and scribes, murmuring. Andrewes sets it at the head of a list of everyone He received, so the complaint becomes the description, and he prints the Greek and the Latin as separate lines so both stand. That is CONVENTIONS §27's rule for a \textbf{separated} doublet, not §41's: the two lines are not a gloss for the reader but the same charge twice, and the English renders both.\\[1pt]\textsuperscript{vii}\,⚠ \textbf{The second half of the list is the same material as printed 307, doing the opposite work.} Those who gainsaid Him, would cast Him down headlong, twice took up stones, blasphemed, preferred Barabbas — every one of these stood two pages back among the mercies \emph{given thanks for}. Here they are entered as \textbf{examples}, that is, as instances of what He bore, and the catalogue runs on into the death itself. \textbf{The \emph{Preces} uses its material more than once and never for the same thing twice}, and the two passages must not be conformed or cross-referred in the text: the repetition is the design.}
\clearpage
\sethead{§11 \emph{Pro morte Christi}}
% ---------- printed 310 ----------
\versoalign
\markright{310}
\setnoted{}
\addcontentsline{toc}{section}{§11 \emph{Pro morte Christi}}
\originalsone{310}{\input{fragments/la310}}{}
\clearpage
\markright{310}
\setnoted{1,5,7,25}
\enleafn{310}{\input{fragments/en310}}{}{{\scriptsize \textsuperscript{i}\,Acts ii. 24 — \emph{the pains of death}, ὠδῖνες, which are birth-pangs; Heb. xii. 2 — \emph{despising the shame}; Gal. iii. 13 — \emph{made a curse for us} \textperiodcentered\ \textsuperscript{ii}\,Zech. xi. 12 — the shepherd's wages weighed out: \emph{so they weighed for my price thirty pieces of silver} \textperiodcentered\ \textsuperscript{iii}\,Mark xiii. 33 — the plate's figure; \emph{pavere et tædere} is verbatim the Vulgate of Mark \textbf{xiv}. 33 \textperiodcentered\ \textsuperscript{iv}\,Luke xxii. 53 — to those come out against Him: \emph{this is your hour, and the power of darkness}\par}\par\vspace{3pt}\textsuperscript{i}\,\textbf{Three words for the cross, in three languages, read across the page — and the first is a word about childbirth.} \emph{Dolor, Pudor, Maledictum} stand against \emph{Ὠδῖνες, Αἰσχύνη, Κατάρα}, each with its own proof-text: Acts for the pang, Hebrews for the shame, Galatians for the curse. \textbf{ὠδῖνες are birth-pangs}, the word for a woman in labour, and Peter uses it at Pentecost of the pains of death that could not hold Him. So the first of the three is not simply pain but travail — pain with an issue — which is why it can stand at the head of a catalogue that ends in resurrection. The English keeps \emph{Pang} and the Greek beside it rather than settling for \emph{sorrow}.\\[1pt]\textsuperscript{iii}\,⚠ \textbf{A wrong chapter, and it is the 1853 editor's, not Andrewes'.} The square brackets are his mark for a reference he supplied, and \textup{[\emph{Mar.} xiii. 33.]} sends the reader to \emph{take ye heed, watch and pray}. The words above it are \emph{Tædere} and \emph{Pavere}, and those two are the Vulgate of Mark \textbf{xiv}. 33 — \emph{coepit pavere et tædere} — verbatim and in that order, so the verse is not in doubt and only the figure is. The following \textup{[\emph{Ibid.}]} inherits the error. \textbf{Kept as printed}, like every other, and named here because a reader who checks it will find nothing.}
\clearpage
% ---------- printed 311 ----------
\versoalign
\markright{311}
\setnoted{}
\originalsone{311}{\input{fragments/la311}}{}
\clearpage
\markright{311}
\setnoted{3,12,16,20}
\enleafn{311}{\input{fragments/en311}}{}{{\scriptsize \textsuperscript{i}\,John xviii. 13 — Annas first, who was father in law to Caiaphas the high priest that year \textperiodcentered\ \textsuperscript{iii}\,Mark xiv. 65 — ῥαπίσματι, the blows of an open hand, kept in Greek because the Latin has no one word for it \textperiodcentered\ \textsuperscript{iv}\,Matt. xxvi. 55, 56 — the plate's figure; the condemnation for blasphemy is verses \textbf{65, 66}\par}\par\vspace{3pt}\textsuperscript{ii}\,⚠ \textbf{The catalogue stops being a catalogue for four lines, and prays.} \emph{Qui coram judice tacuisti, contine os meum} — Thou who wast silent before the judge, hold Thou my mouth; Thou who wast willing to be bound, hold Thou my hands. Each suffering is turned into a petition \textbf{for the member that suffered it}, so the silence answers the tongue and the cords answer the hands, and then the numbering resumes as though nothing had happened. This is how the whole section works and why it must not be levelled into a list: it is a meditation that keeps breaking into prayer at the point where the detail becomes unbearable.\\[1pt]\textsuperscript{iv}\,⚠⚠ \textbf{The 1853 editor's supplied references have an error rate of their own, and this is the second instance on two pages.} \textup{[\emph{Vers.} 55, 56.]} under \emph{blasphemiæ condemnari} points at \emph{are ye come out as against a thief} and \emph{all the disciples forsook him}; the condemnation for blasphemy is Matt. xxvi. \textbf{65, 66}. One page back, his \textup{[\emph{Mar.} xiii. 33.]} should be xiv. 33. \textbf{Both are single-digit slips, both inside his square brackets, and neither is Andrewes'.} This bears on how much weight a bracketed reference can carry anywhere in the volume: the three kinds of attribution are not of equal reliability either, and the editor's is the one now shown to be fallible twice in two leaves. Kept as printed.}
\clearpage
% ---------- printed 312 ----------
\versoalign
\markright{312}
\setnoted{}
\originalsone{312}{\input{fragments/la312}}{}
\clearpage
\markright{312}
\setnoted{1,9,15,20}
\enleafn{312}{\input{fragments/en312}}{}{{\scriptsize \textsuperscript{ii}\,Ps. xxii. 16 — \emph{they pierced my hands and my feet}, on the Authorised Version's numbering \textperiodcentered\ \textsuperscript{iii}\,Matt. xxvi. 46 — the plate's figure; the cry \emph{Eli, Eli} is at Matt. xxvii. \textbf{46} \textperiodcentered\ \textsuperscript{iv}\,Luke xxiii. 11 — Herod and his men of war, who arrayed Him in a gorgeous robe\par}\par\vspace{3pt}\textsuperscript{i}\,⚠⚠ \textbf{This catalogue also stands in PART I at printed 154, in the same Latin words} — \emph{Caput spinis coronatum · Oculi lacrymis perfusi · Aures conviciis oppletæ · Os felle et aceto potum · Cor lancea perforatum} — phrase for phrase, same order, same proof-texts. \textbf{What differs is the grammar, and with it the frame:} Part I sets the list ACCUSATIVE under a brace cued \emph{Per}, one long phrase \emph{pleading} the Passion; here it is NOMINATIVE and numbered, an inventory of what was suffered. \textbf{This is a Part I / Part II recension parallel, the first recorded} — the known ones run between Parts II and III. \textbf{Neither may be conformed:} Part II splits \emph{Manus et pedes perfossos} in two and adds the knees and the breast; Part I keeps the bloody sweat at the head and Part II drops it. The divergences are the evidence.\\[1pt]\textsuperscript{ii}\,⚠⚠ \textbf{And the two recensions cite the same verse by two different numbers.} Part I prints \textup{[\emph{Psal.} xxii. \textbf{17}.]} against the pierced hands and feet; this page prints \textup{[\emph{Psal.} xxii. \textbf{16}.]} against the same words. \textbf{Seventeen is the Prayer Book's figure, sixteen the Authorised Version's.} One catalogue, copied, with its psalm reference renumbered between recensions — so that the 1853 counts from more than one Psalter is shown not by inference across distant pages but from two printings of a single list. It is why the plate's figure stands wherever it is, with the system named.}
\clearpage
% ---------- printed 313 ----------
\versoalign
\markright{313}
\setnoted{}
\originalsone{313}{\input{fragments/la313}}{}
\clearpage
\markright{313}
\setnoted{7,8,11,18,20,21,23,26}
\enleafn{313}{\input{fragments/en313}}{}{{\scriptsize \textsuperscript{i}\,Luke xii. 50 — \emph{I have a baptism to be baptized with, and how am I straitened till it be accomplished} \textperiodcentered\ \textsuperscript{ii}\,Isa. liii. 5 — \emph{with his stripes we are healed} \textperiodcentered\ \textsuperscript{iii}\,John xix. 5 — Pilate bringing Him out in the purple: \emph{Ecce Homo} \textperiodcentered\ \textsuperscript{vi}\,Ps. xxii. 16 — again the Authorised Version's figure, as at printed 312 and against Part I's xxii. 17 \textperiodcentered\ \textsuperscript{vii}\,Mark xv. 28 — \emph{and he was numbered with the transgressors} \textperiodcentered\ \textsuperscript{viii}\,Matt. xxvii. 46 — \emph{Eli, Eli, lama sabachthani}\par}\par\vspace{3pt}\textsuperscript{iv}\,\textbf{A second column of one word each, naming what the suffering was.} Against \emph{to be stripped naked} stands \textbf{Shame}, against \emph{to be stretched on the cross} stands \textbf{Pain} — the same device as the \emph{Dolor · Pudor · Maledictum} of printed 310, and as the orders-and-gifts columns at 304. The 1853 sets these hard right and they are not glosses on the Latin but a separate register running down the page, so they must survive into both layers; dropped, the catalogue loses the only place where it says what it is doing.\\[1pt]\textsuperscript{v}\,⚠ \textbf{The nails are evidenced from AFTER the resurrection, and that is the method of the whole catalogue.} No Gospel narrates the act of nailing; so for \emph{clavis configi} Andrewes cites \textbf{John xx. 25} — the print of the nails, as Thomas demanded to see it eight days later — and for the hands and feet dug through he cites a psalm. The same thing was noted at printed 154, where the Passion is assembled out of five books and not one of them a Gospel account of it. \textbf{He is not narrating the Passion; he is proving it}, and a text that shows the scar serves better than one that would only describe the blow.}
\clearpage
% ---------- printed 314 ----------
\versoalign
\markright{314}
\setnoted{}
\originalsone{314}{\input{fragments/la314}}{}
\clearpage
\markright{314}
\setnoted{1,2,3,5,9}
\enleafn{314}{\input{fragments/en314}}{}{{\scriptsize \textsuperscript{ii}\,John xix. 34 — \emph{one of the soldiers with a spear pierced his side} \textperiodcentered\ \textsuperscript{iii}\,Matt. xxvii. 63 — the chief priests to Pilate the next day, of a man already buried\par}\par\vspace{3pt}\textsuperscript{i}\,⚠⚠ \textbf{The run closing here, printed 310–314, is a FOUR-GOSPEL HARMONY and the plate marks nearly every verse} — which is why those leaves could not be tagged one by one; 312 alone carries twenty-five references. \textbf{The four narratives are walked forward in step} — Matthew xxvi into xxvii, John xviii into xix, Luke xxii into xxiii, Mark xiv into xv. ⚠ \textbf{He is not following one evangelist and supplementing him; he is reading all four abreast.}\\[1pt]\textsuperscript{iv}\,⚠ \textbf{The list of sufferings ends on the one item that has no proof-text, and that is deliberate.} \emph{Ἄγνωστοι βάσανοι} — \textbf{torments unknown} — is entered as number 10 with no reference at all, in a catalogue where every other item down forty numbers carries one. Andrewes has assembled the whole Passion out of Scripture and then closes the account by admitting that Scripture does not hold it: what was recorded is finite, what was suffered is not. \textbf{The blank is the point; nothing may be supplied there}, and the Greek is kept because he reaches for it exactly where the Latin catalogue runs out.\\[1pt]\textsuperscript{v}\,\textbf{The wounds are then made to forgive, member for member.} The glorious head wounded forgives what was offended \emph{by the senses of the head}; the hands dug through forgive unlawful touch and unlawful working; the side pierced forgives the thoughts of the heart; the feet fastened forgive \emph{the going of feet swift to evil}. The Passion catalogue has become a confession keyed to the same anatomy, part answering part — the fuller form of the device that broke out for four lines at printed 311, where the silence before the judge answered the mouth and the cords answered the hands.}
\clearpage
% ---------- printed 315 ----------
\versoalign
\markright{315}
\setnoted{}
\originalsone{315}{\input{fragments/la315}}{}
\clearpage
\markright{315}
\setnoted{5,16,23,24}
\enleafn{315}{\input{fragments/en315}}{}{{\scriptsize \textsuperscript{iii}\,Col. ii. 15 — \emph{having spoiled principalities and powers, he made a shew of them openly, triumphing over them} \textperiodcentered\ \textsuperscript{iv}\,Phil. iii. 10 — \emph{that I may know him, and the power of his resurrection}\par}\par\vspace{3pt}\textsuperscript{i}\,⚠ \textbf{The prayer turns and the wounds change owner.} Through five members the wounded body has been forgiving; then \emph{Et ego, Domine, vulneratus sum animâ} — and I too am wounded, in soul — and the same four dimensions Paul used of the love of Christ are set to measure the injury instead: \emph{the multitude, the length, the breadth, the depth of them, from the head even to the heel.} The last clause is Isaiah's description of Israel, \emph{from the sole of the foot even unto the head there is no soundness in it}, turned on the man praying. It ends \emph{et per Tua sana mea} — \textbf{by Thy wounds heal mine} — which is the whole section in four words, and the reason the anatomy had to be kept exact on both sides.\\[1pt]\textsuperscript{ii}\,\textbf{A plea offered by proxy: he asks Christ to present the Passion to the Father on his behalf.} Seven items are named — the death, the opened side, the blood and water, the begging of the body, the taking down, the burial in another's tomb, the three days — and then, rather than pleading them himself, \emph{obsecro ut digneris hæc omnia pro me Patri Tuo offerre}. This is the \emph{Allegatio}'s method (printed 291–299) handed over: the grounds are the same, but the advocate is now asked to make the argument. And what is added at the end is not an event but a motive — \emph{the bitternesses which Thou didst suffer, the Love above all wherewith Thou didst suffer.}}
\clearpage
\sethead{§12 \emph{Transfiguratio}}
% ---------- printed 316 ----------
\versoalign
\markright{316}
\setnoted{}
\addcontentsline{toc}{section}{§12 \emph{Transfiguratio}}
\originalsone{316}{\input{fragments/la316}}{}
\clearpage
\markright{316}
\setnoted{2,3,4,5,6,9,13,15}
\enleafn{316}{\input{fragments/en316}}{}{{\scriptsize \textsuperscript{i}\,Acts i. 9 — \emph{a cloud received him out of their sight} \textperiodcentered\ \textsuperscript{ii}\,Acts ii. 33 — Peter at Pentecost: \emph{by the right hand of God exalted} \textperiodcentered\ \textsuperscript{iii}\,Eph. iv. 8 — \emph{he led captivity captive, and gave gifts unto men} \textperiodcentered\ \textsuperscript{iv}\,Heb. vii. 25 — \emph{he ever liveth to make intercession for them} \textperiodcentered\ \textsuperscript{v}\,Acts i. 11 — the two men in white apparel: \emph{shall so come in like manner} \textperiodcentered\ \textsuperscript{viii}\,Luke i. 35 — the Annunciation: \emph{the power of the Highest shall overshadow thee}\par}\par\vspace{3pt}\textsuperscript{vi}\,⚠ \textbf{A hymn set among the proof-texts, and quoted entire.} \emph{Veni Creator Spiritus, mentes tuorum visita, imple superna gratia quæ Tu creasti pectora} is the first stanza of the ninth-century hymn — the hymn of Pentecost, and the one sung at the laying on of hands in every ordination Andrewes ever took part in. In a section that has proved every article out of Scripture, the confession of the Holy Spirit opens instead in the Church's own words. \textbf{Our English renders it as verse-shaped prose like everything else here, and does not import an existing English translation of the hymn}: the \emph{Preces} is quoting a Latin text, and a familiar rendering would make the reader hear a hymnbook where the page has a prayer.\\[1pt]\textsuperscript{vii}\,\textbf{The Spirit's history is given an Old Testament first, and it begins at the second verse of Genesis.} Before the visible coming, the brooding upon the waters — \emph{the sending-forth} follows, and the whole Old Testament work is catalogued as its own brace before the New begins. ⚠ And a third register appears in the left margin from here: single words — \textbf{Shadow}, and its fellows below — naming the \emph{mode} of each coming, as \emph{Shame} and \emph{Pain} named the members' sufferings at printed 313. They are the plate's, they are structural, and they must not be read as part of the line they stand against.}
\clearpage
\sethead{§13 \emph{Confessio Laudum}}
% ---------- printed 317 ----------
\versoalign
\markright{317}
\setnoted{}
\addcontentsline{toc}{section}{§13 \emph{Confessio Laudum — Pro Spiritu Sancto}}
\originalsone{317}{\input{fragments/la317}}{}
\clearpage
\markright{317}
\setnoted{3,5,6,7,10,12,16,17,18,19,20}
\enleafn{317}{\input{fragments/en317}}{}{{\scriptsize \textsuperscript{i}\,Acts ii. 3, 4 — Pentecost: \emph{cloven tongues like as of fire} \textperiodcentered\ \textsuperscript{ii}\,Acts iv. 31 — after Peter and John were let go: \emph{the place was shaken where they were assembled} \textperiodcentered\ \textsuperscript{iii}\,Acts x. 44 — while Peter yet spake, the Holy Ghost fell on the Gentiles \textperiodcentered\ \textsuperscript{iv}\,Acts xix. 6 — the twelve at Ephesus who had not so much as heard whether there were a Holy Ghost \textperiodcentered\ \textsuperscript{vi}\,1 Cor. xii. 4–6 — \emph{diversities of gifts… differences of administrations… diversities of operations} \textperiodcentered\ \textsuperscript{vii}\,John xvi. 8 — \emph{he will reprove the world of sin} \textperiodcentered\ \textsuperscript{viii}\,1 John ii. 20 — \emph{ye have an unction from the Holy One} \textperiodcentered\ \textsuperscript{ix}\,John xiv. 26 — \emph{he shall bring all things to your remembrance} \textperiodcentered\ \textsuperscript{x}\,Rom. v. 5 — \emph{the love of God is shed abroad in our hearts by the Holy Ghost} \textperiodcentered\ \textsuperscript{xi}\,Rom. viii. 26 — \emph{the Spirit itself maketh intercession for us with groanings which cannot be uttered}\par}\par\vspace{3pt}\textsuperscript{v}\,\textbf{A Latin pair that only works because of the Greek beneath it.} \emph{Invocatio} answers \emph{Ἐπίκλησις} and \emph{advocatio} answers \emph{Παράκλησις}, and the two Greek words are the same verb — καλέω, to call — with different prefixes: a \textbf{calling-upon} and a \textbf{calling-alongside}. So the Spirit invoked and the Spirit as Advocate are one act described from two ends, which the Latin can imitate and the English cannot: \emph{invocation} and \emph{advocacy} no longer look like relatives in our language. The plate sets Greek and Latin on one line as a gloss for the reader, so by CONVENTIONS §41 the English renders it \textbf{once}, keeping the Greek visible with its sense in a parenthesis, rather than printing the same word twice.\\[1pt]\textsuperscript{vi}\,\textbf{Three threes, and the first is Paul's.} \emph{Gifts, ministrations, operations} is 1 Cor. xii. 4–6 entire, where each is assigned to a different Person — the same Spirit, the same Lord, the same God — and Andrewes uses it as the frame for what follows: the \textbf{gifts} of the Spirit, His \textbf{works}, His \textbf{fruits}, each then opened in its own list. The catalogue's shape is borrowed from the verse it cites, which is why the three bare words on their own lines are not a summary but a heading.}
\clearpage
% ---------- printed 318 ----------
\versoalign
\markright{318}
\setnoted{}
\originalsone{318}{\input{fragments/la318}}{}
\clearpage
\markright{318}
\setnoted{1,2,3,4,11,13,17,20}
\enleafn{318}{\input{fragments/en318}}{}{{\scriptsize \textsuperscript{i}\,Rom. viii. 16 — \emph{the Spirit itself beareth witness with our spirit, that we are the children of God} \textperiodcentered\ \textsuperscript{ii}\,2 Cor. i. 22 — \emph{who hath also sealed us, and given the earnest of the Spirit} \textperiodcentered\ \textsuperscript{iii}\,Eph. i. 14 — \emph{the earnest of our inheritance, until the redemption of the purchased possession} \textperiodcentered\ \textsuperscript{v}\,John xvi. 13 — \emph{he will guide you into all truth} \textperiodcentered\ \textsuperscript{vii}\,Phil. ii. 1 — \emph{if any consolation in Christ, if any bowels and mercies} \textperiodcentered\ \textsuperscript{viii}\,Eph. i. 7 — \emph{the riches of his grace}\par}\par\vspace{3pt}\textsuperscript{iv}\,\textbf{A second Latin word is offered for three of the seven works, and they are refinements rather than synonyms.} Against \emph{Visitatio} stands \emph{Invisere cor} — to visit \textbf{the heart}, naming what is visited; against \emph{Illustratio}, \emph{Illuminatio}; and against \emph{Perfectio}, \textbf{\emph{Perductio}}, which says where the perfecting arrives, a leading-through to the end rather than a completeness in itself. The other four stand bare. \textbf{This right-hand column is Latin glossing Latin}, unlike the \emph{Shadow · Dove · Breath} register two pages back, which named the mode of a coming; both are the plate's and neither belongs to the line beside it.\\[1pt]\textsuperscript{vi}\,\textbf{The \emph{Triumph of Mercy} is the \emph{Allegatio} again, with the grounds narrowed to one attribute.} The same forensic shape as printed 291–299 — numbered heads, each a reason the plea must succeed — except that every head is now some aspect of mercy itself: the Name and its glory, the promised truth and the oath interposed, the consolation of love, the bowels of mercies, and then mercy \textbf{great, manifold, of old, plenteous}, its riches and its overflow. ⚠ \textbf{The numbering runs continuously across the braces}, 1 to 13, so a brace here is a grouping inside one series and not a series of its own; a later pass that renumbers within a brace breaks the count.}
\clearpage
\sethead{§14 \emph{Triumphus misericordiæ}}
% ---------- printed 319 ----------
\versoalign
\markright{319}
\setnoted{}
\addcontentsline{toc}{section}{§14 \emph{Triumphus misericordiæ}}
\originalsone{319}{\input{fragments/la319}}{}
\clearpage
\markright{319}
\setnoted{1,3,4,5,7,11,13,17,21,23}
\enleafn{319}{\input{fragments/en319}}{}{{\scriptsize \textsuperscript{ii}\,1 Tim. i. 14 — \emph{the grace of our Lord was exceeding abundant} \textperiodcentered\ \textsuperscript{iii}\,Rom. v. 20 — \emph{where sin abounded, grace did much more abound} \textperiodcentered\ \textsuperscript{iv}\,Eph. ii. 7 — \emph{the exceeding riches of his grace in his kindness toward us} \textperiodcentered\ \textsuperscript{v}\,Jas. ii. 13 — \emph{mercy rejoiceth against judgment} \textperiodcentered\ \textsuperscript{vi}\,Lam. iii. 22 — \emph{it is of the LORD's mercies that we are not consumed} \textperiodcentered\ \textsuperscript{vii}\,Ps. xxiii. 6 — \emph{goodness and mercy shall follow me all the days of my life} \textperiodcentered\ \textsuperscript{ix}\,Ps. lxxxvi. 13 — \emph{thou hast delivered my soul from the lowest hell} \textperiodcentered\ \textsuperscript{x}\,Luke i. 78 — Zacharias: \emph{the tender mercy of our God}\par}\par\vspace{3pt}\textsuperscript{i}\,⚠⚠ \textbf{A ladder built out of Greek prefixes, and it cannot be climbed in English.} περισσεία, πλεονασμός, then \textbf{ὑπερ}πλεονασμός, \textbf{ὑπερ}περισσεία, \textbf{ὑπερ}βαλλόντως — abundance, increase, and then three compounds in which the same prefix is piled on, each taken from a place where St Paul himself reached for a word that did not exist until he made it: \emph{exceeding abundant}, \emph{much more abound}, \emph{the exceeding riches}. \textbf{The rhetoric is the repetition of ὑπερ-}, which English has to render four different ways — super-, over-, exceeding, surpassing — so the ladder flattens the moment the Greek is dropped. The plate keeps the Greek beside the Latin and \textbf{our English keeps both}, glossing rather than replacing; this is the clearest case in Part II for the edition's whole method.\\[1pt]\textsuperscript{viii}\,⚠ \textbf{Paul's four dimensions measure mercy here, and measured the WOUND at printed 315.} \emph{Μῆκος, Πλάτος, Βάθος} — length, breadth, depth — are Eph. iii. 18's measure of the love of Christ, and Andrewes fills each from the Psalter: from everlasting \textbf{to} everlasting for the length, up to heaven and down to hell for the depth, \emph{upon all} for the breadth. Four pages earlier the same figure was turned on the man praying — \emph{behold the multitude, the length, the breadth, the depth of them, from the head even to the heel} — to measure his hurt. \textbf{The same Pauline measure serves the injury and the mercy, and the two passages must not be conformed or cross-referred in the text.} It is the fourth time in this Part that a text is used twice and never for the same thing.}
\clearpage
% ---------- printed 320 ----------
\versoalign
\markright{320}
\setnoted{}
\originalsone{320}{\input{fragments/la320}}{}
\clearpage
\markright{320}
\setnoted{4,8,9,14,17,24,25,26}
\enleafn{320}{\input{fragments/en320}}{}{{\scriptsize \textsuperscript{i}\,Ecclus. ii. 18 — Ben Sira: \emph{as his majesty is, so is his mercy} \textperiodcentered\ \textsuperscript{ii}\,Matt. x. 29 — \emph{not one of them shall fall on the ground without your Father} \textperiodcentered\ \textsuperscript{iv}\,Luke x. 34 — the Samaritan, who bound up his wounds and set him on his own beast \textperiodcentered\ \textsuperscript{v}\,John xxi. 15 — on the shore, asked three times for the three denials \textperiodcentered\ \textsuperscript{vi}\,2 Cor. v. 20 — \emph{we pray you in Christ's stead, be ye reconciled to God} \textperiodcentered\ \textsuperscript{vii}\,Exod. xxv. 22 — the mercy seat between the cherubims: \emph{there I will meet with thee} \textperiodcentered\ \textsuperscript{viii}\,2 Cor. vi. 2 — \emph{now is the accepted time; now is the day of salvation}\par}\par\vspace{3pt}\textsuperscript{iii}\,⚠ \textbf{Twelve cases in a row, and every one is named by the thing that disqualified him.} Not Peter but \emph{Peter that denied}; not Thomas but \emph{Thomas the doubter}; Paul \emph{the blasphemer}, Magdalene \emph{the sinner}, the son \emph{wicked}, Jonah \emph{that fled}, the woman \emph{taken}, the thief. Where a person is not named the case is: the sheep that strayed, the coin that was lost, the debtor of ten thousand talents, the man left half dead. \textbf{The list could have been a roll of saints and is deliberately a roll of failures}, because the heading is \emph{Mercy} and the argument is what mercy has in fact done. Our English keeps every epithet; dropping one turns an argument back into a commemoration.\\[1pt]\textsuperscript{vii}\,\textbf{Mercy is given a place and a time, and the place is a piece of furniture.} \emph{Propitiatorium} is the \textbf{mercy seat} — the gold cover of the ark, between the cherubim, of which God said \emph{there I will meet with thee and I will commune with thee}; and Hebrews' \emph{throne of grace} is that same seat read as a throne. So the two halves of the brace are one object under two names, Old Testament and New. \textbf{The time is set against it in the same shape}: \emph{the day of salvation}, which in 2 Cor. vi. 2 is \emph{now}. Place and time together make the point of the whole \emph{Triumphus} — that there is somewhere to go and it is open.}
\clearpage
\sethead{§15 \emph{Præoratio}}
% ---------- printed 321 ----------
\versoalign
\markright{321}
\setnoted{}
\addcontentsline{toc}{section}{§15 \emph{Præoratio} + \emph{Deprecatio}}
\originalsone{321}{\input{fragments/la321}}{}
\clearpage
\markright{321}
\setnoted{2,3,8,12,15}
\enleafn{321}{\input{fragments/en321}}{}{{\scriptsize \textsuperscript{i}\,Acts x. 4 — the angel to Cornelius: \emph{thy prayers and thine alms are come up for a memorial before God} \textperiodcentered\ \textsuperscript{ii}\,2 Chr. xxx. 27 — after Hezekiah's passover: \emph{their prayer came up to his holy dwelling place, even unto heaven} \textperiodcentered\ \textsuperscript{iv}\,Lam. i. 12 — \emph{is it nothing to you, all ye that pass by?} \textperiodcentered\ \textsuperscript{v}\,Dan. ix. 19 — \emph{O Lord, hearken and do; defer not}\par}\par\vspace{3pt}\textsuperscript{i}\,\textbf{Two chains of six, and the second answers the first verb for verb.} The prayer is to \textbf{ascend, reach, enter, appear, find grace, draw near}; then God is asked to \textbf{hear, incline, attend and consider, understand, hearken, remember to do}. Every one of the twelve carries its own proof-text, so the movement from the mouth to the ear is built entirely out of other men's prayers — Cornelius's, Hezekiah's, Daniel's, and the Psalter's. ⚠ \textbf{The second chain does not end at hearing.} \emph{Memento ut facias} — remember to do it — is the sixth term, and it is Daniel's; the sequence is not complete until the hearing has become an act.\\[1pt]\textsuperscript{iii}\,⚠ \textbf{A promise about God's word is turned on the man's own prayer.} \emph{Non redeat ad me vacua} is Isaiah lv. 11, where it is spoken of the word that \textbf{goes out from God's mouth} and shall not return unto Him void. Andrewes asks it of the prayer going up from his — the traffic reversed, the guarantee borrowed. He does not argue for the transfer and the English does not soften it; the clause that follows, \emph{but as Thou knowest, and canst, and wilt}, is what he puts in place of an argument.}
\clearpage
% ---------- printed 322 ----------
\versoalign
\markright{322}
\setnoted{}
\originalsone{322}{\input{fragments/la322}}{}
\clearpage
\markright{322}
\setnoted{4,8,9,10,11,12,13,14}
\enleafn{322}{\input{fragments/en322}}{}{{\scriptsize \textsuperscript{ii}\,Ps. lxxxix. 47 — \emph{remember how short my time is}, the third time this verse anchors a plea \textperiodcentered\ \textsuperscript{iii}\,Gen. xviii. 27 — Abraham pleading for Sodom: \emph{I have taken upon me to speak, which am but dust and ashes} \textperiodcentered\ \textsuperscript{iv}\,Ps. xxxix. 12 — \emph{I am a stranger with thee, and a sojourner} \textperiodcentered\ \textsuperscript{v}\,Job iv. 19 — Eliphaz: \emph{them that dwell in houses of clay, whose foundation is in the dust} \textperiodcentered\ \textsuperscript{vi}\,Gen. xlvii. 9 — Jacob to Pharaoh: \emph{few and evil have the days of the years of my life been} \textperiodcentered\ \textsuperscript{vii}\,Prov. xxvii. 1 — \emph{boast not thyself of to morrow, for thou knowest not what a day may bring forth} \textperiodcentered\ \textsuperscript{viii}\,Isa. xxxviii. 13 — Hezekiah in his sickness: \emph{from day even to night wilt thou make an end of me}\par}\par\vspace{3pt}\textsuperscript{i}\,\textbf{The whole warding-off ends on one antithesis, and it is built to be lopsided.} \emph{Vincatur Leo ab Agna infirma, spiritus violentus a carne debili} — let the Lion be conquered by the weak Sheep, the violent spirit by the feeble flesh. The adversary is named by strength and the petitioner by weakness, twice over, and the verb is passive in both halves: nothing is asked to \emph{fight}, only that the stronger be overcome. \textbf{The English keeps the diminutive and the adjectives} — \emph{weak Sheep}, \emph{feeble flesh} — because the disproportion is the prayer.\\[1pt]\textsuperscript{iv}\,⚠⚠ \textbf{A second case of one verse cited by two different numbers in the two Parts.} \emph{I am a stranger with thee, and a sojourner} is \textup{\emph{Psal.} xxxix. \textbf{12}} here and \textup{\emph{Psal.} xxxix. \textbf{14}} at printed 261. Twelve is the Authorised Version's figure exactly; fourteen answers to no system cleanly, the Prayer Book making it thirteen, so Part I's may be a slip of one rather than another Psalter. \textbf{The pairing is what matters}, and it is the same phenomenon as the pierced hands and feet at printed 154 and 312, cited xxii. 17 and xxii. 16 — there unambiguously Prayer Book against Authorised Version. \textbf{Two independent instances now show the two Parts numbering the same psalm differently}, which is why no figure in this edition is ever reconciled to another.}
\clearpage
\sethead{§16 \emph{Allegatio}}
% ---------- printed 323 ----------
\versoalign
\markright{323}
\setnoted{}
\addcontentsline{toc}{section}{§16 \emph{Allegatio} (fragilitatis humanæ)}
\originalsone{323}{\input{fragments/la323}}{}
\clearpage
\markright{323}
\setnoted{2,4,12,14,15,16,17,18,21,26}
\enleafn{323}{\input{fragments/en323}}{}{{\scriptsize \textsuperscript{i}\,Rom. vi. 6 — \emph{that the body of sin might be destroyed} \textperiodcentered\ \textsuperscript{ii}\,1 John v. 19 — \emph{the whole world lieth in wickedness} \textperiodcentered\ \textsuperscript{iii}\,Job xiv. 4 — \emph{who can bring a clean thing out of an unclean? not one} \textperiodcentered\ \textsuperscript{iv}\,Ps. li. 5 — David's Miserere: \emph{in sin did my mother conceive me} \textperiodcentered\ \textsuperscript{v}\,Heb. xii. 15 — \emph{lest any root of bitterness springing up trouble you} \textperiodcentered\ \textsuperscript{vi}\,Rom. xi. 17 — the wild olive graffed in among the branches \textperiodcentered\ \textsuperscript{vii}\,Ps. cvi. 6 — \emph{we have sinned with our fathers, we have committed iniquity} \textperiodcentered\ \textsuperscript{viii}\,1 Kings viii. 47 — Solomon's prayer, putting the words in the mouth of a people in captivity \textperiodcentered\ \textsuperscript{ix}\,Isa. lxiii. 10 — \emph{they rebelled, and vexed his holy Spirit} \textperiodcentered\ \textsuperscript{x}\,Jer. vii. 13 — \emph{I called you, but ye answered not}\par}\par\vspace{3pt}\textsuperscript{iii}\,⚠ \textbf{This is the reference the volume elsewhere gets wrong, and here it is right.} \textup{\emph{Job} xiv. 4} — \emph{who can bring a clean thing out of an unclean?} — stands correctly on this page; at printed 406, where Part III makes its own confession out of the same material, the plate prints \textup{\emph{Job.} xiv. \textbf{14}}, which is \emph{if a man die, shall he live again?} Both stand as printed and neither is mended. \textbf{The corroborating witness is the book itself}, which is worth more than any conjecture, and it is now the fourth such case — with \textup{Is. xxxiii. 15} against §37's \textup{xxxviii. 15}, \textup{Psal. cxxx. 30} against printed 297's \textup{cxxx. 3}, and printed 25 confirming the reading at 409.\\[1pt]\textsuperscript{vii}\,\textbf{The confession is made entirely in other men's words.} Eleven numbered admissions follow, and not one is Andrewes' own sentence: \emph{I have sinned, I have done unjustly} is the Psalter's, \emph{I have dealt wickedly in the covenant} is Solomon putting words into the mouth of a people not yet carried away, \emph{I have grieved the Spirit} is Isaiah's account of the wilderness generation, \emph{I have gone after my own thoughts} is Jeremiah quoting the people's defiance. \textbf{Each is a third-person or plural indictment turned into the first person singular} — the accusation accepted rather than answered. It is the same method as the Passion catalogue at printed 313: he is not describing his case, he is proving it, and the proof-texts are the confession.}
\clearpage
\sethead{§17 \emph{Peccatorum Confessio}}
% ---------- printed 324 ----------
\versoalign
\markright{324}
\setnoted{}
\addcontentsline{toc}{section}{§17 \emph{Peccatorum Confessio}}
\originalsone{324}{\input{fragments/la324}}{}
\clearpage
\markright{324}
\setnoted{1,2,3,4,5,20}
\enleafn{324}{\input{fragments/en324}}{}{{\scriptsize \textsuperscript{i}\,Heb. iii. 13 — \emph{lest any of you be hardened through the deceitfulness of sin} \textperiodcentered\ \textsuperscript{ii}\,1 Kings xvi. 33 — Ahab, who did more to provoke the LORD than all before him; Ps. x. 13 — \emph{wherefore doth the wicked contemn God?} \textperiodcentered\ \textsuperscript{iii}\,Lam. iii. 59 — \emph{O LORD, thou hast seen my wrong: judge thou my cause} \textperiodcentered\ \textsuperscript{iv}\,Ps. l. 21 — \emph{these things hast thou done, and I kept silence}\par}\par\vspace{3pt}\textsuperscript{iii}\,⚠ \textbf{The fourteenth term of the confession is not a sin but a silence.} Thirteen admissions run down the page in the first person — I have hardened myself, I have provoked Thee — and the catalogue ends by turning to what God did: \emph{And all this Thou hast seen. And Thou hast held Thy peace.} \textbf{Both halves are quoted, and both are dangerous.} Jeremiah asks God to judge \textbf{because} He has seen; the psalm's silence is followed, in the same verse, by \emph{thou thoughtest that I was altogether such an one as thyself} — the forbearance mistaken for consent. So the list of sins closes on the one thing in it that is God's act, and leaves it undecided whether the silence is mercy or evidence. Nothing is added, and the section stops here.\\[1pt]\textsuperscript{v}\,\textbf{From here the page stops praying and starts weighing.} \emph{Aggravatio} is the confessor's technical term, and the nine heads are the \textbf{circumstances} by which a sin's gravity is reckoned in the manuals — its measure, quality, iteration, continuance, the person, the manner, the moving cause, the time, the place — the \emph{quis, quid, ubi, quibus auxiliis, cur, quomodo, quando} that the schoolmen took from the rhetoricians' analysis of an act and the confessors took from the schoolmen. ⚠ \textbf{Not one of the nine carries a proof-text, and the scripture band below is empty for the rest of this leaf; that is the right outcome, not a gap} — the page is an instrument, and an instrument is not argued from Scripture. The four that follow, \emph{folly, ingratitude, hardness, contempt}, are aggravations of another kind: not circumstances of the deed but the mind it was done in.\\[1pt]\textsuperscript{vi}\,\textbf{The kinds of sin, and the pairs are the argument.} Each numbered head sets two terms against each other inside a brace and prints its own reference in the cell, so the English keeps them there and the band below has nothing to add. \textbf{Isaiah's pair is one image, not two sins}: \emph{woe unto them that draw iniquity with cords of vanity, and sin as it were with a cart rope} — the thread first and the cable after, the same haulage grown strong enough to be seen. The second pair divides sins of \textbf{need} from sins of \textbf{excess} — the Vulgate's \emph{necessitates cordis mei} against the \emph{superfluity of naughtiness} in James — which is why what looks like a list of sixteen kinds is really a list of sixteen antitheses.}
\clearpage
\sethead{§18 \emph{Aggravatio}}
% ---------- printed 325 ----------
\versoalign
\markright{325}
\setnoted{}
\addcontentsline{toc}{section}{§18 \emph{Aggravatio} + \emph{Peccati Species}}
\originalsone{325}{\input{fragments/la325}}{}
\clearpage
\markright{325}
\setnoted{1,10,11}
\enleafn{325}{\input{fragments/en325}}{}{\textsuperscript{i}\,\textbf{The sixteen kinds are sixteen antitheses, and the brace is the argument.} Every head sets two terms against each other — omission against commission, inward against outward, ignorant against knowing, once against often, hidden against manifest — so what looks like a list of sins is a set of axes along which one sin can be placed. \textbf{The plate prints its references inside the cells and the English keeps them there, so the scripture band stays empty for this leaf.} The first head takes Paul for both its halves at once: \emph{what I would, that do I not; but what I hate, that do I} is a single sentence containing the not-doing and the doing, which is why one reference serves a pair.\\[1pt]\textsuperscript{ii}\,\textbf{Anger and concupiscence are not two sins here but the two appetites.} The division is the schoolmen's \emph{irascibilis} and \emph{concupiscibilis}, the passionate part of the soul cut in two, so head twelve is not naming wrath and lust among the rest but sorting every sin by which appetite carried it. The inner brace then divides concupiscence itself, \emph{of the flesh} from \emph{of the world}.\\[1pt]\textsuperscript{iii}\,⚠ \textbf{The gravest distinction on the page carries no reference, because it needed none.} \emph{By one not yet called} against \emph{by one already called} — sin before the calling and sin after it — was a commonplace of the manuals and of the epistles that speak of those once enlightened. It is set down as flatly as the difference between sleeping and waking two lines below, and the plate leaves the reader to supply what hangs on it.}
\clearpage
% ---------- printed 326 ----------
\versoalign
\markright{326}
\setnoted{}
\originalsone{326}{\input{fragments/la326}}{}
\clearpage
\markright{326}
\setnoted{14,15,18,21,23,24,25,27}
\enleafn{326}{\input{fragments/en326}}{}{{\scriptsize \textsuperscript{i}\,2 Cor. vii. 11 — the marks of godly sorrow, which both columns of this page are numbering \textperiodcentered\ \textsuperscript{ii}\,Rom. vi. 21 — \emph{what fruit had ye then in those things whereof ye are now ashamed?} \textperiodcentered\ \textsuperscript{iii}\,Matt. xi. 28 — \emph{come unto me, all ye that labour and are heavy laden} \textperiodcentered\ \textsuperscript{iv}\,1 Cor. xi. 31 — \emph{if we would judge ourselves, we should not be judged} \textperiodcentered\ \textsuperscript{v}\,Jude 23 — \emph{hating even the garment spotted by the flesh} \textperiodcentered\ \textsuperscript{vi}\,Gal. v. 13 — \emph{use not liberty for an occasion to the flesh} \textperiodcentered\ \textsuperscript{vii}\,Jas. iv. 10 — \emph{humble yourselves in the sight of the Lord, and he shall lift you up} \textperiodcentered\ \textsuperscript{viii}\,Jer. xxxi. 19 — Ephraim, after his turning: \emph{I smote upon my thigh; I was ashamed}\par}\par\vspace{3pt}\textsuperscript{i}\,⚠⚠ \textbf{Five impersonal verbs under a sixth, and then the grammar changes.} The section is headed \emph{Pœnitet} — \emph{it repenteth me} — and the five under it are impersonals too: \emph{Dolet, Pudet, Piget, Horret, Tædet}, it grieves, it shames, it vexes, it horrifies, it wearies. \textbf{In Latin the man is the subject of none of them.} Repentance, in the first list, is something that happens to him, and each is given the Greek noun of the affection and the thing it is felt for — sorrow for the wounds, shame for the stains, indignation for the guilt, fear for the punishment, weariness of the yoke. \textbf{Then the second list turns to nouns of action} — self-judgement, self-avenging, hatred, flight from occasions, self-humbling, the smiting of the breast and of the thigh — and every one is a thing he does. The page moves from what repentance suffers to what repentance performs, and it makes the turn in the grammar before it makes it in the words.}
\clearpage
\sethead{§19 \emph{Pœnitet}}
% ---------- printed 327 ----------
\versoalign
\markright{327}
\setnoted{}
\addcontentsline{toc}{section}{§19 \emph{Pœnitet}}
\originalsone{327}{\input{fragments/la327}}{}
\clearpage
\markright{327}
\setnoted{1,2,4,11,14,16,17,19,20,22,25,27}
\enleafn{327}{\input{fragments/en327}}{}{{\scriptsize \textsuperscript{i}\,Jer. iv. 8 — \emph{gird you with sackcloth, lament and howl} \textperiodcentered\ \textsuperscript{ii}\,Joel ii. 12 — \emph{turn ye even to me with all your heart, with fasting, and with weeping} \textperiodcentered\ \textsuperscript{iii}\,Prov. xvi. 6 — \emph{by mercy and truth iniquity is purged} \textperiodcentered\ \textsuperscript{iv}\,Ps. lxxix. 8 — \emph{let thy tender mercies speedily prevent us: for we are brought very low} \textperiodcentered\ \textsuperscript{v}\,Rom. viii. 32 — \emph{how shall he not with him also freely give us all things?} \textperiodcentered\ \textsuperscript{vi}\,Acts vii. 60 — Stephen, kneeling: \emph{lay not this sin to their charge} \textperiodcentered\ \textsuperscript{vii}\,Ps. xxxii. 2 — \emph{blessed is the man unto whom the LORD imputeth not iniquity} \textperiodcentered\ \textsuperscript{viii}\,Ps. cxxx. 3 — \emph{if thou, LORD, shouldest mark iniquities, O Lord, who shall stand?} \textperiodcentered\ \textsuperscript{ix}\,Ps. cxliii. 2 — \emph{enter not into judgment with thy servant} \textperiodcentered\ \textsuperscript{x}\,Ps. li. 11 — \emph{cast me not away from thy presence} \textperiodcentered\ \textsuperscript{xi}\,Exod. xxxii. 12 — Moses on the mount: \emph{turn from thy fierce wrath, and repent of this evil} \textperiodcentered\ \textsuperscript{xii}\,Ps. lxxvii. 7 — \emph{will the Lord absent himself for ever? will he be no more entreated?}\par}\par\vspace{3pt}\textsuperscript{vii}\,⚠ \textbf{Three Greek verbs stand beside the Latin here, and they are three different ways of not counting a sin.} \emph{Μὴ στήσῃς} is Stephen's, dying — do not \textbf{set} it up against me; \emph{μὴ λογίσῃ} is the Septuagint's, in the psalm Paul builds imputation on — do not \textbf{reckon} it to my account; \emph{μὴ μνησθῇς} — do not \textbf{remember} it. Judicial, mercantile, personal: \textbf{the English of the three is nearly one sentence and the Greek of them is three distinct acts}, which is why the plate keeps the Greek. ⚠ Stephen prayed his for \textbf{other men}; it is asked here for the man praying.\\[1pt]\textsuperscript{xi}\,\textbf{Nine petitions in a row ask God not to act, and then the prayer turns and asks what He is.} From \emph{hold back thine anger} to \emph{cast me not away}, every clause is a negative. The turn comes at \emph{Esto} — \textbf{be Thou} — with three adjectives that rhyme in the Latin and cannot in the English: \emph{placabilis, præstabilis, exorabilis}. \textbf{They are adjectives of capacity, not of action:} the prayer stops asking God to do anything and asks that He be the kind of God who can be moved.}
\clearpage
\sethead{§20 \emph{Miserere}}
% ---------- printed 328 ----------
\versoalign
\markright{328}
\setnoted{}
\addcontentsline{toc}{section}{§20 \emph{Miserere}}
\originalsone{328}{\input{fragments/la328}}{}
\clearpage
\markright{328}
\setnoted{7,9,10,11,12,15,18,19,21,26}
\enleafn{328}{\input{fragments/en328}}{}{{\scriptsize \textsuperscript{ii}\,Job xiv. 16 — the Vulgate's \emph{parce peccatis meis}, where the AV has \emph{dost thou not watch over my sin?} \textperiodcentered\ \textsuperscript{iii}\,John i. 29 — the Baptist's \emph{ὁ αἴρων}: \emph{behold the Lamb of God, which taketh away the sin of the world} \textperiodcentered\ \textsuperscript{iv}\,2 Sam. xii. 13 — Nathan to David: \emph{the LORD also hath put away thy sin} \textperiodcentered\ \textsuperscript{v}\,Isa. xxxviii. 17 — Hezekiah's song in his sickness: \emph{thou hast cast all my sins behind thy back} \textperiodcentered\ \textsuperscript{vi}\,Joel ii. 14 — \emph{who knoweth if he will return and repent, and leave a blessing behind him?} \textperiodcentered\ \textsuperscript{vii}\,2 Cor. v. 18 — \emph{who hath reconciled us to himself by Jesus Christ} \textperiodcentered\ \textsuperscript{viii}\,1 Kings viii. 39 — Solomon at the dedication: \emph{hear thou in heaven, and forgive} \textperiodcentered\ \textsuperscript{ix}\,Ps. li. 9 — \emph{hide thy face from my sins} \textperiodcentered\ \textsuperscript{x}\,Ps. lxxxv. 2 — \emph{thou hast covered all their sin}\par}\par\vspace{3pt}\textsuperscript{i}\,\textbf{The page divides sin into four things it does, and gives each its own ladder.} \emph{Reatus, macula, morbus, servitus} — guilt, stain, sickness, bondage, the scholastic four-fold effect — and each carries its Greek term: guilt is \textbf{ὑπόδικον}, liable to judgement, a courtroom word; stain is \textbf{κηλίς}, a spot, with \textbf{δυσειδές} beside it. \textbf{The four are not four sins but one sin under four descriptions}, so the eleven petitions under the first are eleven ways of asking one verdict to be lifted.\\[1pt]\textsuperscript{ii}\,⚠⚠ \textbf{Almost every text here states in the indicative what God HAS done, and is asked back as an imperative.} Nathan told David \emph{the LORD hath put away thy sin}, and the page prays \emph{Transfer}; Hezekiah sang \emph{thou hast cast all my sins behind thy back}, and it prays \emph{Dissipa}; the psalm says \emph{thou hast covered all their sin}, and it prays \emph{Tege}. \textbf{The Baptist's is the clearest}: \emph{ὁ αἴρων}, the participle that names Christ, becomes the bare imperative \textbf{αἶρε}. ⚠ \textbf{And it is not a device of one page} — printed 22 asks Isaiah's \emph{I have blotted out thy transgressions} back as \emph{blot thou out mine}, 321 asks of a man's prayer the promise made of God's own word. \textbf{What twice looked like a liberty is the method}, and it is nowhere argued for.}
\clearpage
% ---------- printed 329 ----------
\versoalign
\markright{329}
\setnoted{}
\originalsone{329}{\input{fragments/la329}}{}
\clearpage
\markright{329}
\setnoted{1,4,7,11,16,18,23,24,25}
\enleafn{329}{\input{fragments/en329}}{}{{\scriptsize \textsuperscript{i}\,Ps. xix. 12 — \emph{cleanse thou me from secret faults} \textperiodcentered\ \textsuperscript{ii}\,Isa. xliv. 22 — \emph{I have blotted out thy transgressions as a cloud}, asked here as \emph{Dele}, blot thou out \textperiodcentered\ \textsuperscript{iii}\,Luke x. 35 — the Samaritan to the host: \emph{take care of him} \textperiodcentered\ \textsuperscript{iv}\,Matt. xii. 20 — \emph{a bruised reed shall he not break, and smoking flax shall he not quench} \textperiodcentered\ \textsuperscript{v}\,Ps. cxix. 106 — \emph{I have sworn, and I will perform it} \textperiodcentered\ \textsuperscript{vi}\,2 Cor. vii. 11 — the marks of godly sorrow, here made the penitent's own purposes \textperiodcentered\ \textsuperscript{vii}\,Isa. lv. 7 — \emph{let the wicked forsake his way, and the unrighteous man his thoughts} \textperiodcentered\ \textsuperscript{viii}\,2 Tim. ii. 19 — \emph{let every one that nameth the name of Christ depart from iniquity} \textperiodcentered\ \textsuperscript{ix}\,1 Pet. iv. 3 — \emph{the time past of our life may suffice us}\par}\par\vspace{3pt}\textsuperscript{vi}\,⚠ \textbf{The same six Greek words have now done two different jobs on three consecutive leaves.} At printed 326 Paul's marks of godly sorrow — \emph{ἀγανάκτησις, φόβος, ἐπιπόθησις, ζῆλος, σπουδή, ἀπολογία} — were the \textbf{affections} repentance feels, listed under impersonal verbs. Here they are the \textbf{resolutions} it makes, and the verbs are first person singular and active: \emph{I have sworn, I have longed, I desire, I am zealous, I am earnest, I plead.} \textbf{Paul's list is not being repeated; it is being conjugated.} ⚠ The pass has now recorded five texts used twice and never for the same thing, and this is the largest of them — six words at once.}
\clearpage
% ---------- printed 330 ----------
\versoalign
\markright{330}
\setnoted{}
\originalsone{330}{\input{fragments/la330}}{}
\clearpage
\markright{330}
\setnoted{1,7,10,11,12,14,15,18,24,27,29}
\enleafn{330}{\input{fragments/en330}}{}{{\scriptsize \textsuperscript{i}\,Lam. iii. 40 — \emph{let us search and try our ways, and turn again to the LORD} \textperiodcentered\ \textsuperscript{ii}\,Hos. ii. 6 — \emph{I will hedge up thy way with thorns}, said of a wife to be kept from her lovers \textperiodcentered\ \textsuperscript{iii}\,Rom. vii. 19 — \emph{the good that I would I do not} \textperiodcentered\ \textsuperscript{iv}\,Isa. xxxvii. 3 — Hezekiah to Isaiah, with Sennacherib at the gate: \emph{the children are come to the birth, and there is not strength to bring forth} \textperiodcentered\ \textsuperscript{v}\,Luke xxiv. 38 — the risen Christ to the eleven: \emph{why do thoughts arise in your hearts?} \textperiodcentered\ \textsuperscript{vi}\,Rom. vii. 23 — \emph{bringing me into captivity to the law of sin} \textperiodcentered\ \textsuperscript{vii}\,Isa. xxxviii. 14 — Hezekiah sick: \emph{I am oppressed; undertake for me} \textperiodcentered\ \textsuperscript{viii}\,Ps. lxx. 1 — \emph{make haste, O God, to deliver me} \textperiodcentered\ \textsuperscript{ix}\,Isa. xxxvii. 17 — Hezekiah spreading the letter before the LORD: \emph{open thine eyes, and see} \textperiodcentered\ \textsuperscript{x}\,Job xxxvi. 10 — Elihu: \emph{he openeth also their ear to discipline} \textperiodcentered\ \textsuperscript{xi}\,Neh. vi. 9 — Nehemiah at the wall: \emph{strengthen my hands}\par}\par\vspace{3pt}\textsuperscript{viii}\,\textbf{The ladder of help begins with the words every hour of the Office begins with.} \emph{Deus in adjutorium} is the versicle Andrewes said several times a day, and the first four rungs each set a Greek imperative beside the opening words of a psalm — \emph{βοήθει} with \emph{O God, make speed to help me}, \emph{ἀνάληψαι} with \emph{Save me, O God}, \emph{ἀντιλαβοῦ} with \emph{Let God arise}, \emph{ἐπίστρεψον} with \emph{Arise, O Lord}. \textbf{The Greek names the act and the Latin names the psalm it is asked in}, so a reader who knows the psalter hears four whole psalms behind four words. ⚠ The last four rungs turn the request inward — \emph{open mine ears, lighten mine eyes, confirm my hands}: having asked God to open His eyes and incline His ear, the prayer asks the same two things for itself.}
\clearpage
% ---------- printed 331 ----------
\versoalign
\markright{331}
\setnoted{}
\originalsone{331}{\input{fragments/la331}}{}
\clearpage
\markright{331}
\setnoted{2,3,5,7,8,15,17,21,22,28}
\enleafn{331}{\input{fragments/en331}}{}{{\scriptsize \textsuperscript{i}\,Ps. xxvi. 2 — the Vulgate's \emph{ure renes meos}, burn my reins, where the AV has \emph{try} \textperiodcentered\ \textsuperscript{ii}\,Ps. cxix. 120 — the Vulgate's \emph{nail my flesh with thy fear}, for the AV's \emph{my flesh trembleth} \textperiodcentered\ \textsuperscript{iii}\,Ps. xxxii. 9 — \emph{whose mouth must be held in with bit and bridle} \textperiodcentered\ \textsuperscript{iv}\,2 Pet. ii. 22 — \emph{the dog is turned to his own vomit again} \textperiodcentered\ \textsuperscript{v}\,1 Cor. x. 13 — \emph{there hath no temptation taken you but such as is common to man} \textperiodcentered\ \textsuperscript{vi}\,Gen. iv. 7 — to Cain, before the murder: \emph{sin lieth at the door} \textperiodcentered\ \textsuperscript{vii}\,Gen. xxii. 14 — Abraham naming the place: \emph{in the mount of the LORD it shall be seen} \textperiodcentered\ \textsuperscript{viii}\,Exod. xxxiv. 6, 7 — the Name proclaimed to Moses in the cleft of the rock \textperiodcentered\ \textsuperscript{ix}\,Lev. xxvi. 29 — the covenant curses, and the promise of return hidden in them \textperiodcentered\ \textsuperscript{x}\,Deut. xxx. 1, 2, 3, 6 — the restoration clause: \emph{the LORD thy God will circumcise thine heart}\par}\par\vspace{3pt}\textsuperscript{i}\,⚠ \textbf{Four of these petitions ask to be hurt, and the Latin Bible is why.} \emph{Ure renes}, burn my reins, is the Vulgate where the AV has \emph{try my reins}; \emph{Confige carnes}, nail my flesh with thy fear, is the Vulgate where the AV has \emph{my flesh trembleth}. \textbf{The English Bible has an examination in both places and the Latin has a wound} — and the ladder goes on in the Latin's key, the bit and bridle, the hedge of thorns.\\[1pt]\textsuperscript{vi}\,⚠ \textbf{The soul's rest is quarried out of the covenant curses, and the texts stand in canonical order.} Genesis, Exodus, Leviticus, Deuteronomy, the Psalter — sin at the door, the ten righteous, the mount where the Lord provides, the Name proclaimed in the cleft of the rock, then the two threatening chapters. \textbf{Leviticus xxvi and Deuteronomy xxx are curse-chapters}, and what he takes from them is the clause at the end of each: \emph{until they confess}, \emph{then shall they pray}, \emph{I will remember my covenant}, \emph{he will circumcise thine heart}.}
\clearpage
\sethead{§21 \emph{Quies animæ}}
% ---------- printed 332 ----------
\versoalign
\markright{332}
\setnoted{}
\addcontentsline{toc}{section}{§21 \emph{Quies animæ}}
\originalsone{332}{\input{fragments/la332}}{}
\clearpage
\markright{332}
\setnoted{2,11,19,22,26,28,30}
\enleafn{332}{\input{fragments/en332}}{}{{\scriptsize \textsuperscript{i}\,Ps. xlii. 12 — \emph{hope thou in God: for I shall yet praise him} \textperiodcentered\ \textsuperscript{ii}\,Ps. ciii. 9 — \emph{he will not always chide, neither keepeth he his anger for ever} \textperiodcentered\ \textsuperscript{iv}\,Isa. i. 18, 19 — \emph{though your sins be as scarlet, they shall be as white as snow} \textperiodcentered\ \textsuperscript{v}\,Isa. xlii. 3 — \emph{a bruised reed shall he not break}, the verse Matthew quotes at printed 329 \textperiodcentered\ \textsuperscript{vi}\,Isa. xliii. 25 — \emph{I, even I, am he that blotteth out thy transgressions for mine own sake} \textperiodcentered\ \textsuperscript{vii}\,Isa. xliv. 22 — \emph{I have blotted out thy transgressions as a cloud}, standing here as God speaks it\par}\par\vspace{3pt}\textsuperscript{iii}\,⚠⚠ \textbf{The direction of speech reverses here, and it is the answer to the \emph{Miserere}.} For six leaves the man has been speaking — asking, resolving, confessing. From \emph{Mercy that triumpheth} to the end of the section \textbf{God speaks}, in the first person, and the prayer says nothing at all: \emph{I am he that blotteth out thine iniquities} · \emph{I will bear} · \emph{before they call, I will hear}. The petitions are not answered one by one; they are answered by being \textbf{stopped}, and the man is left listening to promises he has spent six leaves turning into requests.\\[1pt]\textsuperscript{vii}\,⚠⚠ \textbf{The volume proves its own method from inside itself, on one verse.} Isaiah xliv. 22 stands at printed 22 as \emph{Dele … et dissipa} — \textbf{blot thou out MINE}, an imperative in the man's mouth — and it stands here in the indicative and in God's, \emph{I blot out, as a cloud, thine iniquities}, with \emph{return unto me, and I will redeem thee} left attached. \textbf{Same verse, both ways, in one book}, and neither place notes the other. What printed 328 shows as a habit is confirmed here by the plainest possible witness: the edition never has to argue that the imperatives are turned indicatives, because the book turns one back.}
\clearpage
% ---------- printed 333 ----------
\versoalign
\markright{333}
\setnoted{}
\originalsone{333}{\input{fragments/la333}}{}
\clearpage
\markright{333}
\setnoted{1,4,13,16,22,30}
\enleafn{333}{\input{fragments/en333}}{}{{\scriptsize \textsuperscript{i}\,Isa. xlvi. 4 — \emph{even to your old age I am he; and even to hoar hairs will I carry you} \textperiodcentered\ \textsuperscript{ii}\,Isa. liii. 4 — \emph{surely he hath borne our griefs, and carried our sorrows} \textperiodcentered\ \textsuperscript{iii}\,Isa. lxv. 24 — \emph{before they call, I will answer; and while they are yet speaking, I will hear} \textperiodcentered\ \textsuperscript{iv}\,Ezek. xviii. 23 — \emph{have I any pleasure at all that the wicked should die?} \textperiodcentered\ \textsuperscript{v}\,Ezek. xxxiii. 11 — \emph{as I live, saith the Lord GOD, I have no pleasure in the death of the wicked} \textperiodcentered\ \textsuperscript{vi}\,Isa. lv. 7 — \emph{let him return unto the LORD, and he will have mercy upon him}\par}\par\vspace{3pt}\textsuperscript{ii}\,\textbf{The Servant song is the only text in the catena quoted whole.} Everything around it is a single verse — a promise, a question, a clause — but Isaiah liii is given three verses running, wounded, bruised, the chastisement, the stripes, and the sheep gone astray. \textbf{The chapter is not cited for a sentence; it is set down as the ground the rest of the promises stand on}, and it is placed at the centre of the section rather than at its head.\\[1pt]\textsuperscript{v}\,⚠ \textbf{Ezekiel says this twice and the page prints both times, not one.} Chapter xviii asks it as a question and chapter xxxiii swears it as an oath — \emph{as I live} — and the two stand here nine lines apart with their surrounding verses, the turning, the ruin, the wicked who departs and lives. \textbf{The repetition is the prophet's own and the leaf keeps it.} A later hand tidying the duplication would delete an oath.}
\clearpage
\sethead{§22 \emph{Confessio Peccati}}
% ---------- printed 334 ----------
\versoalign
\markright{334}
\setnoted{}
\addcontentsline{toc}{section}{§22 \emph{Confessio Peccati}}
\originalsone{334}{\input{fragments/la334}}{}
\clearpage
\markright{334}
\setnoted{3,6,8,14,15,16,17,19,21}
\enleafn{334}{\input{fragments/en334}}{}{{\scriptsize \textsuperscript{i}\,Ps. lxix. 5 — \emph{my sins are not hid from thee} \textperiodcentered\ \textsuperscript{ii}\,Job xxxi. 33 — Job's oath of clearing: \emph{if I covered my transgressions as Adam} \textperiodcentered\ \textsuperscript{iii}\,Ps. cxli. 4 — \emph{incline not my heart to any evil thing, to practise wicked works} \textperiodcentered\ \textsuperscript{iv}\,Ps. cxix. 176 — \emph{I have gone astray like a lost sheep} \textperiodcentered\ \textsuperscript{v}\,Jer. xxxi. 18 — Ephraim: \emph{as a bullock unaccustomed to the yoke} \textperiodcentered\ \textsuperscript{vi}\,Prov. xxvi. 11 — \emph{as a dog returneth to his vomit, so a fool returneth to his folly} \textperiodcentered\ \textsuperscript{vii}\,2 Pet. ii. 22 — Peter quoting that proverb and adding the washed sow \textperiodcentered\ \textsuperscript{viii}\,Josh. vii. 19, 20 — Joshua to Achan: \emph{give glory to the LORD, and make confession} — and Achan's \emph{thus and thus have I done} \textperiodcentered\ \textsuperscript{ix}\,Matt. xii. 20 — Matthew's \emph{he shall not break}, printed here as \emph{break thou not}\par}\par\vspace{3pt}\textsuperscript{viii}\,\textbf{The confession borrows the words of a man who was stoned for the sin he confessed.} Achan's \emph{I have sinned, and thus and thus have I done} is used here as the form of an honest confession, and it is, precisely: the one confession in Scripture that names the deed in detail. \textbf{The page takes the formula and not the outcome.} ⚠ The valley of Achor, where that story ends, has already been prayed against twice in this volume, at printed 205 and 282; here the man who made it is spoken with rather than deprecated.\\[1pt]\textsuperscript{ix}\,⚠⚠ \textbf{Here the turn from indicative to imperative is visible IN THE GREEK, and the plate does it in the open.} Matthew writes \emph{κάλαμον συντετριμμένον οὐ κατεάξει} — he \textbf{shall} not break the bruised reed, a statement about Christ. The plate prints \emph{κάλαμον συντετριμμένον} \textbf{μὴ κατεάξῃς} — a prohibition addressed to Him. ⚠ \textbf{This verse has now stood on three consecutive leaves and never twice the same way}: at printed 329 as \emph{Extirpa}, at 332 as Isaiah's own promise in God's mouth, and here as a petition in the second person. \textbf{Whatever else the notes record about this edition's method, it can be shown on one verse in three openings}, and the third of them shows it in the original.}
\clearpage
% ---------- printed 335 ----------
\versoalign
\markright{335}
\setnoted{}
\originalsone{335}{\input{fragments/la335}}{}
\clearpage
\markright{335}
\setnoted{6,10,11,12,18}
\enleafn{335}{\input{fragments/en335}}{}{{\scriptsize \textsuperscript{i}\,Rom. ix. 1 — Paul swearing to his grief for his kinsmen: \emph{my conscience also bearing me witness} \textperiodcentered\ \textsuperscript{ii}\,Job vii. 20 — \emph{I am become a burden to myself} \textperiodcentered\ \textsuperscript{iii}\,Ps. li. 17 — \emph{a broken and a contrite heart, O God, thou wilt not despise}; Rom. viii. 26 — \emph{groanings which cannot be uttered} \textperiodcentered\ \textsuperscript{iv}\,Isa. xxiv. 16 — \emph{my leanness, my leanness, woe unto me}; Rom. ii. 5 — \emph{after thy hardness and impenitent heart}; Jer. ix. 1 — \emph{oh that mine head were waters} \textperiodcentered\ \textsuperscript{v}\,2 Cor. vii. 11 — the \emph{indignation} among the marks of godly sorrow, on its fourth leaf running\par}\par\vspace{3pt}\textsuperscript{i}\,⚠ \textbf{Paul's oath was sworn about other men's damnation and is taken here for the speaker's own sin.} \emph{I say the truth in Christ, I lie not, my conscience also bearing me witness} — Paul is about to say he could wish himself accursed \textbf{for his brethren}. The page keeps the oath and its solemnity entire and changes what is sworn to: \emph{that I have so sinned against thee.} \textbf{Fourth text in Part II turned inward that was spoken outward}, and as always nothing on the page marks the change.}
\clearpage
% ---------- printed 336 ----------
\versoalign
\markright{336}
\setnoted{}
\originalsone{336}{\input{fragments/la336}}{}
\clearpage
\markright{336}
\setnoted{1,5,7,13,15,18}
\enleafn{336}{\input{fragments/en336}}{}{{\scriptsize \textsuperscript{i}\,Job xlii. 6 — Job at the end, out of the whirlwind: \emph{I abhor myself, and repent in dust and ashes} \textperiodcentered\ \textsuperscript{ii}\,Ps. xliv. 15 — \emph{my confusion is continually before me, and the shame of my face hath covered me} \textperiodcentered\ \textsuperscript{iii}\,Heb. xii. 21 — Moses at the mountain: \emph{I exceedingly fear and quake} \textperiodcentered\ \textsuperscript{iv}\,Jas. i. 14 — \emph{every man is tempted, when he is drawn away of his own lust, and enticed} \textperiodcentered\ \textsuperscript{v}\,Prov. v. 12 — \emph{how have I hated instruction, and my heart despised reproof!} \textperiodcentered\ \textsuperscript{vi}\,Ps. lv. 5 — \emph{fearfulness and trembling are come upon me, and horror hath overwhelmed me}\par}\par\vspace{3pt}\textsuperscript{iii}\,⚠ \textbf{Five things he confesses not having reverenced, and the fifth is not frightening.} The incomprehensibleness of the glory, the dread power, the awfulness of the presence, the exact justice — and then \emph{thy lovely goodness}, \textbf{bonitas amabilis}. The list is built to make the reader expect a fifth terror and gives him an attraction instead. \textbf{What has gone unreverenced is not only what should have frightened him but what should have drawn him}, and the confession is the graver for it.}
\clearpage
% ---------- printed 337 ----------
\versoalign
\markright{337}
\setnoted{}
\originalsone{337}{\input{fragments/la337}}{}
\clearpage
\markright{337}
\setnoted{6,22,26}
\enleafn{337}{\input{fragments/en337}}{}{{\scriptsize \textsuperscript{iii}\,Mark iv. 38 — the disciples waking Him in the storm: \emph{carest thou not that we perish?}\par}\par\vspace{3pt}\textsuperscript{i}\,\textbf{The judgement is set out in two columns and they are not a list twice over.} The left names what the court \textbf{is} — an incorruptible Judge, a dreadful tribunal, a defence without excuse, arguments not to be escaped — and the right names what the sentence \textbf{is}: an everlasting Gehenna, the fire that is not quenched, the restless worm, chaos immeasurable. \textbf{Procedure on one side and punishment on the other}, and the leaf reads across, so every step of the trial has its outcome standing beside it. ⚠ The two columns are the plate's own arrangement and must survive into both layers.\\[1pt]\textsuperscript{ii}\,⚠⚠ \textbf{The page ends by lodging a legal APPEAL, and it appeals from God to God.} \emph{I have deserved death from thee … yet still I appeal unto thee, O Lord — from the throne of justice to the throne of grace.} An appeal in law moves a cause from a lower court to a higher one, and here both courts are the same God's; nothing is denied and no ground of the sentence is contested. \textbf{The plea is not that the judgement is wrong but that it is not the last word available.} Then \emph{Admit, O Lord, this appeal; if thou admit it not, we perish} — and the last line answers its own verb out of the Gospel.\\[1pt]\textsuperscript{iii}\,⚠ \textbf{The clincher of a courtroom argument is a question asked in a boat.} \emph{Perimus} stands twice in three lines — \emph{if thou admit it not, we perish}, then Mark's \emph{doth it not concern thee, if we perish?} — so the appeal's own word is caught up and handed back by Scripture, and the section ends not on a plea but on the disciples' reproach. ⚠ \textbf{This verse has now been used four times in the volume and never once the same way} — at printed 172 as one of the Gospel cries, at 286 and 294 for two different things again, and here as the last word of a legal appeal. \textbf{Never conform or cross-refer them.}}
\clearpage
% ---------- printed 338 ----------
\versoalign
\markright{338}
\setnoted{}
\originalsone{338}{\input{fragments/la338}}{}
\clearpage
\markright{338}
\setnoted{1,4,8,10,12,13,15,20,22,23}
\enleafn{338}{\input{fragments/en338}}{}{{\scriptsize \textsuperscript{i}\,1 Tim. ii. 4 — \emph{who will have all men to be saved} \textperiodcentered\ \textsuperscript{ii}\,Tit. iii. 11 — \emph{he that is such is subverted, and sinneth, being condemned of himself} \textperiodcentered\ \textsuperscript{iii}\,Gen. xxxii. 10 — Jacob at the ford, before meeting Esau: \emph{I am not worthy of the least of all thy mercies} \textperiodcentered\ \textsuperscript{iv}\,Luke xv. 19 — the prodigal, rehearsing his speech: \emph{make me as one of thy hired servants} \textperiodcentered\ \textsuperscript{v}\,Matt. xv. 27 — the woman of Canaan: \emph{yet the dogs eat of the crumbs} \textperiodcentered\ \textsuperscript{vi}\,Matt. ix. 21 — the woman with the issue of blood: \emph{if I may but touch his garment, I shall be whole} \textperiodcentered\ \textsuperscript{vii}\,1 Pet. v. 6 — \emph{humble yourselves therefore under the mighty hand of God} \textperiodcentered\ \textsuperscript{viii}\,Ps. cxliii. 6 — \emph{my soul thirsteth after thee, as a thirsty land} \textperiodcentered\ \textsuperscript{ix}\,Luke xviii. 13 — the publican, who would not lift up his eyes \textperiodcentered\ \textsuperscript{x}\,Ps. cxxx. 1 — \emph{out of the depths have I cried unto thee, O LORD}\par}\par\vspace{3pt}\textsuperscript{iii}\,⚠⚠ \textbf{Four sayings of unworthiness, each from a different person, and every one of them was answered.} Jacob at the ford, the prodigal on the road, the woman of Canaan, the woman with the issue of blood — and they descend: not worthy of the mercies, not worthy to be a hired servant, not worthy of the crumbs, not worthy to touch the hem. \textbf{Jacob got the blessing, the prodigal the robe, the Canaanite her crumb, the woman her healing.} The ladder of unworthiness is built entirely out of approaches that succeeded, which is the argument it is making without stating it. ⚠ \textbf{And the fourth is not a saying of unworthiness at all.} The woman said \emph{if I may but touch his garment, I shall be whole} — the boldest confidence in the Gospels — and the page turns it into \emph{I am not worthy to touch the hem of thy garment.}}
\clearpage
% ---------- printed 339 ----------
\versoalign
\markright{339}
\setnoted{}
\originalsone{339}{\input{fragments/la339}}{}
\clearpage
\markright{339}
\setnoted{1,3,7,9,13,14,23,26}
\enleafn{339}{\input{fragments/en339}}{}{{\scriptsize \textsuperscript{i}\,Ps. lxxix. 9 — \emph{help us, O God of our salvation, for the glory of thy name} \textperiodcentered\ \textsuperscript{ii}\,Ps. xxv. 10 — \emph{pardon mine iniquity; for it is great} \textperiodcentered\ \textsuperscript{iii}\,Eph. ii. 4 — \emph{God, who is rich in mercy} \textperiodcentered\ \textsuperscript{iv}\,Rom. v. 20 — \emph{where sin abounded, grace did much more abound} \textperiodcentered\ \textsuperscript{v}\,1 Tim. i. 15 — \emph{sinners; of whom I am chief} \textperiodcentered\ \textsuperscript{vi}\,Jas. ii. 13 — \emph{mercy rejoiceth against judgment}, printed here as \emph{superexaltet misericordia} \textperiodcentered\ \textsuperscript{vii}\,Dan. ix. 19 — Daniel in the exile: \emph{O Lord, hearken and do; defer not, for thine own sake} \textperiodcentered\ \textsuperscript{viii}\,1 Tim. iii. 16 — \emph{great is the mystery of godliness: God was manifest in the flesh}\par}\par\vspace{3pt}\textsuperscript{iv}\,⚠ \textbf{A rising ladder of words for mercy, and it is the second of its kind in Part II.} \emph{Multitudo · magna multitudo · divitiæ · abundantia · superabundantia} — five terms, each stronger than the last, each with its own proof-text, ending on Paul's compound. \textbf{This is what printed 319 does in Greek} with περισσεία · πλεονασμός · ὑπερπλεονασμός · ὑπερπερισσεία, a ladder built out of prefixes; here the same figure is built in Latin out of different words. ⚠ \textbf{The two must not be conformed or cross-referred in the text} — but they are one habit, and the summit of both is \emph{superabundare}, which is why the sentence that follows sets the man's abounding transgression against it.}
\clearpage
\sethead{§23 \emph{In magnum pietatis sacramentum}}
% ---------- printed 340 ----------
\versoalign
\markright{340}
\setnoted{}
\addcontentsline{toc}{section}{§23 \emph{In magnum pietatis sacramentum}}
\originalsone{340}{\input{fragments/la340}}{}
\clearpage
\markright{340}
\setnoted{8,11,12,14,16,18,20,22,24}
\enleafn{340}{\input{fragments/en340}}{}{{\scriptsize \textsuperscript{iii}\,Tit. ii. 14 — \emph{that he might purify unto himself a peculiar people} \textperiodcentered\ \textsuperscript{iv}\,Heb. ii. 15 — \emph{deliver them who through fear of death were all their lifetime subject to bondage} \textperiodcentered\ \textsuperscript{v}\,Col. ii. 15 — \emph{he made a shew of them openly, triumphing over them} \textperiodcentered\ \textsuperscript{vi}\,1 Cor. xv. 55 — \emph{O death, where is thy sting? O grave, where is thy victory?} \textperiodcentered\ \textsuperscript{vii}\,Heb. vi. 20 — \emph{whither the forerunner is for us entered, even Jesus} \textperiodcentered\ \textsuperscript{viii}\,Rom. viii. 34 — \emph{who is even at the right hand of God, who also maketh intercession for us} \textperiodcentered\ \textsuperscript{ix}\,Heb. xii. 2 — \emph{looking unto Jesus the author and finisher of our faith}\par}\par\vspace{3pt}\textsuperscript{i}\,\textup{Charismata} and \textup{Fructus} are not two words for one thing: \textbf{the gifts are given for other people, the fruits grow in the man himself} — 1 Cor. xii against Gal. v. 22. A man may have the gift without the fruit, which is the burden of 1 Cor. xiii.\\[1pt]\textsuperscript{ii}\,⚠⚠ \textbf{Four of the six titles are the Vulgate's own word, lifted from the verse printed beside them.} Tit. ii. 14 \emph{ut \textbf{mundaret} sibi populum} gives \textup{mundator}; Heb. ii. 15 \emph{ut \textbf{liberaret} eos} gives \textup{liberator}; Col. ii. 15 \emph{palam \textbf{triumphans} illos} gives \textup{triumphator}; Heb. vi. 20 \emph{ubi \textbf{præcursor} … introivit Iesus} gives \textup{præcursor}, the only place the word occurs in the New Testament. \textbf{The catalogue is quarried, not composed.} ⚠ \textbf{Exactly two depart, and both depart somewhere.} Rom. viii. 34 says only \emph{interpellat pro nobis}; \textup{Advocatus} is at 1 John ii. 1. And Heb. xii. 2 reads \emph{auctorem fidei et \textbf{consummatorem}} where the plate sets \textup{Instaurator} — the verb of Eph. i. 10, \emph{instaurare omnia in Christo}. \textbf{At the second Advent the title is not the finisher of our faith but the restorer of all things}; the English keeps "finisher" because the reference is to Hebrews and Part III renders it so.\\[1pt]\textsuperscript{ix}\,⚠ \textbf{The same six titles, in this order and in these mysteries, stand again at printed 402–403}, in Part III's Harley recension, turned from thanksgiving into belief — \textbf{and there they carry not one reference.} This is the only page where the catena behind them is visible. The two lists are not conformed and must not be.}
\clearpage
% ---------- printed 341 ----------
\versoalign
\markright{341}
\setnoted{}
\originalsone{341}{\input{fragments/la341}}{}
\clearpage
\markright{341}
\setnoted{1,2,7,8,9,25}
\enleafn{341}{\input{fragments/en341}}{}{{\scriptsize \textsuperscript{ii}\,Rev. ix. 11 — the angel of the pit, \emph{whose name in the Hebrew tongue is Abaddon}, destruction. ⚠ \textbf{The row sets one Hebrew name against another}: against \emph{Abaddon}, destruction, stands \emph{JESUS}, salvation — the only line in capitals, and an antithesis only a reader of Hebrew hears. \textperiodcentered\ \textsuperscript{iii}\,Gal. iv. 19 — Paul in travail a second time: \emph{until Christ be formed in you} \textperiodcentered\ \textsuperscript{iv}\,Rom. viii. 29 — \emph{predestinate to be conformed to the image of his Son}\par}\par\vspace{3pt}\textsuperscript{i}\,⚠⚠ \textbf{The braces are not two enemies but one enemy and his translation} — \textup{\{ Satanæ / Adversario \}} is the Hebrew name and its Latin, \textup{\{ Diabolo / Calumniatori \}} the Greek and its Latin. Read as a list the table has nine adversaries; read as set, five names for one. ⚠ \textbf{Three of the five pairs are a courtroom} — Adversary met by Mediator, Slanderer by Advocate, Accuser by Intercessor. Each is answered \textbf{by the answering office, not by a greater power}; the two outer rows, destroyer/Saviour and captor/Redeemer, hold the legal three inside them.\\[1pt]\textsuperscript{v}\,⚠ \textbf{The two remedies are crossed.} Cold in prayer, he is sent to the Throne, the appearing and the intercession; hot in evil desire, to the judgement-seat and the last trump. \textbf{The cure for coldness is not more fear, and the cure for heat is not more comfort.} \textup{Cathedra} is the Throne and \textup{tribunal} the judgement-seat: the same Christ sits in both, and which one to remember depends on what is wrong with the man.\\[1pt]\textsuperscript{vi}\,The four cells of the unction print their own references: \textbf{compunction}, Acts ii. 37, \emph{pricked in their heart}; \textbf{clear knowledge}, 2 Cor. ii. 6 as printed; \textbf{fervent prayer}, Jas. v. 16; \textbf{love shed abroad}, Rom. v. 5. ⚠⚠ \textbf{Three take their noun from the Vulgate of their own verse} — \emph{compunctio} from \emph{compuncti sunt corde}, \emph{diffusio charitatis} from \emph{caritas Dei diffusa est} — \textbf{which is what convicts the fourth.} 2 Cor. \textbf{ii.} 6 concerns a man punished by the many and has nothing to do with knowledge; 2 Cor. \textbf{iv.} 6 reads \emph{ad illuminationem \textbf{scientiæ} claritatis Dei} and supplies \emph{scientia} exactly as the others are supplied. \textbf{A single wrong numeral, ii for iv. Kept as printed; the true verse is 2 Cor. iv. 6} — the quarrying habit of the facing leaf doing forensic work one page later.}
\clearpage
% ---------- printed 342 ----------
\versoalign
\markright{342}
\setnoted{}
\originalsone{342}{\input{fragments/la342}}{}
\clearpage
\markright{342}
\setnoted{1,2,11,12,13,15}
\enleafn{342}{\input{fragments/en342}}{}{{\scriptsize \textsuperscript{i}\,Eph. i. 13, 14 — \emph{sealed with that holy Spirit of promise, which is the earnest of our inheritance}. ⚠ \textbf{Not one figure twice}: a seal marks a thing as owned; an \emph{arrhabo}, the Greek ἀρραβών, is \textbf{a first instalment of the sum itself}, not a token for it. \textperiodcentered\ \textsuperscript{iii}\,Acts ii. 38 — \emph{be baptized … for the remission of sins, and ye shall receive the gift of the Holy Ghost} \textperiodcentered\ \textsuperscript{iv}\,Tit. iii. 7 — \emph{made heirs according to the hope of eternal life}\par}\par\vspace{3pt}\textsuperscript{ii}\,⚠ \textbf{Four things forbidden toward the Holy Ghost, and they ascend} — quench him, 1 Thess. v. 19, as a man puts out a fire; strive against him, Jer. l. 24, said of \textbf{Babylon}; grieve him, Eph. iv. 30, which can only be said of one who loves; and Heb. x. 29, \emph{done despite unto the Spirit of grace}, the verse of falling away past renewal. \textbf{From smothering a flame to insulting a person} — and the middle two are where the Spirit stops being treated as a power. ⚠ \textbf{\textup{Jer. l. 24} is a roman FIFTY, not the figure one}; the reference is correct and must not be "corrected" to Jer. i.\\[1pt]\textsuperscript{v}\,⚠ \textbf{The apostles asked this together; he asks it alone.} Luke xvii. 5 is \emph{adauge \textbf{nobis} fidem}, prayed by the twelve at once; the plate sets \emph{adauge \textbf{mihi}}. \textbf{A prayer made in company is taken into a closet} — as at 327 with Stephen's prayer for other men, and 335 with Paul's oath about his kinsmen, and never marked. The mustard seed on the next line is the Lord's answer, ver. 6.\\[1pt]\textsuperscript{vi}\,⚠⚠ \textbf{Not a description of faith but a concordance of the word.} Nine epithets, each with its reference, and \textbf{every one is the adjective its own verse attaches to \emph{fides}}: dead, Jas. ii. 20; unfeigned, 2 Tim. i. 5, \emph{fide non ficta}, which the plate reverses to \emph{non hypocritica}; working by love, Gal. v. 6, \emph{fides quæ per caritatem operatur}; working with works, Jas. ii. 22, \emph{fides cooperabatur operibus}; ministering virtue, 2 Pet. i. 5; overcoming the world, 1 John v. 4, \emph{hæc est victoria … fides nostra}; most holy, Jude 20, \emph{sanctissimæ fidei vestræ}. \textbf{He has collected every place the word is qualified, set the three that condemn against the five that commend, and made the definition out of the collection.} ⚠ Third leaf running built this way — the six titles at 340, the four cells at 341 — \textbf{the method of the section, not a mannerism.} Only \emph{non temporariam} is not an epithet of faith in its text: Rev. ii. 13 praises Pergamos for holding fast in the days of Antipas, so the fault is named by its opposite.}
\clearpage
% ---------- printed 343 ----------
\versoalign
\markright{343}
\setnoted{}
\originalsone{343}{\input{fragments/la343}}{}
\clearpage
\markright{343}
\setnoted{1,3,5,7,10,18,28}
\enleafn{343}{\input{fragments/en343}}{}{{\scriptsize \textsuperscript{i}\,Deut. xxxii. 4 — \emph{a God of truth and without iniquity} — and Is. ix. 6, \emph{the Prince of Peace}. ⚠ \textbf{The titles are chosen for what the next line asks}: he addresses God as truth and peace, then asks for peace and truth. \textbf{The petition is made out of the address.} \textperiodcentered\ \textsuperscript{iii}\,Acts iv. 32 — \emph{the multitude of them that believed were of one heart and of one soul} \textperiodcentered\ \textsuperscript{iv}\,Is. xlii. 3, and Matt. xii. 20 quoting it — \emph{a bruised reed shall he not break}. ⚠ \textbf{Prophecy and fulfilment cited together}, and what both state as a fact about the Servant becomes a name to call him by. \textperiodcentered\ \textsuperscript{v}\,Dan. ix. 16 — \emph{let thine anger be turned away from thy city} — unmarked in the plate, but Daniel's at the exile, and \textbf{the volume printed ix. 19 of the same prayer at 339}. ⚠ It takes up Daniel's intercession in reverse, the later verse first. \textperiodcentered\ \textsuperscript{vii}\,Ps. xvii. 15 — the close is a psalm the plate does not mark: \emph{I shall behold thy face in righteousness: I shall be satisfied … with thy likeness.} The Latin is the Vulgate's \emph{in justitia \textbf{apparebo} … \textbf{satiabor} cum apparuerit \textbf{gloria} tua}, word for word — \textbf{a quotation the reader will otherwise hear only as a cadence.}\par}\par\vspace{3pt}\textsuperscript{ii}\,⚠⚠ \textbf{Hezekiah said \emph{in my days}; the plate says \emph{in ours}.} Is. xxxix. 8 is verbatim Vulgate — \emph{sit pax et veritas in diebus} — but for the last word, where \emph{meis} has become \textbf{\textup{nostris}}. The verse is the king's answer to the word that his sons would be carried to Babylon, and it is the most self-regarding sentence in his life. \textbf{Made plural, the complacency becomes an intercession.} ⚠ \textbf{The facing leaf does the opposite:} at 342 the apostles' \emph{adauge \textbf{nobis} fidem} is prayed as \emph{mihi}. \textbf{Two consecutive pages alter one pronoun in opposite directions}, each toward what the prayer needs, neither marking that it has.\\[1pt]\textsuperscript{vi}\,⚠ \textbf{He undertakes to pray for people he has never met, on the sole ground that they love him} — \emph{mei amantes, etsi mihi ignotos} — and asks that they be brought into the Kingdom \emph{even as me also}. The petition assumes what it cannot check. \textbf{It is the one intercession here with no scriptural warrant printed beside it}, and it asks for them what he has spent the section asking for himself.}
\clearpage
% ---------- printed 344 ----------
\versoalign
\markright{344}
\setnoted{}
\originalsone{344}{\input{fragments/la344}}{}
\clearpage
\markright{344}
\setnoted{1,8,21,25}
\enleafn{344}{\input{fragments/en344}}{}{\textsuperscript{i}\,\textbf{The scripture foot is empty here and on the next leaf, and that is correct}: every line is a psalm quoted whole with its reference beside it, so a tag would set the verse under the verse. ⚠⚠ \textbf{The figures are a hybrid, checkable twice on this page.} \emph{Ps.} lxxxix. \textbf{16} — the AV divides that psalm so the verse is \textbf{15}; the Vulgate, calling the psalm \textbf{88}, has 16. \emph{Psal.} lvii. \textbf{12} — AV \textbf{11}; Vulgate, calling it 56, has 12. \textbf{The plate keeps the Vulgate's verse-figure under the Hebrew psalm-number}, which is neither system entire — a fourth kind of numbering after the plain AV at 300, the plain Vulgate at 274 and the BCP at 433. \textbf{An introduction claiming one Psalter is answerable for this page.}\\[1pt]\textsuperscript{ii}\,⚠ \textup{et in sæcula, et in sæcula sæculorum} — \textbf{not emphasis, not a doubled line.} Ps. cxlv. 21 is Vulgate \emph{in sæculum et in sæculum sæculi}, for the Greek's \emph{unto the age, and unto the age of the age}. The English doubles because the Latin does; \textbf{a later pass must not tidy one away.}\\[1pt]\textsuperscript{iii}\,⚠ \textbf{Here the plate is against both English Psalters and the English follows the plate.} Ps. cxxxviii. 1 in the Hebrew has \emph{before the \textbf{gods} will I sing praise} — so the AV, and so Coverdale. The Latin is the Vulgate's \emph{in conspectu \textbf{Angelorum}}, the Septuagint's reading, and our English keeps \textbf{angels}. ⚠ Not a departure from CONVENTIONS §9 but the case it exists for: \textbf{the register is Coverdale's, the text is whatever Andrewes quoted.}\\[1pt]\textsuperscript{iv}\,⚠ \textbf{The one line not from the Psalter, and where the leaf turns.} Rev. iv. 11 is the song of the four-and-twenty elders, arriving straight after he calls himself an unworthy sinner. \textbf{Having exhausted the earth's book of praise he borrows heaven's} — and the section ends in the twelve-membered doxology of Rev. v, so this is the hinge, not an interruption.}
\clearpage
% ---------- printed 345 ----------
\versoalign
\markright{345}
\setnoted{}
\originalsone{345}{\input{fragments/la345}}{}
\clearpage
\markright{345}
\setnoted{2,3,8,10,12,22}
\enleafn{345}{\input{fragments/en345}}{}{{\scriptsize \textsuperscript{i}\,Luke ii. 14 — the angels over Bethlehem: \emph{on earth peace, good will toward men} \textperiodcentered\ \textsuperscript{iii}\,Rev. v. 12 — \emph{Worthy is the Lamb that was slain to receive power, and riches, and wisdom, and strength, and honour, and glory, and blessing}; Rev. v. 13 — \emph{be unto him that sitteth upon the throne, and unto the Lamb} \textperiodcentered\ \textsuperscript{iv}\,Matt. xxi. 9 — the crowds at the entry into Jerusalem: \emph{Hosanna in the highest}\par}\par\vspace{3pt}\textsuperscript{ii}\,⚠⚠ \textbf{The plate cites two verses and sets twelve terms; neither verse has twelve.} Rev. v. 12 gives seven — power, riches, wisdom, strength, honour, glory, blessing — and v. 13 four, all already among the seven. \textbf{Five are therefore from elsewhere in the book:} \emph{might} and \emph{thanksgiving} stand together at Rev. vii. 12, and \emph{salvation} in the cry at Rev. vii. 10, \textbf{whose own words — \emph{unto our God which sitteth upon the throne, and unto the Lamb} — are the two lines this grid runs into.} ⚠ The Apocalypse's own lists run to \textbf{seven}, twice; this one runs to \textbf{twelve}, and the page does not say why.\\[1pt]\textsuperscript{v}\,⚠ \textbf{The Confession of Praise opens by arguing that he has no right to make it} — sinning and unrepentant, it would be \emph{more fitting} to lie prostrate and beg pardon than praise God with a polluted mouth. He praises anyway, on one ground: \emph{de ingenita Tua bonitate confidens}, \textbf{trusting a goodness that is God's by nature and not a response to anything in him.} The whole catalogue hangs from that clause.\\[1pt]\textsuperscript{vi}\,⚠⚠ \textbf{The catalogue is the Roman law of persons.} \emph{Liber} against \emph{Verna} is not free against slave but free against \textbf{a slave born in the house}; \emph{Legitimus} against \emph{Spurius} is the law of lawful issue; \emph{Homo civilis} against \emph{Barbarus}, citizen against outsider. \textbf{He goes through the divisions of the \emph{status personarum} one by one, giving thanks for a legal condition as another man might for health.} ⚠⚠ \textbf{The same catalogue stands again at printed 428}, in Part III §4's \emph{Quod} brace — \textbf{with the antitheses stripped out} (\emph{Vivo · Ratione præditus · Civilis · Liber · Ingenuus · Honesta stirpe}). Here every good is defined by what he might have been instead; there they are simply goods. \emph{Ingenuus} is the very term answering this page's \emph{non Verna}. \textbf{Not to be conformed: the antitheses are the whole difference.}}
\clearpage
\sethead{§24 \emph{Confessio Laudis}}
% ---------- printed 346 ----------
\versoalign
\markright{346}
\setnoted{}
\addcontentsline{toc}{section}{§24 \emph{Confessio Laudis}}
\originalsone{346}{\input{fragments/la346}}{}
\clearpage
\markright{346}
\setnoted{1,6,12,16,23,26}
\enleafn{346}{\input{fragments/en346}}{}{{\scriptsize \textsuperscript{v}\,1 Pet. i. 3 — \emph{begotten us again unto a lively hope by the resurrection of Jesus Christ}; and in the editor's own bracket \textbf{[4]}, \emph{an inheritance incorruptible, and undefiled, and that fadeth not away, reserved in heaven for you} — \textbf{the page quotes both verses nearly entire.} \textperiodcentered\ \textsuperscript{vi}\,Eph. i. 3 — \emph{blessed us with all spiritual blessings in heavenly places in Christ}\par}\par\vspace{3pt}\textsuperscript{i}\,⚠ \textbf{A proverb, not a description.} \emph{Mali corvi malum ovum} — the Greek κακοῦ κόρακος κακὸν ᾠόν, \emph{a bad crow's bad egg}, said of a man whose badness is inherited and expected. \textbf{The line thanks God not for respectable parents but for not being the sort of son people say that about.}\\[1pt]\textsuperscript{ii}\,⚠⚠ \textbf{\textup{Non expositus} is the technical word for a child EXPOSED} — carried out and left to die, lawful and ordinary in the world whose vocabulary this catalogue uses. It stands between \emph{educatus}, reared, and \emph{literatus}, taught. ⚠ \textup{Mechanicus} here is \textbf{a manual labourer}, not a machinist; \textup{bardus} two lines above is dull-witted. \textbf{The list thanks God for a place in a social order, in that order's own terms.}\\[1pt]\textsuperscript{iii}\,⚠ \textbf{The most concrete line in the catalogue names what poverty actually costs}: an honest estate, \emph{so that I have need neither to flatter, nor to borrow.} \textbf{Not wealth — enough not to have to fawn and not to have to owe.} The two are set as one item.\\[1pt]\textsuperscript{iv}\,\textbf{The gifts are divided as the schools divided them} — \emph{bona gratiae, naturae, fortunae} — and only the first is subdivided, into \textbf{redemption and calling}. ⚠ The same three stand again at printed 429 as \emph{Pro bonis \{ naturæ … \}}, the fourth recension parallel. \textbf{Not to be conformed.}}
\clearpage
% ---------- printed 347 ----------
\versoalign
\markright{347}
\setnoted{}
\originalsone{347}{\input{fragments/la347}}{}
\clearpage
\markright{347}
\setnoted{3,8,9,14,22}
\enleafn{347}{\input{fragments/en347}}{}{{\scriptsize \textsuperscript{i}\,2 Cor. i. 4 — \emph{who comforteth us in all our tribulation}; 2 Cor. i. 5 — \emph{as the sufferings of Christ abound in us, so our consolation also aboundeth}. \textbf{Written out of the affliction in Asia in which he despaired of life, and quoted nearly whole above.} \textperiodcentered\ \textsuperscript{ii}\,Dan. ii. 23 — Daniel's thanksgiving the night the king's dream was shewn him, with the executioner already appointed for the wise men of Babylon. \textbf{His words stand in the four lines above}; the bracket is the editor's.\par}\par\vspace{3pt}\textsuperscript{iii}\,⚠ \textbf{Three pairs, and the columns are creation against redemption} — what God made and marked (hands, countenance, name) against what Christ bought (blood, purchase, adoption), \textbf{each left-hand title answered by the cost of the right.} The two unpaired lines run the full measure and close the table Father, Son, Spirit, Church.\\[1pt]\textsuperscript{iv}\,⚠⚠⚠ \textbf{\textup{Divisio} is the book explaining its own shape, and it is the most important page in Part II for anyone writing about this volume.} What follows is not a prayer but \textbf{an analysis of what a prayer contains}; the leaf prints the parts, so they are not repeated here. ⚠ \textbf{Most of these words are actual section headings in this volume} — \emph{Confessio peccatorum}, \emph{Confessio fidei}, \emph{Confessio laudis}, \emph{Deprecatio}, \emph{Intercessio}, \emph{Supplicatio}, \emph{Gratiarum actio}, \emph{Commendatio} — \textbf{so the page is a table of contents for a book that never prints one.} \emph{Compellatio}, \emph{Comprecatio} and \emph{Threni} appear here only.\\[1pt]\textsuperscript{v}\,⚠⚠ \textbf{THE COLUMNS HERE ARE NOT A PAIRING} — unlike the table four lines above, which \emph{is} read across. \textbf{The \textup{Divisio} is one list broken into two columns, read down the left then down the right.} The proof is the join: \textup{Deprecatio} stands at the foot of the left column and its members, \emph{of sin} and \emph{of punishment}, at the head of the right. \textbf{Read across, it marries the address to sin.}}
\clearpage
% ---------- printed 348 ----------
\versoalign
\markright{348}
\setnoted{}
\originalsone{348}{\input{fragments/la348}}{}
\clearpage
\markright{348}
\setnoted{1,5,14,26}
\enleafn{348}{\input{fragments/en348}}{}{{\scriptsize \textsuperscript{iv}\,Gen. xxxii. 10 — Jacob at the ford, before meeting Esau: \emph{I am not worthy of the least of all thy mercies, and of all the truth which thou hast shewed unto thy servant.} ⚠ \textbf{The volume has used this verse once already, at printed 338}, where it stood first in a ladder of four sayings of unworthiness that were every one of them answered. \textbf{Here it opens a confession instead. Never conform or cross-refer the two places.}\par}\par\vspace{3pt}\textsuperscript{i}\,⚠⚠ \textbf{The \textup{Divisio} continues, and here the columns ARE parallel} — unlike 347's, which is one list wrapped. \textbf{Left is the FORM of the prayer}, right is \textbf{what one prays THROUGH}, headed by the bare preposition \textup{Per}. ⚠ \textbf{And the right column is the whole economy, in order and Trinitarian}: Creation, Nurture, Governance, Preservation, Disposing; then Redemption opened into Conception, Nativity, Life, Passion and Death, Resurrection, Ascension; then, alone at the foot, \textbf{Inspiration} — the Son's part given eleven words to the Father's five and the Spirit's one.\\[1pt]\textsuperscript{ii}\,⚠⚠ \textbf{The three little words between the columns are three different ways of joining a petition to God's acts.} \emph{Quia} — \textbf{because} thou didst these things. \emph{Propter vel Per} — \textbf{for the sake of}, or \textbf{through}, them. \emph{Secundum} — \textbf{according to} them, as their measure. \textbf{The same confession may be grounded on God's work as its cause, as its channel, or as its standard}, and the page sets the three side by side without choosing. \emph{Et annunciabo} closes the list: the prayer ends by undertaking to tell.\\[1pt]\textsuperscript{iii}\,⚠ \textbf{\textup{Ἄσκησις} is the one Greek word on the leaf}, and it governs what follows: the \emph{exercise}, daily or nightly — penitence for evils, thankfulness for good things — and then a \textbf{particular} exercise, of age, of calling, of estate. ⚠⚠ \textbf{The brace answers all three heads at once, and its worked example is a single man: \emph{si senex; in Clero; Episcopus} — an old man, in the clergy, a bishop.} Andrewes was Bishop of Winchester and died in his seventy-first year. \textbf{The one instance the scheme troubles to give is his own condition.}}
\clearpage
\sethead{§25 \emph{Confessio Peccati et Protestatio Gratitudinis}}
% ---------- printed 349 ----------
\versoalign
\markright{349}
\setnoted{}
\addcontentsline{toc}{section}{§25 \emph{Confessio Peccati et Protestatio Gratitudinis}}
\originalsone{349}{\input{fragments/la349}}{}
\clearpage
\markright{349}
\setnoted{1,4,7,8,10,12,16}
\enleafn{349}{\input{fragments/en349}}{}{{\scriptsize \textsuperscript{ii}\,2 Sam. vii. 20 — David after Nathan brought him the promise of a house: \emph{what can David say more unto thee?} \textperiodcentered\ \textsuperscript{iii}\,2 Sam. ix. 8 — \textbf{Mephibosheth}, lame in both feet, Saul's grandson, called to eat at David's table continually: \emph{what is thy servant, that thou shouldest look upon such a dead dog as I am?} ⚠ \textbf{The bitterest self-naming in the Old Testament, and it is spoken by a man being shewn kindness.} \textperiodcentered\ \textsuperscript{iv}\,John xiii. 1 — the Supper: \emph{having loved his own which were in the world, he loved them unto the end} \textperiodcentered\ \textsuperscript{v}\,Ps. cxvi. 12 — \emph{what shall I render unto the LORD for all his benefits toward me?} \textperiodcentered\ \textsuperscript{vi}\,1 Thess. iii. 9 — Paul of the Thessalonians: \emph{what thanks can we render to God again … for all the joy wherewith we joy for your sakes before our God?}\par}\par\vspace{3pt}\textsuperscript{i}\,⚠⚠ \textbf{Five questions in a row, and not one of them is answered.} A section headed \emph{Protestation of Thankfulness} opens by finding nothing to say: David after the promise, Mephibosheth at the king's table, the psalmist, Paul. \textbf{Each asks what can be said or given back, and the page simply sets the next question under the last.} ⚠ \textbf{The third is not a question in its own text.} John xiii. 1 is a statement — \emph{having loved his own, he loved them unto the end} — and the plate turns it into \emph{For that thou hast loved me hitherto?} \textbf{The volume's habit of changing a verse's mood, elsewhere indicative into imperative, here makes an indicative interrogative.}\\[1pt]\textsuperscript{vii}\,⚠⚠ \textbf{From here the private prayer takes the shape of the EUCHARISTIC ANAPHORA.} A \emph{spiritual sacrifice} is offered, received \emph{upon thy spiritual Altar}, and the Spirit is asked to be \textbf{sent down in return} — which is an \textbf{epiclesis}, the invocation that follows the offering in the Greek liturgies. ⚠ And \emph{voluntary and involuntary} is their formula too, ἑκουσίων καὶ ἀκουσίων, which no English rite has. \textbf{A man alone in his chamber is praying the shape of the Church's central action over himself.}}
\clearpage
% ---------- printed 350 ----------
\versoalign
\markright{350}
\setnoted{}
\originalsone{350}{\input{fragments/la350}}{}
\clearpage
\markright{350}
\setnoted{6,8,11,23}
\enleafn{350}{\input{fragments/en350}}{}{{\scriptsize \textsuperscript{i}\,Eph. iii. 20 — \emph{unto him that is able to do exceeding abundantly above all that we ask or think, according to the power that worketh in us}; Eph. iii. 21 — \emph{unto him be glory in the church by Christ Jesus throughout all ages, world without end} \textperiodcentered\ \textsuperscript{ii}\,Ps. lxiii. 5 — \emph{my soul shall be satisfied as with marrow and fatness; and my mouth shall praise thee with joyful lips}\par}\par\vspace{3pt}\textsuperscript{iii}\,⚠⚠ \textbf{The Lord's Prayer is not simply said here — it is TAKEN APART, in exactly the way the \textup{Divisio} two leaves back took prayer apart.} The seven petitions are \textbf{numbered}, their verbs pulled out into a left column and their objects set right. ⚠ \textbf{And the brace is doing grammar}: \textup{\{ Nomen / Regnum \} Tuum} sets one \emph{Tuum} to serve two petitions, which English cannot copy and must say twice. \textbf{The prayer is being read as a structure, by a man who has just finished drawing the structure of prayer in general.} ⚠ The volume returns to it twice more — §26 \emph{Oratio Dominica variata} and §31's petition-by-petition meditation — \textbf{and the three are not to be conformed; this is the plainest of them.}\\[1pt]\textsuperscript{iv}\,⚠ \textbf{The doxology is here, and the Latin Bible does not have it.} \emph{Quia Tuum est regnum, potentia et gloria} is absent from the Vulgate at Matt. vi. 13 and from the Roman rite, which keeps it apart from the prayer; it stands in the Greek text and in the Greek liturgies. \textbf{On a Latin page, its presence is a mark of which prayer-book was actually to hand.}}
\clearpage
\sethead{§26 \emph{Oratio Dominica Variata}}
% ---------- printed 351 ----------
\versoalign
\markright{351}
\setnoted{}
\addcontentsline{toc}{section}{§26 \emph{Oratio Dominica Variata}}
\originalsone{351}{\input{fragments/la351}}{}
\clearpage
\markright{351}
\setnoted{3,5,8,10,12,14,15,18,20,21,23,25,27}
\enleafn{351}{\input{fragments/en351}}{}{{\scriptsize \textsuperscript{ii}\,Gen. xv. 1 — the word to \textbf{Abram} in the vision, before he had a son: \emph{I am thy shield, and thy exceeding great reward} \textperiodcentered\ \textsuperscript{iii}\,Num. xxiv. 13 — ⚠⚠ \textbf{Balaam}, the prophet hired to curse Israel: \emph{I cannot go beyond the commandment of the LORD, to do either good or bad of mine own mind.} \textbf{The petition \emph{thy will be done} is put into the mouth of a man who wanted to do otherwise and could not.} \textperiodcentered\ \textsuperscript{iv}\,Gen. xxviii. 20 — \textbf{Jacob's vow at Bethel}, with the stone still under his head: \emph{if God will give me bread to eat, and raiment to put on} \textperiodcentered\ \textsuperscript{v}\,1 Kings viii. 36 — \textbf{Solomon} at the dedication of the temple: \emph{forgive the sin of thy servants} \textperiodcentered\ \textsuperscript{vi}\,Matt. vi. 34 — \emph{take therefore no thought for the morrow} \textperiodcentered\ \textsuperscript{vii}\,Ps. xl. 12 — \emph{innumerable evils have compassed me about} \textperiodcentered\ \textsuperscript{viii}\,Ps. cxiii. 2 — \emph{blessed be the name of the LORD from this time forth and for evermore} \textperiodcentered\ \textsuperscript{ix}\,Job xxxiv. 30 — Elihu: \emph{that the hypocrite reign not, lest the people be ensnared} — ⚠ \textbf{set II's petition for the KINGDOM is a prayer against bad government.} \textperiodcentered\ \textsuperscript{x}\,1 Sam. iii. 18 — \textbf{Eli}, when the child Samuel had told him the whole doom of his house: \emph{it is the LORD: let him do what seemeth him good.} \textbf{Thy will be done, said by a man who has just been told he is ruined.} \textperiodcentered\ \textsuperscript{xi}\,Job xxxi. 40 — ⚠⚠ \textbf{Job's self-curse turned into the petition for daily bread by one word.} Closing his oath of innocence Job says \emph{let thistles grow instead of wheat, and cockle instead of barley}; the plate sets \emph{pro tritico \textbf{ne} proveniat urtica, pro hordeo \textbf{ne} lolium} — \textbf{let them NOT.} The curse a man called down on himself if he lied becomes the fourth petition. \textperiodcentered\ \textsuperscript{xii}\,Job vii. 20 — \emph{I have sinned; what shall I do unto thee, O thou preserver of men?} — Job's words made the petition for pardon \textperiodcentered\ \textsuperscript{xiii}\,Job xxxi. 1 — \emph{I made a covenant with mine eyes} — ⚠ \textbf{but the plate widens it twice}: the covenant is made \emph{cum sensibus meis}, with all the senses, and its object is \emph{ut ne cogitem quidem de malo}, that he should not so much as think of evil, where Job's own concern was narrower.\par}\par\vspace{3pt}\textsuperscript{i}\,⚠⚠⚠ \textbf{The prayer said plainly on the leaf before is now said again SEVEN TIMES OVER, petition by petition, in other men's words} — thy kingdom come by the word to Abram, thy will be done by \textbf{Balaam}, daily bread by \textbf{Jacob's vow at Bethel}, forgiveness by \textbf{Solomon at the dedication}. \textbf{Then set II begins the seven again with different speakers.} ⚠ \textbf{This is the method of §§33–34}, the creed said in the words of ten confessors — \textbf{applied here to the Lord's Prayer, and applied earlier. The petitions are never restated in Andrewes' own words at all.}}
\clearpage
% ---------- printed 352 ----------
\versoalign
\markright{352}
\setnoted{}
\originalsone{352}{\input{fragments/la352}}{}
\clearpage
\markright{352}
\setnoted{2,4,5,6,9,11,12,15,19,22,26}
\enleafn{352}{\input{fragments/en352}}{}{{\scriptsize \textsuperscript{i}\,Job v. 19 — Eliphaz: \emph{he shall deliver thee in six troubles: yea, in seven there shall no evil touch thee} \textperiodcentered\ \textsuperscript{ii}\,Exod. xxviii. 36 — the golden plate of the mitre, \emph{graven, like the engravings of a signet, HOLINESS TO THE LORD} \textperiodcentered\ \textsuperscript{iii}\,1 Pet. ii. 9 — \emph{ye are a chosen generation, a royal priesthood, an holy nation} \textperiodcentered\ \textsuperscript{iv}\,John x. 9 — \emph{I am the door: by me if any man enter in, he shall be saved, and shall go in and out} \textperiodcentered\ \textsuperscript{v}\,Matt. iv. 4 — ⚠⚠ \textbf{the petition for daily bread said in the words that refuse bread.} Christ to the tempter: \emph{man shall not live by bread alone, but by every word that proceedeth out of the mouth of God.} \textbf{This is the second time in two leaves that the fourth petition is made out of a text against bread} — set II built it from Job's self-curse of thistles for wheat. \textbf{The petition is never simply asked.} \textperiodcentered\ \textsuperscript{vi}\,Ps. xxv. 18 — \emph{look upon mine affliction and my pain; and forgive all my sins} \textperiodcentered\ \textsuperscript{vii}\,Ps. xcv. 8 — \emph{harden not your heart, as in the provocation}; Heb. xii. 15 — \emph{lest any root of bitterness springing up trouble you}. ⚠ \textbf{Two texts for the sixth petition, and both name a place in the wilderness where a people went wrong.} \textperiodcentered\ \textsuperscript{viii}\,2 Sam. xxiv. 16 — the angel stretching out his hand upon Jerusalem, and the LORD saying \emph{It is enough: stay now thine hand} \textperiodcentered\ \textsuperscript{ix}\,Ps. cxiii. 2 — \emph{blessed be the name of the LORD from this time forth and for evermore}. ⚠ \textbf{It opened set II on the leaf before and opens set IV here: the one verse the sets reuse is the one that blesses the Name}, which is the petition they all begin from. \textperiodcentered\ \textsuperscript{x}\,Ps. cxlii. 5 — \emph{thou art my refuge and my portion in the land of the living} \textperiodcentered\ \textsuperscript{xi}\,Ps. cxliii. 10 — \emph{thy spirit is good; lead me into the land of righteousness}\par}\par\vspace{3pt}\textsuperscript{ii}\,⚠⚠ \textbf{Set III opens the seven again, and its first petition is a thing engraved, not said.} \emph{Sanctitas Domino} — \textbf{HOLINESS TO THE LORD} — is the inscription cut on the golden plate of the high priest's mitre, worn on the forehead. \textbf{For \emph{hallowed be thy Name} the page does not quote a prayer but a piece of metal}, and the two that follow keep the priestly setting: a \textbf{royal priesthood}, and going in and out \textbf{by the door}.}
\clearpage
% ---------- printed 353 ----------
\versoalign
\markright{353}
\setnoted{}
\originalsone{353}{\input{fragments/la353}}{}
\clearpage
\markright{353}
\setnoted{2,5,7,11,12,15,17,21,25}
\enleafn{353}{\input{fragments/en353}}{}{{\scriptsize \textsuperscript{i}\,Ps. cxlv. 15 — \emph{thou givest them their meat in due season}; Ps. cxlv. 16 — \emph{thou openest thine hand, and satisfiest the desire of every living thing} \textperiodcentered\ \textsuperscript{ii}\,Ps. li. 1 — \emph{have mercy upon me, O God, according to thy lovingkindness} \textperiodcentered\ \textsuperscript{iii}\,Ps. lxxxix. 23 — \emph{the enemy shall not exact upon him; nor the son of wickedness afflict him} \textperiodcentered\ \textsuperscript{v}\,Prov. xviii. 10 — \emph{the name of the LORD is a strong tower: the righteous runneth into it, and is safe} \textperiodcentered\ \textsuperscript{vi}\,Prov. viii. 15 — \emph{by me kings reign}; Prov. xxi. 1 — \emph{the king's heart is in the hand of the LORD, as the rivers of water: he turneth it whithersoever he will}. ⚠ \textbf{\textup{Flecte ad bonum} is in neither verse — the page adds the petition to the two texts.} \textperiodcentered\ \textsuperscript{viii}\,Prov. xxx. 7 — \textbf{Agur}: \emph{two things have I required of thee; deny me them not before I die}; Prov. xxx. 8 — \emph{give me neither poverty nor riches; feed me with food convenient for me} \textperiodcentered\ \textsuperscript{ix}\,Prov. xx. 9 — \emph{who can say, I have made my heart clean?}; Ps. xli. 5 — the plate's figure, which is ⚠ \textbf{the Vulgate's verse-number under the Hebrew psalm-number}: the AV and the Prayer Book have it at xli. 4, \emph{heal my soul; for I have sinned against thee}. \textbf{Second page to show that hybrid}, after printed 344.\par}\par\vspace{3pt}\textsuperscript{iv}\,⚠⚠ \textbf{Set V says the whole prayer out of the WISDOM books, almost all of it Proverbs} — where sets I and II used men speaking in a story, this one uses a book of sentences. ⚠ \textbf{And its fourth petition is at last the natural one.} After Job's self-curse at 351 and \emph{man shall not live by bread alone} at 352, set V reaches \textbf{Agur's} \emph{feed me with food convenient for me} — \textbf{the one place in Scripture where a man asks for exactly daily bread. Three leaves, three treatments, and the plain sense comes last.}\\[1pt]\textsuperscript{vii}\,⚠ \textbf{Prov. xix. 21 turned into the third petition, borrowing the Pater Noster's own verb.} The proverb states a fact — \emph{there are many devices in a man's heart; nevertheless the counsel of the LORD, that shall stand} — and the page asks it: \emph{let there not be many devices in our hearts; but let thy counsel stand, \textbf{et fiat}}. \textbf{\textup{Fiat} is the word of \emph{thy will be done}, imported into the proverb standing in for it.}}
\clearpage
% ---------- printed 354 ----------
\versoalign
\markright{354}
\setnoted{}
\originalsone{354}{\input{fragments/la354}}{}
\clearpage
\markright{354}
\setnoted{2,3,5,7,9,12,14,17,18,27,29}
\enleafn{354}{\input{fragments/en354}}{}{{\scriptsize \textsuperscript{ii}\,Rom. ii. 24 — \emph{the name of God is blasphemed among the Gentiles through you} \textperiodcentered\ \textsuperscript{iii}\,Isa. lx. 12 — \emph{the nation and kingdom that will not serve thee shall perish; yea, those nations shall be utterly wasted} \textperiodcentered\ \textsuperscript{iv}\,Isa. xlvi. 10 — \emph{my counsel shall stand, and I will do all my pleasure} \textperiodcentered\ \textsuperscript{v}\,Isa. lv. 10 — \emph{that it may give seed to the sower, and bread to the eater}. ⚠ \textbf{\textup{Baculum panis}, the STAFF of bread, is not in that verse} — it is the prophets' phrase for the staff that is broken in famine, brought in here to ask that it be given. \textperiodcentered\ \textsuperscript{vi}\,Isa. lxiv. 9 — \emph{be not wroth very sore, O LORD, neither remember iniquity for ever: behold, see, we beseech thee, we are all thy people.} ⚠ \textbf{The volume has used this verse once already, at printed 109}, in the confession — \textbf{here it is the fifth petition. Never conform the two places.} \textperiodcentered\ \textsuperscript{vii}\,Ezek. xiv. 4 — \emph{that putteth the stumblingblock of his iniquity before his face} \textperiodcentered\ \textsuperscript{viii}\,Luke v. 5 — Peter at the draught of fishes: \emph{Master, we have toiled all the night, and have taken nothing} \textperiodcentered\ \textsuperscript{x}\,2 Cor. iii. 14 — \emph{the vail untaken away in the reading of the old testament} \textperiodcentered\ \textsuperscript{xi}\,Ps. cxix. 12 — \emph{blessed art thou, O LORD: teach me thy statutes}\par}\par\vspace{3pt}\textsuperscript{i}\,⚠⚠ \textbf{The later sets are each quarried from ONE region of Scripture, and set VI completes the sweep} — \textbf{IV} the Psalms, \textbf{V} Proverbs, \textbf{VI} the Prophets, where the first three ranged. \textbf{The Lord's Prayer is said through the whole Bible, a region at a time.} ⚠⚠ \textbf{And set VI is the harshest}: \emph{thy kingdom come} is asked as \emph{let all the nations and kings perish that serve not thy kingdom}. \textbf{Nothing earlier is spoken in that voice, and the page does not soften it.}\\[1pt]\textsuperscript{ix}\,⚠⚠ \textbf{A new section opens with a key, the plainest statement in the volume of how Andrewes reads a miracle.} Five pairs read across: the world is the sea, men the fishes, the Church the boat, the preacher the fisherman, the Word the net. \textbf{The draught of fishes set out as a working diagram for a man about to preach}, hung on Peter's \emph{we have toiled all the night and taken nothing}. ⚠ \textbf{This table IS read across} — unlike the \textup{Divisio} at 347.}
\clearpage
\sethead{§27 \emph{In Prædicatione}}
% ---------- printed 355 ----------
\versoalign
\markright{355}
\setnoted{}
\addcontentsline{toc}{section}{§27 \emph{In Prædicatione · Ad Verbum · Ad Templi ingressum}}
\originalsone{355}{\input{fragments/la355}}{}
\clearpage
\markright{355}
\setnoted{2,3,4,8,17,18,19,28}
\enleafn{355}{\input{fragments/en355}}{}{{\scriptsize \textsuperscript{i}\,1 Sam. x. 26 — of the band that went with Saul: \emph{whose hearts God had touched} \textperiodcentered\ \textsuperscript{ii}\,Eph. i. 8 — the plate's figure. ⚠ \textbf{The line wants Eph. i. 18}, \emph{the eyes of your understanding being enlightened}, which is what \emph{illumina sensus cordis} renders; i. 8 is \emph{he hath abounded toward us in all wisdom and prudence}. \textbf{Kept as printed} — it stands on the volume's running list of wrong references with \textup{[Jer. l. 24.]} and \textup{[Dan. vi. 4.]} \textperiodcentered\ \textsuperscript{iii}\,Ps. li. 15 — \emph{O Lord, open thou my lips}; Ps. lxxi. 8 — \emph{let my mouth be filled with thy praise}. ⚠ \textbf{Ps. li. 15 returns alone nine lines below} — the one verse this page says twice, and the page is about nothing but the mouth. \textperiodcentered\ \textsuperscript{v}\,Isa. l. 4 — \emph{the Lord GOD hath given me the tongue of the learned, that I should know how to speak a word in season}; Eph. iv. 29 — \emph{that it may minister grace unto the hearers} \textperiodcentered\ \textsuperscript{vi}\,Eph. vi. 19 — Paul asking prayer for himself: \emph{that utterance may be given unto me, that I may open my mouth boldly} \textperiodcentered\ \textsuperscript{vii}\,Ps. lxxxi. 10 — \emph{open thy mouth wide, and I will fill it} \textperiodcentered\ \textsuperscript{viii}\,Ps. xxviii. 2 — \emph{when I lift up my hands toward thy holy oracle}\par}\par\vspace{3pt}\textsuperscript{iv}\,⚠⚠ \textbf{\textup{Carbo duplicis naturæ} — the COAL OF THE TWOFOLD NATURE — is a title for Christ, and it is not in Isaiah.} The prophet's coal, taken from the altar with tongs and laid on his mouth, is read as \textbf{the coal that is wood and fire at once, and so God and man at once}; the tongs are what handles what cannot be touched. ⚠ \textbf{The reading is Greek} — the ἄνθραξ of the Eastern liturgies, where the same coal is what the communicant receives — \textbf{the fourth mark of that rite in six leaves}, after 349's anaphora and \emph{voluntary and involuntary} and 350's doxology, \textbf{and the strongest, because it carries a whole Christology in one noun.}\\[1pt]\textsuperscript{vii}\,⚠⚠ \textbf{Ps. lxxxi. 10 is God's own command, answered as though it had been obeyed.} \emph{Open thy mouth wide, and I will fill it} — the psalm; \emph{Dilato os meum, Domine: \textbf{Tu imple}} — \textbf{I have opened my mouth wide: do thou fill it.} The condition is reported as performed and the promise called in. ⚠ \textbf{Not an indicative made imperative, but a promise TAKEN UP and its half discharged.}}
\clearpage
% ---------- printed 356 ----------
\versoalign
\markright{356}
\setnoted{}
\originalsone{356}{\input{fragments/la356}}{}
\clearpage
\markright{356}
\setnoted{1,10,12,17,20,21,23,26,27,29}
\enleafn{356}{\input{fragments/en356}}{}{{\scriptsize \textsuperscript{i}\,Ps. xlviii. 9 — \emph{we have thought of thy lovingkindness, O God, in the midst of thy temple} \textperiodcentered\ \textsuperscript{ii}\,Ps. xxvi. 8 — \emph{LORD, I have loved the habitation of thy house, and the place where thine honour dwelleth} \textperiodcentered\ \textsuperscript{iii}\,Ps. xxvi. 7 — \emph{that I may publish with the voice of thanksgiving, and tell of all thy wondrous works} \textperiodcentered\ \textsuperscript{iv}\,Ps. xxvii. 4 — \emph{one thing have I desired of the LORD, that will I seek after; that I may dwell in the house of the LORD all the days of my life} \textperiodcentered\ \textsuperscript{v}\,Ps. xxvii. 8 — \emph{when thou saidst, Seek ye my face; my heart said unto thee, Thy face, LORD, will I seek} \textperiodcentered\ \textsuperscript{vi}\,Ps. cxviii. 19 — \emph{open to me the gates of righteousness: I will go into them, and I will praise the LORD} \textperiodcentered\ \textsuperscript{viii}\,John xx. 28 — \textbf{Thomas}, who would not believe till he had seen: \emph{my Lord and my God}. ⚠⚠ \textbf{At the door of a building where nothing is to be seen, the worshipper says the words of the one apostle who insisted on seeing} — and the next line is the blessing on those who do not. \textperiodcentered\ \textsuperscript{ix}\,John xx. 29 — \emph{blessed are they that have not seen, and yet have believed} \textperiodcentered\ \textsuperscript{x}\,Ps. xxxii. 6 — the plate's figure. ⚠ \emph{I said, I will confess my sins unto the LORD} stands at \textbf{verse 5} in the Authorised Version and at \textbf{6} in the Prayer Book Psalter, which is what the plate follows here — \textbf{a third numbering system on show in this volume}, beside the plain Vulgate and the hybrid at 344 and 353.\par}\par\vspace{3pt}\textsuperscript{vii}\,⚠⚠ \textbf{The same verse is said twice on this leaf, once in each language} — Ps. xxvi. 8 stands in Latin at the door, \emph{Domine, dilexi decorem domus Tuæ}, and returns fourteen lines later in \textbf{Greek}, Κύριε, ἠγάπησα εὐπρέπειαν οἴκου σου. \textbf{The English renders both identically}, which is the rule for a separated doublet; it is not a gloss and neither line may be dropped. ⚠ \textbf{He does not merely cite the Greek Psalter — he prays in it}, which joins the coal at 355, the anaphora and the Greek formula at 349, and the doxology at 350.}
\clearpage
\sethead{§28 \emph{[Θωμᾶς}}
% ---------- printed 357 ----------
\versoalign
\markright{357}
\setnoted{}
\addcontentsline{toc}{section}{§28 \emph{[Θωμᾶς — Pœnitetne?]}}
\originalsone{357}{\input{fragments/la357}}{}
\clearpage
\markright{357}
\setnoted{8,9,14,15,17,19,22,26,28}
\enleafn{357}{\input{fragments/en357}}{}{{\scriptsize \textsuperscript{ii}\,Dan. vi. 4 — the plate's figure. ⚠ \textbf{Daniel kneeling three times a day is at vi. 10}; vi. 4 is the presidents seeking occasion against him. \textbf{Kept as printed} — it stands on the volume's running list with \textup{[Jer. l. 24.]} and \textup{[Eph. i. 8.]} \textperiodcentered\ \textsuperscript{iii}\,Job xlii. 6 — \emph{wherefore I abhor myself, and repent in dust and ashes} \textperiodcentered\ \textsuperscript{iv}\,Ps. lxiii. 6 — ⚠ \textbf{the reference supplies a word, not a sense.} \emph{When I remember thee upon my bed} is the psalmist doing well; the page takes only \emph{lectis} from it and makes the bed \textbf{the thing to be avoided} — if not on the ground and in ashes, yet at least not lying down. \textperiodcentered\ \textsuperscript{v}\,Luke xvi. 19 — \textbf{the rich man of the parable}: \emph{clothed in purple and fine linen, and fared sumptuously every day} \textperiodcentered\ \textsuperscript{vi}\,Dan. i. 8 — Daniel at the king's table, who \emph{would not defile himself with the portion of the king's meat} \textperiodcentered\ \textsuperscript{vii}\,Lev. v. 16 — the trespass offering: \emph{he shall make amends … and shall add the fifth part thereto} \textperiodcentered\ \textsuperscript{viii}\,Deut. xxvi. 12 — the third year's tithe, \emph{the year of tithing}. ⚠ \textbf{The line offers a THIRTIETH part} (\emph{tricesimam partem}), which is not what that verse commands, nor any tithe. \textbf{Kept as printed and on the running list.} \textperiodcentered\ \textsuperscript{ix}\,2 Cor. viii. 3 — of the Macedonians: \emph{to their power, yea, and beyond their power, they were willing of themselves}\par}\par\vspace{3pt}\textsuperscript{i}\,⚠⚠⚠ \textbf{Eleven pairs, and every one of them is a bargain struck downwards: \emph{if not … yet}.} He sets a scriptural ceiling, concedes he cannot reach it, and names the least he will do instead — seven times a day like David, or at any rate three like Daniel; Solomon's long prayer, or at any rate the Publican's short one; Christ's whole night, or at any rate one hour. ⚠⚠ \textbf{And the floor is repeatedly set by a BAD example, not a good one.} If not on the ground in ashes, \textbf{yet not in beds}; if not in sackcloth, \textbf{yet not in purple and fine linen} — which is the rich man of the parable. \textbf{Where the ceiling is a saint, the floor is a sinner: if I cannot be Job, let me at least not be Dives.} ⚠ \textbf{The whole ladder ends by giving up on measure altogether} — \emph{tamen quoad potentiam}, yet according to my power (2 Cor. viii. 3). \textbf{The 21 references on this leaf are the densest in Part II outside the Passion, and the English lines name their own exemplars, so the band selects.}}
\clearpage
\sethead{§29 \emph{Actus Adorationis}}
% ---------- printed 358 ----------
\versoalign
\markright{358}
\setnoted{}
\addcontentsline{toc}{section}{§29 \emph{Actus Adorationis}}
\originalsone{358}{\input{fragments/la358}}{}
\clearpage
\markright{358}
\setnoted{2,6,10,12,15,17}
\enleafn{358}{\input{fragments/en358}}{}{{\scriptsize \textsuperscript{ii}\,Rom. viii. 32 — \emph{he that spared not his own Son, but delivered him up for us all} \textperiodcentered\ \textsuperscript{iii}\,1 John i. 7 — \emph{the blood of Jesus Christ his Son cleanseth us from all sin} \textperiodcentered\ \textsuperscript{iv}\,Eph. iv. 8 — \emph{when he ascended up on high, he led captivity captive} \textperiodcentered\ \textsuperscript{v}\,Matt. iii. 16 — \emph{the Spirit of God descending like a dove, and lighting upon him} \textperiodcentered\ \textsuperscript{vi}\,Acts ii. 3 — \emph{there appeared unto them cloven tongues like as of fire, and it sat upon each of them}\par}\par\vspace{3pt}\textsuperscript{i}\,⚠⚠ \textbf{This is the LITANY's opening, in Latin, and it is a counterweight to the Greek marks of the last eight leaves.} \emph{O God the Father, of heaven · O God the Son, Redeemer of the world · O God the Holy Ghost · O holy, blessed Trinity} — \textbf{the Prayer Book's first four invocations, in the Prayer Book's order.} ⚠ And the English Litany's odd \emph{of heaven}, which is not idiomatic English, is a rendering of the very \emph{de cœlis} standing here: \textbf{both this page and the Litany are drawing on the same Latin.} \textbf{The volume's liturgical debts run in two directions at once, and an introduction that argues only the Greek will be answerable for this leaf.} ⚠ One divergence, stated without conjecture: where the Litany's third invocation names the Spirit \emph{proceeding from the Father and the Son}, this one names him \textbf{the Comforter}.}
\clearpage
% ---------- printed 359 ----------
\versoalign
\markright{359}
\setnoted{}
\originalsone{359}{\input{fragments/la359}}{}
\clearpage
\markright{359}
\setnoted{7,12,19,20,21,22,26,28}
\enleafn{359}{\input{fragments/en359}}{}{{\scriptsize \textsuperscript{i}\,Rom. xiv. 9 — \emph{to this end Christ both died, and rose, and revived, that he might be Lord both of the dead and living} \textperiodcentered\ \textsuperscript{iii}\,Gen. xxiv. 7 — Abraham to his servant, before the journey for Rebekah: \emph{he shall send his angel before thee} \textperiodcentered\ \textsuperscript{iv}\,Matt. ii. 9 — \emph{the star, which they saw in the east, went before them} \textperiodcentered\ \textsuperscript{v}\,Matt. xiv. 30 — Peter on the water: \emph{beginning to sink, he cried, saying, Lord, save me}. ⚠ \textbf{The four examples are two and two: Abraham's servant and the Wise Men were GUIDED, Peter and Paul were RESCUED — one pair for the setting out, one for what may happen on the way.} \textperiodcentered\ \textsuperscript{vi}\,Acts xxvii. 44 — the wreck off Malta: \emph{and so it came to pass, that they escaped all safe to land} \textperiodcentered\ \textsuperscript{vii}\,Ps. lxviii. 1 — \emph{let God arise, let his enemies be scattered} \textperiodcentered\ \textsuperscript{viii}\,Ps. cxix. 115 — \emph{I will keep the commandments of my God}\par}\par\vspace{3pt}\textsuperscript{ii}\,⚠⚠⚠ \textbf{A prayer for the dead, in a bishop of the Church of England's private manual, in the Requiem's own words.} \emph{Defunctis requiem et lucem perpetuam} is the \emph{Requiem æternam dona eis, Domine, et lux perpetua luceat eis} of the Mass for the departed. ⚠ \textbf{The page has just laid the ground for it}: Christ is \emph{Lord alike of the living and of the dead}, and the two halves of one company are named in turn — \emph{whose we are}, whom this world yet holds in the flesh, and \emph{whose also are they}, whom the world to come has already received unclothed of the body. \textbf{The petition then divides its gifts along that line}: to the living mercy and grace, to the departed rest and light. \textbf{The English renders it plainly and does not soften it}, which is the edition's rule and is the only honest course with a text of this kind.}
\clearpage
\sethead{§30 \emph{Oratio peregre profecturi}}
% ---------- printed 360 ----------
\versoalign
\markright{360}
\setnoted{}
\addcontentsline{toc}{section}{§30 \emph{Oratio peregre profecturi · Septem Opera Misericordiæ}}
\originalsone{360}{\input{fragments/la360}}{}
\clearpage
\markright{360}
\setnoted{2,11}
\enleafn{360}{\input{fragments/en360}}{}{\textsuperscript{i}\,\textbf{The scripture foot is empty on this leaf because the leaf cites nothing, and that is correct.} ⚠⚠ \textbf{Both columns are the medieval MEMORY-VERSES, one verb to a work, made to be counted on the fingers} — \emph{Visito, poto, cibo, redimo, tego, colligo, condo} for the corporal works, \emph{Consule, plecte, doce, solare, remitte, fer, ora} for the spiritual. \textbf{They are not composed here; they are quoted, and every schoolman knew them by heart.} ⚠ The usual spiritual verse has \emph{castiga} where this sets \textbf{\textup{plecte}}, a harder word — to punish rather than to correct. \textbf{Their presence says what kind of book this is: a working manual with a checklist in it}, and it is another debt to the Latin West, beside the Litany on the leaf before.\\[1pt]\textsuperscript{ii}\,⚠⚠ \textbf{§31 is the volume's THIRD treatment of the Lord's Prayer, and the method is confession.} §25 set it out plainly (printed 350) and §26 said it seven times over in other men's words (351–354); \textbf{here each petition is stated, then admitted to have been failed, then asked.} \emph{Holy art thou … to be held holy by all men, and by some more than by others, and by me first of all above many} — \textbf{and then: \emph{Yet I have not so held it.}} The petition becomes an indictment before it becomes a request. ⚠ \textbf{The three treatments are not to be conformed}; this is the only one that turns the prayer against the man praying it.}
\clearpage
\sethead{§31 \emph{Πατὲρ ἡμῶν}}
% ---------- printed 361 ----------
\versoalign
\markright{361}
\setnoted{}
\addcontentsline{toc}{section}{§31 \emph{Πατὲρ ἡμῶν} (the Lord's Prayer, petition by petition)}
\originalsone{361}{\input{fragments/la361}}{}
\clearpage
\markright{361}
\setnoted{7,11,14,15,18,19,21,25}
\enleafn{361}{\input{fragments/en361}}{}{{\scriptsize \textsuperscript{iii}\,Matt. xviii. 24 — the servant brought in \emph{which owed him ten thousand talents} \textperiodcentered\ \textsuperscript{v}\,Dan. ix. 8 — the plate's figure. ⚠ \textbf{The words are verse 7's}: printed 35 sets the same pair at \textbf{ix. 7}, and printed 133 prints the two verses in order under their own \textup{Vers.} marks. \textbf{Kept as printed.} \textperiodcentered\ \textsuperscript{vi}\,Hos. xiii. 9 — \emph{O Israel, thou hast destroyed thyself}, which the Latin says in the first person \textperiodcentered\ \textsuperscript{vii}\,Ps. cxxx. 3 — \emph{if thou, LORD, shouldest mark iniquities, O Lord, who shall stand?} ⚠ \textbf{The same two lines stand verbatim at printed 41 under a figure reading cxxx. 30} — the closest witness the volume has against that misprint. \textperiodcentered\ \textsuperscript{viii}\,\textup{Vers. 7, 8} — i.e. \textbf{Ps. cxxx. 7}, \emph{with him is plenteous redemption}; \textbf{Ps. cxxx. 8}, \emph{he shall redeem Israel from all his iniquities} — asked here with the name of Israel gone\par}\par\vspace{3pt}\textsuperscript{i}\,⚠⚠ \textbf{The third petition drives off two wills the plate never references — John i. 13's} \emph{neque ex voluntate carnis, neque ex voluntate viri}. ⚠⚠ \textbf{Then it turns the prayer's own geography onto the man saying it}: \emph{fiat … de terra hac, quæ ego sum} — \textbf{thy will be done upon this earth, which I am}, the brace holding the clod twice, \textbf{by} this earth and \textbf{upon} it. \textbf{That is §31's method entire: every clause is made particular before it is asked.}\\[1pt]\textsuperscript{ii}\,⚠⚠ \textbf{The fourth petition is asked twice here, and the second asking is the Eucharist}: the temporal reading braced — \emph{health, peace, sufficiency} — then, on its own line, \emph{Da panem Angelorum ad salutem æternam}, the manna of Ps. lxxviii. 25, which the plate does not reference. ⚠ \textbf{The volume has never yet let this petition be simply asked} (351, 352, 353); \textbf{this is the fourth refusal and the only one that names what it wants instead.}\\[1pt]\textsuperscript{iv}\,\textbf{Four figures for four failures, and levelling them loses the descent}: the huge sum of debt, \emph{prolapsiones fœdas}, \emph{recidivas crebras} (a physician's word for the sickness that returns), and \emph{volutabra diutina}, \textbf{long wallowings} — the \emph{volutabrum} is the swine's wallowing-place. \textbf{From the counting-house to the sty.}\\[1pt]\textsuperscript{vi}\,⚠ \textbf{The two verses travel together in this book and are here reversed}, so that \emph{to thee righteousness, to me the shame} is now answered by \emph{my destruction is of myself}. ⚠ \textbf{And the Hosea stands in a third Latin form} — \emph{mihi est ex me} (35), \emph{a me est} (the Evening prayers), \emph{mihi a me est} (here); \textbf{none conformed to the others.}}
\clearpage
% ---------- printed 362 ----------
\versoalign
\markright{362}
\setnoted{}
\originalsone{362}{\input{fragments/la362}}{}
\clearpage
\markright{362}
\setnoted{2,3,9,12,21,26}
\enleafn{362}{\input{fragments/en362}}{}{{\scriptsize \textsuperscript{i}\,Ps. lxxxvi. 13 — ⚠⚠ \textbf{the psalm says \emph{thou HAST delivered}; the line asks \emph{Libera}, deliver thou} — the indicative prayed back as a command \textperiodcentered\ \textsuperscript{ii}\,Ps. xlii. 7 — \emph{deep calleth unto deep … thy billows are gone over me} \textperiodcentered\ \textsuperscript{iv}\,Ps. ciii. 14 — \emph{he knoweth our frame; he remembereth that we are dust}, asked here as a petition\par}\par\vspace{3pt}\textsuperscript{ii}\,⚠⚠ \textbf{The verse is inverted}: the psalm's cry of drowning gains a purpose-clause it lacks — \emph{Ut de abysso eripiat}, \textbf{that it may pluck me out of the deep} — so the second deep is God's mercy. ⚠ \textbf{Four lines on, the petition asks for the sins he cannot feel}, and asks not pardon for them but \textbf{light}: \emph{ut possim confiteri}, that I may be \textbf{able} to confess.\\[1pt]\textsuperscript{iii}\,⚠⚠⚠ \textbf{The sixth petition is re-said three times to get God out of the way of the tempting.} \emph{Et ne inducas} — \textbf{quoted headless, with no object at all} — then \emph{ne me induci sinas}, suffer me not to be led in, then \emph{ne me sinas intrare}, suffer me not to enter. \textbf{God leading, God permitting, and at last the man walking in himself} — the middle being the ancient reading Cyprian and Augustine defended, \textbf{the third going past both.}\\[1pt]\textsuperscript{v}\,\textbf{Augustine, and famous}: \emph{Hic ure, hic seca, ut in æternum parcas}, re-cast around the two ages as \textbf{there spare, here burn, here cut.}\\[1pt]\textsuperscript{vi}\,⚠⚠ \textbf{The last petition descends out of theology into the man's own week and ends in the present tense.} The evils are anatomised downward — flesh, Fiend, penal evils, the world to come, the world that now is — and then stop being categories: \textbf{the malice of THIS sickness, \emph{wherewith I wrestle},} and the business \textbf{I am entangled in.} ⚠ \textbf{Both are present indicatives about the man holding the book} — the same seam as the \emph{Divisio} at 347–348.}
\clearpage
% ---------- printed 363 ----------
\versoalign
\markright{363}
\setnoted{}
\originalsone{363}{\input{fragments/la363}}{}
\clearpage
\markright{363}
\setnoted{4,5,13,24}
\enleafn{363}{\input{fragments/en363}}{}{{\scriptsize \textsuperscript{i}\,Dan. ix. 16 — Daniel in the exile: \emph{O Lord, according to all thy righteousness, let thine anger be turned away}\par}\par\vspace{3pt}\textsuperscript{i}\,⚠⚠ \textbf{This sentence stands three times in Part II in three Latin dresses, and the plate tags it once — here.} Printed 343 (\emph{ut auferatur ira Tua}) and 373 (\emph{Ut avertatur a me justissima indignatio}) are unmarked. ⚠ \textbf{The Hosea two leaves back behaves the same way: a concordance built on the volume's own marks finds one of three and misses two.}\\[1pt]\textsuperscript{ii}\,⚠ \textbf{The brace makes two superlatives commute} — \emph{sæpissime, gravissime} and \emph{gravissime, sæpissime} — and both limbs are printed in both orders. ⚠⚠ \textbf{\emph{Domino Reverendo, familiaque} is not resolved}: "the Reverend Lord, and his household" may be his diocesan, his patron, or the master of a house he served in, and \textbf{this witness cannot settle it.}\\[1pt]\textsuperscript{iii}\,⚠ \textbf{§32 begins here, and its heading is a technical word.} Ἀνακεφαλαίωσις is the noun of Eph. i. 10's \emph{ἀνακεφαλαιώσασθαι}, to gather all things under one head in Christ — Irenæus' term, not a synonym for "summary." \textbf{What it gathers is the Creed stripped to its events}, Conception to Judgement, with no doctrine, no Church, no forgiveness of sins among them; and it closes not \emph{credo} but \textbf{\emph{Fac me horum participem}, make me partaker of them.}\\[1pt]\textsuperscript{iv}\,⚠⚠⚠ \textbf{The \emph{Interpellatio Eucharistica} is the most doctrinally loaded line in Part II, and it is seven verbs long}: \emph{Recolo, gratias ago, admoneo, recordor, commemoro, offero, vel peto ut offeras}. ⚠ \textbf{\emph{Admoneo} points the memorial at GOD} — I put thee in mind — which is the patristic understanding of the anamnesis and not the ordinary Western one. ⚠⚠ \textbf{And \emph{offero} is qualified in the same breath by \emph{vel peto ut offeras}: the sacrifice is offered, and then in one word the offerer is left an open question.} A whole controversy hangs on \emph{vel}, and the English keeps it there and settles nothing. \textbf{This leaf belongs with the anaphora at 349 in any account of what rite the book was written out of.}}
\clearpage
\sethead{§32 \emph{Articulorum Fidei Ἀνακεφαλαίωσις}}
% ---------- printed 364 ----------
\versoalign
\markright{364}
\setnoted{}
\addcontentsline{toc}{section}{§32 \emph{Articulorum Fidei Ἀνακεφαλαίωσις · Interpellatio Eucharistica}}
\originalsone{364}{\input{fragments/gr364}}{}
\clearpage
\markright{364}
\setnoted{1,6,7,11,15,18}
\enleafn{364}{\input{fragments/en364}}{}{{\scriptsize \textsuperscript{vi}\,Ps. xliv. 1 — \emph{our fathers have told us, what work thou didst in their days, in the times of old}\par}\par\vspace{3pt}\textsuperscript{i}\,\textbf{The scripture band is nearly empty here because the plate cites almost nothing, and that is correct} — §32 pleads the mysteries themselves, not texts about them, so \textbf{what the leaf needs is its Greek opened.} ⚠ \textbf{The two columns are not an antithesis} — the Person on the left, the life on the right. ⚠ \textbf{\emph{Ἀπαρχὰς αἵματος}, the firstfruits of thy blood}, is the Circumcision counted as the first shedding, eight days old; \emph{ἀπαρχή} is the firstfruits word of 1 Cor. xv. 20. \textbf{The article is given a reason for standing in a plea about the Passion.}\\[1pt]\textsuperscript{ii}\,⚠⚠ \textbf{The left brace is Philippians ii turned into nouns and run BACKWARDS.} \emph{Ἐκκένωσις} and \emph{ταπείνωσις} are the substantives of \emph{ἐκένωσεν} and \emph{ἐταπείνωσεν ἑαυτόν}, and after them come the enfleshing and \emph{σύλληψις ἐν γαστρί}, the conception in the womb: \textbf{from the self-emptying down to the womb, later cause to earlier — the order of descent, not of narrative.} ⚠ \emph{Ἐκκένωσις} is rendered \textbf{emptying-out}, not "exinanition" — the κένωσις opened at printed 305.\\[1pt]\textsuperscript{iii}\,⚠⚠ \textbf{The Son of man with nowhere to lay his head returns here doing the opposite of what it did at printed 238.} There it was turned round to indict the man praying — \emph{no fit place to lay it here either, in me}; \textbf{here it is simply pleaded, one of Christ's own hardships beside hunger and thirst and watching.} ⚠ \textbf{Not cross-referred and not to be conformed}: this volume uses a text twice and never for the same thing.\\[1pt]\textsuperscript{iv}\,⚠ \textbf{Two pleas are in the words of the events themselves}, neither referenced by the plate: \emph{κατάληψιν ὡς λῃστοῦ}, \textbf{apprehension as of a thief}, Christ's own complaint in the garden; and \emph{συνοχὴ ἕως σταυροῦ}, the \emph{συνέχομαι} of \emph{how am I straitened till it be accomplished}. ⚠ Then \textbf{three places under one preposition} — garden, pavement and place of a skull — \textbf{left untranslated, as the plate has them.}\\[1pt]\textsuperscript{v}\,⚠⚠ \textbf{\emph{Ἐξομολόγησις} here is NOT confession of sin} but the word's older sense, \textbf{acknowledgement and so praise}, which is why its limbs are \emph{of Grace, a Song} and \emph{of Praises, a Psalm}. \textbf{A reader bringing the penitential sense across from §17 and §22 misreads the close of the section.} ⚠ The line above gives prayer an affirmative and a negative form, \textbf{in Latin under Greek labels.}\\[1pt]\textsuperscript{vi}\,⚠ \textbf{§33 opens here on one Latin line and then runs Greek to the end} — twenty-one deliverances of the fathers, each a name and what he was delivered from. \textbf{The argument is the psalm's own}: what was done of old is the ground for asking it now.}
\clearpage
\sethead{§33 \emph{Sicut patres, et nos libera, Domine}}
% ---------- printed 365 ----------
\versoalign
\markright{365}
\setnoted{}
\addcontentsline{toc}{section}{§33 \emph{Sicut patres, et nos libera, Domine}}
\originalsone{365}{\input{fragments/gr365}}{}
\clearpage
\markright{365}
\setnoted{}
\enleaf{365}{\input{fragments/en365}}{}
\clearpage
\sethead{§34 \emph{Præparatio ad Εὐλογίαν}}
% ---------- printed 366 ----------
\versoalign
\markright{366}
\setnoted{}
\addcontentsline{toc}{section}{§34 \emph{Præparatio ad Εὐλογίαν}}
\originalsone{366}{\input{fragments/la366}}{}
\clearpage
\markright{366}
\setnoted{}
\enleaf{366}{\input{fragments/en366}}{}
\clearpage
% ---------- printed 367 ----------
\versoalign
\markright{367}
\setnoted{}
\originalsone{367}{\input{fragments/la367}}{}
\clearpage
\markright{367}
\setnoted{}
\enleaf{367}{\input{fragments/en367}}{}
\clearpage
% ---------- printed 368 ----------
\versoalign
\markright{368}
\setnoted{}
\originalsone{368}{\input{fragments/gr368}}{}
\clearpage
\markright{368}
\setnoted{}
\enleaf{368}{\input{fragments/en368}}{}
\clearpage
% ---------- printed 369 ----------
\versoalign
\markright{369}
\setnoted{}
\originalsone{369}{\input{fragments/gr369}}{}
\clearpage
\markright{369}
\setnoted{}
\enleaf{369}{\input{fragments/en369}}{}
\clearpage
\sethead{§35 \emph{Oratio Tho. Bradwardini}}
% ---------- printed 370 ----------
\versoalign
\markright{370}
\setnoted{}
\addcontentsline{toc}{section}{§35 \emph{Oratio Tho. Bradwardini}}
\originalsone{370}{\input{fragments/la370}}{}
\clearpage
\markright{370}
\setnoted{}
\enleaf{370}{\input{fragments/en370}}{}
\clearpage
\sethead{§36 [\emph{Pro} intercession, headless]}
% ---------- printed 371 ----------
\versoalign
\markright{371}
\setnoted{}
\addcontentsline{toc}{section}{§36 [\emph{Pro} intercession, headless] + \emph{De Pœnitentia}}
\originalsone{371}{\input{fragments/la371}}{}
\clearpage
\markright{371}
\setnoted{}
\enleaf{371}{\input{fragments/en371}}{}
\clearpage
% ---------- printed 372 ----------
\versoalign
\markright{372}
\setnoted{}
\originalsone{372}{\input{fragments/la372}}{}
\clearpage
\markright{372}
\setnoted{}
\enleaf{372}{\input{fragments/en372}}{}
\clearpage
\sethead{§37 \emph{Confessio Peccatorum}}
% ---------- printed 373 ----------
\versoalign
\markright{373}
\setnoted{}
\addcontentsline{toc}{section}{§37 \emph{Confessio Peccatorum}}
\originalsone{373}{\input{fragments/la373}}{}
\clearpage
\markright{373}
\setnoted{}
\enleaf{373}{\input{fragments/en373}}{}
\clearpage
% ---------- printed 374 ----------
\versoalign
\markright{374}
\setnoted{}
\originalsone{374}{\input{fragments/la374}}{}
\clearpage
\markright{374}
\setnoted{}
\enleaf{374}{\input{fragments/en374}}{}
\clearpage
% ---------- printed 375 ----------
\versoalign
\markright{375}
\setnoted{}
\originalsone{375}{\input{fragments/la375}}{}
\clearpage
\markright{375}
\setnoted{}
\enleaf{375}{\input{fragments/en375}}{}
\clearpage
% ---------- printed 376 ----------
\versoalign
\markright{376}
\setnoted{}
\originalsone{376}{\input{fragments/la376}}{}
\clearpage
\markright{376}
\setnoted{}
\enleaf{376}{\input{fragments/en376}}{}
\clearpage
% ---------- printed 377 ----------
\versoalign
\markright{377}
\setnoted{}
\originalsone{377}{\input{fragments/la377}}{}
\clearpage
\markright{377}
\setnoted{}
\enleaf{377}{\input{fragments/en377}}{}
\clearpage
% ---------- printed 378 ----------
\versoalign
\markright{378}
\setnoted{}
\originalsone{378}{\input{fragments/la378}}{}
\clearpage
\markright{378}
\setnoted{}
\enleaf{378}{\input{fragments/en378}}{}
\clearpage
% ---------- printed 379 ----------
\versoalign
\markright{379}
\setnoted{}
\originalsone{379}{\input{fragments/la379}}{}
\clearpage
\markright{379}
\setnoted{}
\enleaf{379}{\input{fragments/en379}}{}
\clearpage
\sethead{§38 \emph{Oratiuncula post Confessionem}}
% ---------- printed 380 ----------
\versoalign
\markright{380}
\setnoted{}
\addcontentsline{toc}{section}{§38 \emph{Oratiuncula post Confessionem} + \emph{Oratio Præparatoria}}
\originalsone{380}{\input{fragments/la380}}{}
\clearpage
\markright{380}
\setnoted{}
\enleaf{380}{\input{fragments/en380}}{}
\clearpage
\sethead{§39 \emph{Ante Concionem Cautela cum Precatione ex Fulgentii l…}}
% ---------- printed 381 ----------
\versoalign
\markright{381}
\setnoted{}
\addcontentsline{toc}{section}{§39 \emph{Ante Concionem Cautela cum Precatione ex Fulgentii l. i. ad Mon.}}
\originalsone{381}{\input{fragments/la381}}{}
\clearpage
\markright{381}
\setnoted{}
\enleaf{381}{\input{fragments/en381}}{}
\clearpage
\sethead{§40 Ἀφορμαὶ \emph{Meditationum}}
% ---------- printed 382 ----------
\versoalign
\markright{382}
\setnoted{}
\addcontentsline{toc}{section}{§40 Ἀφορμαὶ \emph{Meditationum} (three catenæ: ante Preces Pœnitentiales · ante Intercessionum · ante Εὐχαριστίαν ἢ Εὐλογίαν)}
\originalsone{382}{\input{fragments/la382}}{}
\clearpage
\markright{382}
\setnoted{}
\enleaf{382}{\input{fragments/en382}}{}
\clearpage
% ---------- printed 383 ----------
\versoalign
\markright{383}
\setnoted{}
\originalsone{383}{\input{fragments/la383}}{}
\clearpage
\markright{383}
\setnoted{}
\enleaf{383}{\input{fragments/en383}}{}
\clearpage
% ---------- printed 384 ----------
\versoalign
\markright{384}
\setnoted{}
\originalsone{384}{\input{fragments/la384}}{}
\clearpage
\markright{384}
\setnoted{}
\enleaf{384}{\input{fragments/en384}}{}
\clearpage
\sethead{§41 \emph{Oratio præparatoria ante} Εὐλογίαν}
% ---------- printed 385 ----------
\versoalign
\markright{385}
\setnoted{}
\addcontentsline{toc}{section}{§41 \emph{Oratio præparatoria ante} Εὐλογίαν + \emph{Monita et Meditationes præparatoriæ in Vespertina ad Deum elevatione mentis} (→ \emph{Scrutinium et Inquisitio, vel Examen})}
\originalsone{385}{\input{fragments/la385}}{}
\clearpage
\markright{385}
\setnoted{}
\enleaf{385}{\input{fragments/en385}}{}
\clearpage
% ---------- printed 386 ----------
\versoalign
\markright{386}
\setnoted{}
\originalsone{386}{\input{fragments/la386}}{}
\clearpage
\markright{386}
\setnoted{}
\enleaf{386}{\input{fragments/en386}}{}
\clearpage
% ---------- printed 387 ----------
\versoalign
\markright{387}
\setnoted{}
\originalsone{387}{\input{fragments/la387}}{}
\clearpage
\markright{387}
\setnoted{}
\enleaf{387}{\input{fragments/en387}}{}
\clearpage
\sethead{§42 ΥΜΝΟΣ ΕΩΘΙΝΟΣ / \emph{Hymnus Matutinus}}
% ---------- printed 388 ----------
\versoalign
\markright{388}
\setnoted{}
\addcontentsline{toc}{section}{§42 ΥΜΝΟΣ ΕΩΘΙΝΟΣ / \emph{Hymnus Matutinus} + ΥΜΝΟΣ ΕΣΠΕΡΙΝΟΣ / \emph{Hymnus Vespertinus}}
\originalsone{388}{\input{fragments/gr388}}{}
\clearpage
\markright{388}
\setnoted{}
\enleaf{388}{\input{fragments/en388}}{}
\clearpage
% ---------- printed 389 ----------
\versoalign
\markright{389}
\setnoted{}
\originalsone{389}{\input{fragments/la389}}{}
\clearpage
% ---------- printed 390 ----------
\versoalign
\markright{390}
\setnoted{}
\originalsone{390}{\input{fragments/gr390}}{}
\clearpage
\markright{390}
\setnoted{}
\enleaf{390}{\input{fragments/en390}}{}
\clearpage
% ---------- printed 391 ----------
\versoalign
\markright{391}
\setnoted{}
\originalsone{391}{\input{fragments/la391}}{}
\clearpage
\sethead{§43 ΕΙΣ ΤΗΝ ΤΟΥ ΧΡΙΣΤΟΥ ΣΤΑΥΡΩΣΙΝ, Μονοστροφικά / \emph{In Chr…}}
% ---------- printed 392 ----------
\versoalign
\markright{392}
\setnoted{}
\addcontentsline{toc}{section}{§43 ΕΙΣ ΤΗΝ ΤΟΥ ΧΡΙΣΤΟΥ ΣΤΑΥΡΩΣΙΝ, Μονοστροφικά / \emph{In Christum Crucifixum, Monostrophica}}
\originalsone{392}{\input{fragments/gr392}}{}
\clearpage
\markright{392}
\setnoted{}
\enleaf{392}{\input{fragments/en392}}{}
\clearpage
% ---------- printed 393 ----------
\versoalign
\markright{393}
\setnoted{}
\originalsone{393}{\input{fragments/la393}}{}
\clearpage
% ---------- printed 394 ----------
\versoalign
\markright{394}
\setnoted{}
\originalsone{394}{\input{fragments/gr394}}{}
\clearpage
\markright{394}
\setnoted{}
\enleaf{394}{\input{fragments/en394}}{}
\clearpage
% ---------- printed 395 ----------
\versoalign
\markright{395}
\setnoted{}
\originalsone{395}{\input{fragments/la395}}{}
\clearpage

\backmatter
% Plain style: the fancy heads carry the 1853 shoulder-note in \rightmark, which
% the back matter has no use for and would otherwise fight with its own \markboth.
\pagestyle{plain}
% Appendix: an illustrative note on Andrewes' sources. DRAFT — Wilson's to edit.
% GROUNDING. Every page reference here was checked against our own transcripts,
% and every attribution is either (a) printed in the 1853 plate, (b) supplied by
% the 1853 editor in square brackets, or (c) identified by Brightman 1903 and
% recorded in apparatus/BRIGHTMAN-collation.md. The distinction is load-bearing
% and the text keeps it. The count of scripture references is from the plate.
% Nothing here is an attribution of our own making.
\begingroup
\setlength{\parskip}{0.55\baselineskip}
\setlength{\parindent}{0pt}
\addcontentsline{toc}{section}{Appendix: Some of the Sources}
{\centering\Large\scshape Appendix\par}
\vspace{0.4\baselineskip}
{\centering\itshape Some of the sources Andrewes drew on\par}
\vspace{1.5\baselineskip}

This note is illustrative and not a source apparatus. It names enough of what
Andrewes gathered to show the kind of book he was making, and stops well short of
what a full accounting would require.

Three sorts of attribution are mixed in what follows, and they are not of equal
weight. Some authorities are printed in the 1853 itself, standing in the margin at
the head of the sentence they belong to --- \emph{Cyprian.}, \emph{Beda.},
\emph{Seneca.} Some are the 1853 editor's own identifications, which he set in
square brackets, and which are printed here in brackets for that reason. And some
were identified neither by Andrewes nor by his editor but by F.\,E. Brightman in
1903, working from the sources; where a line below rests on Brightman it says so.
The page numbers are the 1853's, as printed on the inner edge of every leaf here.

\vspace{0.8\baselineskip}
{\scshape Scripture, before anything else.}\quad
The 1853 prints better than two thousand scripture references. That number is the
plainest fact about this book: it is a cento, and the greater part of what
Andrewes prays is quotation. The text he quotes is the Septuagint and the Vulgate,
not any English Bible --- which is why the translation here reaches for
Coverdale's Psalter ahead of the Authorised Version, as the preface explains. The
Prayer of Manasses is drawn on at printed 171 (the editor's
identification, [\emph{Orat.\ Manassis.}]).

\vspace{0.8\baselineskip}
{\scshape The liturgies.}\quad
The Liturgy of S.\ Chrysostom is named twice in Part I, at printed 33 and 55, in
the plate itself. The \emph{Sursum corda} that opens the epigraph leaf of Part III
at printed 398 is older than any of the liturgies Andrewes could have held: it
stands in the Hippolytean canons and in Cyprian's treatise on the Lord's Prayer,
an identification we owe to Brightman.

\vspace{0.8\baselineskip}
{\scshape The creeds and the Prayer Book.}\quad
The Nicene Creed is named in the plate at printed 177. The \emph{Te Deum} is drawn
on at printed 25 and again at 431, on the editor's identification. And at printed
199 the margin reads [\emph{Litan.\ Angl.}] --- the English Litany. A bishop's
private book, kept in Greek and Latin, quietly quoting the vernacular liturgy he
had to lead on Sundays.

\vspace{0.8\baselineskip}
{\scshape The Greek fathers.}\quad
Chrysostom again at printed 417. Theophylact opens Part III at printed 398, and
Brightman located the passage exactly: the commentary on S.\ Luke xviii, where the
Greek \emph{βάθρον καὶ κρηπίς} is rendered by the Latin \emph{fundamentum et
basis} --- foundation and ground --- which is the first line of that leaf. Gregory
of Nyssa stands in the catena at printed 386, from the first homily on the Lord's
Prayer.

\vspace{0.8\baselineskip}
{\scshape The Latin fathers.}\quad
Part II gathers them in a short space, each named in the plate at the head of its
sentence: Cyprian at printed 381, then Bede, Arnobius and Jerome at 382, then
Augustine twice at 383 --- once from \emph{contra Maximinum} and once from the
first book of the \emph{Confessions} --- and Augustine again at 417 and 422.

One of these is worth a note of its own. At printed 380 the plate attributes a
passage \emph{ex Fulgentii l.\ i.\ ad Mon.}, to the first book of Fulgentius of
Ruspe \emph{to Monimus}. Our own reading suspected the passage was Augustine's,
and the attribution was kept as printed rather than mended. Brightman, working
independently from the sources, gives Fulgentius too. The plate was right and the
doubt was ours.

\vspace{0.8\baselineskip}
{\scshape A medieval archbishop.}\quad
Printed 370 opens \emph{Oratio Tho.\ Bradwardini} --- Thomas Bradwardine,
archbishop of Canterbury for five weeks in 1349. Brightman identifies the work as
\emph{de causa Dei contra Pelagium}, and the edition Andrewes is likely to have
used as Savile's, London 1618.

\vspace{0.8\baselineskip}
{\scshape Pagan antiquity, and no apology for it.}\quad
Cicero is named in the plate at printed 385 and again at 386, the second time for
Cato, who called himself to account each night for the day's business ---
\emph{de senectute}, on Brightman's identification. Seneca is named at printed
385. At 386 Andrewes cites \emph{Ausonius out of Pythagoras}, a quotation at one
remove and marked as such.

Two things about this stretch repay attention. The first is that nothing in the
page ranks these authorities below the Fathers or above them. Bede, Augustine,
Cicero and Seneca stand in one column, each with his sentence, and the reader is
left to make of that what he will. The second is smaller and more human: the
sentence Seneca gives --- that you must catch yourself out before you can mend
yourself --- Seneca himself had from Epicurus. The chain is Epicurus to Seneca to
Andrewes, and a reader coming on that page would very likely take the thought for
a Father's.

\vspace{0.8\baselineskip}
{\scshape Jewish and contemporary.}\quad
At printed 386, in the middle of his own sentence rather than in a margin,
Andrewes writes \emph{Recte igitur Rab.\ J.} --- rightly, then, Rabbi J. --- on
not putting repentance off until tomorrow. Hebrew stands in his text in more
places than the 1853 prints; what it left out is recorded in the apparatus, and
the preface says why it has not been restored here. On the following page,
printed 387, he cites \emph{Rev.\ Usserius} --- James Ussher, then a
contemporary --- with page numbers, as one cites a book lately read.

\vspace{0.8\baselineskip}
{\scshape A caution about the margins.}\quad
The attributions printed in the plate are not always exact, and they have been
left as they are. The clearest instance is the \emph{Seneca.} at printed 385,
which stands at the head of two sentences. The first is Seneca's, from the
twenty-eighth moral letter. The second --- that a sore never found festers and
goes uncured --- is not in that letter and has not been located anywhere else.
It may be Andrewes' own amplification of Seneca's metaphor; it may equally come
from a collection of sentences of the kind he plainly used. The honest statement
is that it is unattested elsewhere, and that one margin covers two sentences of
which only the first is certainly Seneca's. A margin in this book is a pointer,
not a warranty.

\endgroup
\clearpage


\clearpage
\immediate\closeout\fitfile
\end{document}
